% Géométrie analytique dans l'espace — Mathématiques 1ère E — Cours signé OnBuch+
% Programme officiel MINESEC. Compiler avec : tectonic N19-geometrie-analytique-espace.tex
\newcommand{\DOCMATIERE}{Mathématiques}
\newcommand{\DOCNIVEAU}{1ère E}
\newcommand{\DOCMODULE}{Géométrie dans l'espace}
\newcommand{\DOCLECON}{N19}
\newcommand{\DOCTITRE}{Géométrie analytique dans l'espace}
% OnBuch+ — préambule « cours signés » ENRICHI (compilé avec tectonic / XeLaTeX)
% Système visuel « Pop » d'OnBuch+ : fond crème (jamais blanc), bordures encre
% épaisses, ombres « dures » décalées, titres Archivo Black en capitales.
% Version MATHÉMATIQUES : graphes (pgfplots/tikz), boîtes pédagogiques
% (définition, propriété, méthode, attention, le savais-tu, exercice résolu,
% exercice, TP, bilan), page de garde et signature OnBuch+ ancrée en bas.
%
% Macros attendues dans main.tex AVANT \input{preamble} :
%   \DOCMATIERE  \DOCNIVEAU  \DOCMODULE  \DOCLECON (numéro)  \DOCTITRE
\documentclass[11pt]{article}

\usepackage{fontspec}
\usepackage[french]{babel}
\usepackage[a4paper,margin=2cm,top=3.4cm,bottom=2.8cm,headheight=1.4cm,footskip=1.3cm]{geometry}
\usepackage{amsmath,amssymb,amsfonts}
\usepackage{mathtools} % \xmapsto, \coloneqq... souvent utilisés par les modèles
\usepackage{cancel} % simplifications barrées
% \abs{z} (module, valeur absolue) et \norm{u} (norme), utilisés par les modèles
\providecommand{\abs}{}\let\abs\relax\DeclarePairedDelimiter{\abs}{\lvert}{\rvert}
\providecommand{\norm}{}\let\norm\relax\DeclarePairedDelimiter{\norm}{\lVert}{\rVert}
\usepackage[table]{xcolor} % [table] : fournit aussi \rowcolors (alternance de lignes)
\usepackage{pagecolor}
\usepackage{tikz}
\usetikzlibrary{arrows.meta,positioning,calc,shapes.geometric,shapes.misc,shapes.arrows,decorations.pathmorphing,decorations.pathreplacing,decorations.markings,patterns,shadows,mindmap,trees,fit,backgrounds,matrix,intersections,angles,quotes}
\usepackage{pgfplots}
\usepackage{pgf-pie} % diagrammes circulaires \pie (statistiques)
\pgfplotsset{compat=1.18}
\usepgfplotslibrary{fillbetween,groupplots} % aires entre courbes (calcul intégral)
\usepackage{enumitem}
\usepackage{fancyhdr}
% [most] charge toutes les bibliothèques tcolorbox (listings, minted, raster,
% documentation...) et gonfle inutilement la mémoire TeX sur un document long
% (~300 boîtes) : on ne charge que ce qui est réellement utilisé.
\usepackage[skins,breakable]{tcolorbox}
\usepackage{graphicx}
\usepackage{colortbl}
\usepackage{booktabs}
\usepackage{tabularx}
\usepackage{diagbox} % case d'en-tête diagonale des tableaux à double entrée
% Colonnes extensibles centrée / gauche / droite (C, L, R), souvent utilisées par les modèles
\newcolumntype{C}{>{\centering\arraybackslash}X}
\newcolumntype{L}{>{\raggedright\arraybackslash}X}
\newcolumntype{R}{>{\raggedleft\arraybackslash}X}
\usepackage{array}
\usepackage{multicol}
\usepackage{siunitx}
% parse-numbers=false : accepte aussi \SI{2,5 \times 10^{-3}}{...}
\sisetup{locale=FR,output-decimal-marker={,},group-minimum-digits=4,parse-numbers=false}
\DeclareSIUnit\ohm{\ensuremath{\symup{\Omega}}} % U+2126 absent de la police
\usepackage{qrcode}
% \degree : commande fréquemment utilisée par les modèles pour un angle ou une
% température en dehors de \SI{}{} (non fournie par les paquets chargés).
\providecommand{\degree}{^{\circ}}
% \qed (fin de démonstration), utilisé par les modèles sans amsthm
\providecommand{\qed}{\hfill$\square$}
% \barre : double barre des valeurs interdites dans un tableau de variations (tkz-tab)
\providecommand{\barre}{\ensuremath{\|}}
% environnement proof (amsthm non chargé) utilisé par les modèles
\ifdefined\proof\else\newenvironment{proof}[1][Démonstration]{\par\noindent\textit{#1.}\ }{\qed\par}\fi
\usepackage{lastpage}
\usepackage{hyperref}
\hypersetup{hidelinks}

% ============================================================
% Palette « Pop » OnBuch+ — mêmes hex que la classe P (plus_ui.dart)
% ============================================================
\definecolor{poporange}{HTML}{F59321}
\definecolor{popdark}{HTML}{E07A0C}
\definecolor{popcream}{HTML}{FFF6E8}
\definecolor{popink}{HTML}{1C1714}
\definecolor{poppurple}{HTML}{7A5AE0}
\definecolor{popgreen}{HTML}{2E9E6A}
\definecolor{poppink}{HTML}{E0457B}
\definecolor{popblue}{HTML}{2F7DE1}
\definecolor{popgold}{HTML}{C98A12}
\definecolor{popmuted}{HTML}{8A7F76}
\definecolor{popline}{HTML}{EADFD2}
\definecolor{popgray}{HTML}{8A7F76}
\colorlet{popgrey}{popgray}
% Alias tolérants (le modèle nomme parfois les couleurs différemment)
\colorlet{popwhite}{white}
\colorlet{popblack}{popink}
\colorlet{popred}{poppink}
\colorlet{popyellow}{popgold}
% Teintes claires pour fonds de tableaux / zones
\colorlet{popblueL}{popblue!10!white}
\colorlet{poporangeL}{poporange!14!white}
\colorlet{popgreenL}{popgreen!12!white}
\colorlet{poppurpleL}{poppurple!12!white}
\colorlet{poppinkL}{poppink!12!white}
\colorlet{popgoldL}{popgold!14!white}

\pagecolor{popcream}

% ============================================================
% Typographie
% ============================================================
\defaultfontfeatures{Ligatures=TeX}
\setmainfont{PlusJakartaSans-Medium.ttf}[Path=../../../fonts/,BoldFont=PlusJakartaSans-Bold.ttf]
% Maths en sans-serif (Fira Math) avec chiffres et lettres droites de Plus
% Jakarta, pour que les formules se fondent dans le texte.
\usepackage[math-style=ISO,bold-style=ISO]{unicode-math}
\setmathfont{FiraMath-Regular.otf}[Path=../../../fonts/]
\setmathfont{PlusJakartaSans-Medium.ttf}[Path=../../../fonts/,range={up/num,up/{Latin,latin},\mathup}]
\usepackage{icomma} % virgule décimale sans espace en mode math
% Caractères mathématiques Unicode absents des polices de texte (Plus Jakarta,
% Archivo Black) : rendus par leur équivalent mathématique, en texte comme en
% formule — sinon ils s'affichent en carré vide (ex. « Arithmétique dans ℤ »).
\usepackage{newunicodechar}
\newunicodechar{ℝ}{\ensuremath{\mathbb{R}}}
\newunicodechar{ℤ}{\ensuremath{\mathbb{Z}}}
\newunicodechar{ℂ}{\ensuremath{\mathbb{C}}}
\newunicodechar{ℕ}{\ensuremath{\mathbb{N}}}
\newunicodechar{ℚ}{\ensuremath{\mathbb{Q}}}
\newunicodechar{𝔻}{\ensuremath{\mathbb{D}}}
\newunicodechar{α}{\ensuremath{\alpha}}
\newunicodechar{β}{\ensuremath{\beta}}
\newunicodechar{γ}{\ensuremath{\gamma}}
\newunicodechar{δ}{\ensuremath{\delta}}
\newunicodechar{ε}{\ensuremath{\varepsilon}}
\newunicodechar{θ}{\ensuremath{\theta}}
\newunicodechar{λ}{\ensuremath{\lambda}}
\newunicodechar{μ}{\ensuremath{\mu}}
\newunicodechar{σ}{\ensuremath{\sigma}}
\newunicodechar{τ}{\ensuremath{\tau}}
\newunicodechar{φ}{\ensuremath{\varphi}}
\newunicodechar{ψ}{\ensuremath{\psi}}
\newunicodechar{ω}{\ensuremath{\omega}}
\newunicodechar{Σ}{\ensuremath{\Sigma}}
% majuscules produites par \MakeUppercase dans les titres (α-aminés -> Α)
\newunicodechar{Α}{\ensuremath{\alpha}}
\newunicodechar{Β}{\ensuremath{\beta}}
\newunicodechar{Γ}{\ensuremath{\Gamma}}
\newunicodechar{Δ}{\ensuremath{\Delta}}
\newunicodechar{Ε}{\ensuremath{\varepsilon}}
\newunicodechar{Θ}{\ensuremath{\Theta}}
\newunicodechar{Λ}{\ensuremath{\Lambda}}
\newunicodechar{Μ}{\ensuremath{\mu}}
\newunicodechar{Τ}{\ensuremath{\tau}}
\newunicodechar{Φ}{\ensuremath{\Phi}}
\newunicodechar{Ψ}{\ensuremath{\Psi}}
\newunicodechar{Ω}{\ensuremath{\Omega}}
\newunicodechar{↦}{\ensuremath{\mapsto}}
\newunicodechar{∈}{\ensuremath{\in}}
\newunicodechar{∉}{\ensuremath{\notin}}
\newunicodechar{∧}{\ensuremath{\wedge}}
\newunicodechar{⇔}{\ensuremath{\Leftrightarrow}}
\newunicodechar{⟺}{\ensuremath{\Longleftrightarrow}}
\newunicodechar{⇒}{\ensuremath{\Rightarrow}}
\newunicodechar{⟹}{\ensuremath{\Longrightarrow}}
\newunicodechar{∘}{\ensuremath{\circ}}
\newunicodechar{∪}{\ensuremath{\cup}}
\newunicodechar{∩}{\ensuremath{\cap}}
\newunicodechar{≡}{\ensuremath{\equiv}}
\newunicodechar{⊂}{\ensuremath{\subset}}
\newunicodechar{⊥}{\ensuremath{\perp}}
\newunicodechar{∃}{\ensuremath{\exists}}
\newunicodechar{∀}{\ensuremath{\forall}}
\newunicodechar{∛}{\ensuremath{\sqrt[3]{\,}}}
\newunicodechar{∜}{\ensuremath{\sqrt[4]{\,}}}
\newunicodechar{⃗}{\ensuremath{{}^{\to}}}
\newunicodechar{ⁿ}{\ifmmode^{n}\else\textsuperscript{\ensuremath{n}}\fi}
\newunicodechar{ˣ}{\ifmmode^{x}\else\textsuperscript{\ensuremath{x}}\fi}
\newunicodechar{ᵇ}{\ifmmode^{b}\else\textsuperscript{\ensuremath{b}}\fi}
\newunicodechar{ʳ}{\ifmmode^{r}\else\textsuperscript{\ensuremath{r}}\fi}
\newunicodechar{ᵘ}{\ifmmode^{u}\else\textsuperscript{\ensuremath{u}}\fi}
\newunicodechar{ʸ}{\ifmmode^{y}\else\textsuperscript{\ensuremath{y}}\fi}
\newunicodechar{ᵃ}{\ifmmode^{a}\else\textsuperscript{\ensuremath{a}}\fi}
\newunicodechar{ᵉ}{\ifmmode^{e}\else\textsuperscript{\ensuremath{e}}\fi}
\newunicodechar{ᵏ}{\ifmmode^{k}\else\textsuperscript{\ensuremath{k}}\fi}
\newunicodechar{ᵖ}{\ifmmode^{p}\else\textsuperscript{\ensuremath{p}}\fi}
\newunicodechar{ᴺ}{\ifmmode^{N}\else\textsuperscript{\ensuremath{N}}\fi}
\newunicodechar{ᶜ}{\ifmmode^{c}\else\textsuperscript{\ensuremath{c}}\fi}
\newunicodechar{ʰ}{\ifmmode^{h}\else\textsuperscript{\ensuremath{h}}\fi}
\newunicodechar{ᵠ}{\ifmmode^{q}\else\textsuperscript{\ensuremath{q}}\fi}
\newunicodechar{ₙ}{\ifmmode_{n}\else\textsubscript{\ensuremath{n}}\fi}
\newunicodechar{ₐ}{\ifmmode_{a}\else\textsubscript{\ensuremath{a}}\fi}
\newunicodechar{ᵢ}{\ifmmode_{i}\else\textsubscript{\ensuremath{i}}\fi}
\newunicodechar{ᵣ}{\ifmmode_{r}\else\textsubscript{\ensuremath{r}}\fi}
\newunicodechar{ₖ}{\ifmmode_{k}\else\textsubscript{\ensuremath{k}}\fi}
\newunicodechar{ᵦ}{\ifmmode_{\beta}\else\textsubscript{\ensuremath{\beta}}\fi}
\newunicodechar{ₓ}{\ifmmode_{x}\else\textsubscript{\ensuremath{x}}\fi}
\newfontfamily\popdisplay{ArchivoBlack-Regular.ttf}[Path=../../../fonts/]
\newfontfamily\poplogo{SpaceGrotesk-Bold.ttf}[Path=../../../fonts/]
\newfontfamily\popsemi{PlusJakartaSans-SemiBold.ttf}[Path=../../../fonts/]
\newfontfamily\popxbold{PlusJakartaSans-ExtraBold.ttf}[Path=../../../fonts/]
\newfontfamily\popxboldls{PlusJakartaSans-ExtraBold.ttf}[Path=../../../fonts/,LetterSpace=3.0]

\setlist[enumerate]{leftmargin=1.7em,itemsep=5pt,topsep=4pt}
\setlist[itemize]{leftmargin=1.7em,itemsep=4pt,topsep=4pt,label={\color{poporange}\textbullet}}
\setlength{\parindent}{0pt}
\setlength{\parskip}{5pt}
\renewcommand{\arraystretch}{1.25}

% pgfplots : style Pop par défaut
\pgfplotsset{
  every axis/.append style={axis line style={popink,line width=1pt},tick style={popink},
    grid style={popline,line width=0.5pt},label style={font=\small},tick label style={font=\footnotesize}},
  popcurve/.style={poporange,line width=1.8pt,no markers,smooth},
}
% Même style disponible dans un \draw TikZ simple (hors pgfplots)
\tikzset{popcurve/.style={poporange,line width=1.8pt}}

% ============================================================
% Pill / squiggle
% ============================================================
\newcommand{\poppill}[3]{%
  \tikz[baseline=-0.55ex]\node[draw=popink,line width=1.2pt,rounded corners=2.6pt,%
    fill=#2,inner xsep=6.5pt,inner ysep=3pt]{%
    \popxboldls\fontsize{7}{7}\selectfont\color{#3}\MakeUppercase{#1}};%
}
\newcommand{\popsquiggle}[1]{%
  \tikz[baseline=-0.4ex]{
    \draw[#1,line width=2.6pt,line cap=round]
      plot [smooth] coordinates {(0,0.09) (0.34,0.01) (0.68,0.21) (1.02,0.01) (1.36,0.10)};
  }%
}
% Mot-clé surligné (marqueur orange sous le texte)
\newcommand{\cle}[1]{{\popxbold\color{popdark}#1}}

% ============================================================
% En-tête / pied
% ============================================================
\pagestyle{fancy}
\fancyhf{}
\renewcommand{\headrulewidth}{0pt}
\renewcommand{\footrulewidth}{0pt}
\lhead{\raisebox{-0.15ex}{\poplogo\fontsize{13}{13}\selectfont\color{popink}OnBuch\color{poporange}+}}
\rhead{\poppill{\DOCMATIERE\ \textbullet\ \DOCNIVEAU\ \textbullet\ Leçon \DOCLECON}{white}{popink}}
\lfoot{\popxbold\fontsize{7.6}{7.6}\selectfont\color{popmuted}\MakeUppercase{Cours signé OnBuch+}\ \textbullet\ {\popsemi\DOCTITRE}}
\rfoot{\tikz[baseline=-0.6ex]\node[rounded corners=8.5pt,fill=popink,minimum height=17pt,minimum width=17pt,inner xsep=5pt,inner sep=0]{\popxbold\fontsize{8}{8}\selectfont\color{popcream}\thepage\,/\,\pageref*{LastPage}};}

% ============================================================
% Titres
% ============================================================
\newcommand{\chaptitre}[2]{%
  \begin{tcolorbox}[enhanced,colback=poporange,colframe=popink,boxrule=2.6pt,arc=11pt,
      drop shadow=popink,left=15pt,right=15pt,top=12pt,bottom=11pt,before skip=3pt,after skip=18pt]
    \poppill{#2}{popink}{popcream}\par\vspace{9pt}
    {\raggedright\hyphenpenalty=10000\exhyphenpenalty=10000\popdisplay\fontsize{22}{24}\selectfont\color{popcream}\MakeUppercase{#1}\par}
  \end{tcolorbox}
}
% Section numérotée : pastille orange avec le numéro + titre Archivo
\newcounter{popsec}
\newcommand{\coursec}[1]{%
  \stepcounter{popsec}\setcounter{popsub}{0}%
  \par\vspace{16pt}\noindent
  \begin{minipage}{\linewidth}
  \tikz[baseline=-0.7ex]\node[draw=popink,line width=1.6pt,fill=poporange,rounded corners=4pt,minimum size=22pt,inner sep=0]{\popdisplay\fontsize{12}{12}\selectfont\color{popcream}\thepopsec};%
  \hspace{8pt}{\popdisplay\fontsize{15}{17}\selectfont\color{popink}\MakeUppercase{#1}}\par
  \vspace{4pt}\noindent{\color{poporange}\rule{36pt}{3.6pt}}
  \end{minipage}\par\nopagebreak\vspace{8pt}
}
\newcounter{popsub}
\newcommand{\courssub}[1]{%
  \stepcounter{popsub}%
  \par\vspace{9pt}\noindent{\popxbold\fontsize{11.5}{13}\selectfont\color{popink}{\color{poporange}\thepopsec.\thepopsub}\hspace{6pt}#1}\par\nopagebreak\vspace{4pt}
}

% ============================================================
% Boîtes pédagogiques — même recette : fond blanc, bordure épaisse colorée,
% ombre dure encre, badge plein. Titre optionnel [#1].
% ============================================================
\tcbset{popbase/.style={breakable,enhanced,colback=white,boxrule=2.3pt,arc=9pt,
  shadow={3pt}{-3pt}{0pt}{popink},left=14pt,right=14pt,top=11pt,bottom=11pt,before skip=11pt,after skip=15pt,
  fontupper=\popsemi\fontsize{10.6}{14.5}\selectfont\color{popink},
  fontlower=\popsemi\fontsize{10.6}{14.5}\selectfont\color{popink}}}
% Coche dessinée (le glyphe ✓ n'existe pas dans les polices du document)
\newcommand{\popcheck}[1]{\tikz[baseline=-0.5ex]\draw[#1,line width=1.6pt,line cap=round,line join=round](-0.13,0.0)--(-0.04,-0.09)--(0.13,0.1);}
\newcommand{\popboxhead}[3]{\poppill{#1}{#2}{white}\ifx\relax#3\relax\else\hspace{6pt}{\popxbold\fontsize{10.6}{12}\selectfont\color{popink}#3}\fi\par\vspace{7pt}}

\newtcolorbox{aretenir}[1][]{popbase,colframe=popgreen,before upper={\popboxhead{À retenir}{popgreen}{#1}}}
\newtcolorbox{exemplebox}[1][]{popbase,colframe=poporange,before upper={\popboxhead{Exemple}{poporange}{#1}}}
\newtcolorbox{definition}[1][]{popbase,colframe=popblue,before upper={\popboxhead{Définition}{popblue}{#1}}}
\newtcolorbox{propriete}[1][]{popbase,colframe=poppurple,before upper={\popboxhead{Propriété}{poppurple}{#1}}}
\newtcolorbox{methode}[1][]{popbase,colframe=popgold,before upper={\popboxhead{Méthode}{popgold}{#1}}}
\newtcolorbox{attention}[1][]{popbase,colframe=poppink,before upper={\popboxhead{Attention}{poppink}{#1}}}
\newtcolorbox{experience}[1][]{popbase,colframe=popblue,colback=popblueL,before upper={\popboxhead{Expérience}{popblue}{#1}}}
% « Le savais-tu ? » : carte violette pleine (contexte, culture, Cameroun)
\newtcolorbox{savaistu}[1][]{popbase,colback=poppurple,colframe=popink,
  fontupper=\popsemi\fontsize{10.4}{14.2}\selectfont\color{white},
  before upper={\poppill{Le savais-tu\,?}{white}{poppurple}\ifx\relax#1\relax\else\hspace{6pt}{\popxbold\color{white}#1}\fi\par\vspace{7pt}}}
% Exercice résolu : énoncé en haut, \tcblower puis la solution sur fond crème
\newcounter{popexo}\newcounter{popexr}
\newtcolorbox{exoresolu}[1][]{popbase,colframe=poporange,colbacklower=popcream,
  segmentation style={popink,line width=1.2pt,dashed},
  before upper={\stepcounter{popexr}\popboxhead{Exercice résolu \thepopexr}{poporange}{#1}},
  before lower={\poppill{Solution}{popink}{popcream}\par\vspace{6pt}}}
% Exercice (non résolu en place) — la correction est dans la partie Corrigés
\newtcolorbox{exercice}[1][]{popbase,colframe=popink,
  before upper={\stepcounter{popexo}\popboxhead{Exercice \thepopexo}{popink}{#1}}}
% Corrigé
\newtcolorbox{corrige}[1][]{popbase,colframe=popgreen,colback=popgreenL,
  before upper={\popboxhead{Corrigé}{popgreen}{#1}}}
% Objectifs / prérequis / bilan
\newtcolorbox{objectifs}{popbase,colframe=popink,colback=poporangeL,
  before upper={\popboxhead{Objectifs de la leçon}{popink}{}}}
\newtcolorbox{prerequis}{popbase,colframe=popink,colback=popblueL,
  before upper={\popboxhead{Prérequis}{popblue}{}}}
\newtcolorbox{bilan}{popbase,colframe=popink,colback=white,boxrule=2.8pt,
  before upper={\popboxhead{Fiche bilan}{poporange}{L'essentiel à connaître pour l'examen}}}
% Figure encadrée avec légende
\newtcolorbox{popfigure}[1][]{enhanced,breakable=false,colback=white,colframe=popline,boxrule=1.4pt,arc=8pt,
  left=10pt,right=10pt,top=10pt,bottom=8pt,before skip=10pt,after skip=12pt,halign=center,
  lower separated=false,halign lower=center,
  fontlower=\popsemi\fontsize{9}{11.5}\selectfont\color{popmuted},
  #1}
\newcommand{\legende}[1]{\par\vspace{6pt}{\popsemi\fontsize{9}{11.5}\selectfont\color{popmuted}#1}}

% ============================================================
% Signature OnBuch+
% ============================================================
\newcommand{\poplogoinline}[1]{{\poplogo\fontsize{#1}{#1}\selectfont\color{popink}OnBuch\color{poporange}+}}
\newcommand{\signatureonbuch}{%
  \begin{tcolorbox}[enhanced,colback=popink,colframe=popink,boxrule=2.2pt,arc=10pt,
      drop shadow=poporange,left=14pt,right=14pt,top=10pt,bottom=10pt,before skip=0pt,after skip=0pt]
    \tikz[baseline=-0.7ex]\node[circle,fill=popgreen,draw=popcream,line width=1.3pt,minimum size=22pt,inner sep=0]{\popcheck{white}};%
    \hspace{9pt}\begin{minipage}[c]{0.62\linewidth}
      {\popxbold\fontsize{12}{13}\selectfont\color{popcream}Signé\ \textbullet\ {\poplogo OnBuch\color{poporange}+}}\par\vspace{2pt}
      {\raggedright\popsemi\fontsize{8}{10.5}\selectfont\color{popline}\MakeUppercase{Équipe pédagogique OnBuch+}\\\MakeUppercase{Conforme au programme officiel MINESEC}\par}
    \end{minipage}\hfill
    {\poplogo\fontsize{22}{22}\selectfont\color{popcream}OnBuch\color{poporange}+}
  \end{tcolorbox}%
}

% ============================================================
% Page de garde
% #1 titre · #2 matière · #3 niveau · #4 module · #5 n° leçon · #6 durée ·
% #7 description / objectifs (texte ou itemize)
% ============================================================
\newcommand{\pagegarde}[7]{%
  \thispagestyle{empty}
  \newgeometry{a4paper,margin=1.7cm,top=1.6cm,bottom=1.4cm}%
  \begin{tikzpicture}[remember picture,overlay]
    \fill[poporange!18] ([xshift=-3.2cm,yshift=-3.6cm]current page.north east) circle (4.6cm);
    \fill[poppurple!14] ([xshift=2.4cm,yshift=7.6cm]current page.south west) circle (3.8cm);
  \end{tikzpicture}%
  {\poplogo\fontsize{30}{30}\selectfont\color{popink}OnBuch\color{poporange}+}\hfill
  \raisebox{0.6ex}{\poppill{Cours signé \textbullet\ Premium}{popink}{popcream}}\par
  \vspace{1pt}\popsquiggle{poppurple}\par
  \vspace{8pt}
  \begin{tcolorbox}[enhanced,colback=poporange,colframe=popink,boxrule=2.8pt,arc=13pt,
      drop shadow=popink,left=18pt,right=18pt,top=12pt,bottom=13pt,after skip=10pt]
    \poppill{#2 \textbullet\ #3}{popink}{popcream}\hspace{5pt}\poppill{#4}{white}{popink}\par\vspace{8pt}
    {\popdisplay\fontsize{14}{14}\selectfont\color{popink}LEÇON #5}\par\vspace{4pt}
    {\raggedright\hyphenpenalty=10000\exhyphenpenalty=10000\popdisplay\fontsize{25}{27}\selectfont\color{popcream}\MakeUppercase{#1}\par}
    \vspace{8pt}
    {\popxbold\fontsize{9.5}{9.5}\selectfont\color{popink}\MakeUppercase{Durée conseillée : #6}}
  \end{tcolorbox}
  \begin{tcolorbox}[enhanced,colback=white,colframe=popink,boxrule=2.2pt,arc=10pt,
      drop shadow=popink,left=16pt,right=16pt,top=10pt,bottom=10pt,after skip=8pt]
    \poppill{Dans cette leçon}{poppurple}{white}\par\vspace{5pt}
    \popsemi\fontsize{9.3}{12.6}\selectfont\color{popink}#7
  \end{tcolorbox}
  \noindent\begin{minipage}[t]{0.487\linewidth}
    \begin{tcolorbox}[enhanced,colback=popgreen,colframe=popink,boxrule=2.2pt,arc=10pt,
        drop shadow=popink,left=11pt,right=11pt,top=10pt,bottom=10pt,height=3.1cm,valign=center]
      \tcbox[colback=white,colframe=popink,boxrule=1.3pt,arc=5pt,boxsep=0pt,nobeforeafter,
        left=3pt,right=3pt,top=3pt,bottom=3pt,valign=center]{\qrcode[height=1.9cm]{https://play.google.com/store/apps/details?id=com.luvvix.onbuch&hl=fr}}%
      \hspace{8pt}\begin{minipage}[c]{0.5\linewidth}
        {\popdisplay\fontsize{11}{12.5}\selectfont\color{white}\MakeUppercase{Télécharge OnBuch}}\par\vspace{4pt}
        {\raggedright\popsemi\fontsize{8.4}{11}\selectfont\color{white}Gratuit \textbullet\ Google Play\\Tuteur IA, annales, concours, résultats\par}
      \end{minipage}
    \end{tcolorbox}
  \end{minipage}\hfill
  \begin{minipage}[t]{0.487\linewidth}
    \begin{tcolorbox}[enhanced,colback=poppurple,colframe=popink,boxrule=2.2pt,arc=10pt,
        drop shadow=popink,left=11pt,right=11pt,top=10pt,bottom=10pt,height=3.1cm,valign=center]
      \tikz[baseline=-0.7ex]\node[circle,fill=white,draw=popink,line width=1.3pt,minimum size=1.6cm,inner sep=0]{\popdisplay\fontsize{24}{24}\selectfont\color{poppurple}+};%
      \hspace{8pt}\begin{minipage}[c]{0.6\linewidth}
        {\popdisplay\fontsize{11}{12.5}\selectfont\color{white}\MakeUppercase{Passe à OnBuch+}}\par\vspace{4pt}
        {\raggedright\popsemi\fontsize{8.4}{11}\selectfont\color{white}Tous les cours signés, Tuteur IA illimité, fiches PDF — onglet OnBuch+ de l'app\par}
      \end{minipage}
    \end{tcolorbox}
  \end{minipage}
  \vspace{8pt}
  \popcontenu
  \vfill
  \signatureonbuch
  \restoregeometry
  \newpage
}

% Rangée « Ce cours contient » (page de garde)
\newcommand{\poptile}[3]{% #1 couleur #2 gros texte #3 légende
  \begin{tcolorbox}[enhanced,colback=#1,colframe=popink,boxrule=2pt,arc=9pt,shadow={2.5pt}{-2.5pt}{0pt}{popink},
      left=6pt,right=6pt,top=9pt,bottom=9pt,halign=center,valign=center,height=1.75cm,width=\linewidth,nobeforeafter]
    {\popdisplay\fontsize{12.5}{13}\selectfont\color{white}\MakeUppercase{#2}}\par\vspace{3pt}
    {\centering\popsemi\fontsize{7.8}{9.5}\selectfont\color{white}#3\par}
  \end{tcolorbox}}
\newcommand{\popcontenu}{%
  \par\noindent{\popxboldls\fontsize{7.5}{7.5}\selectfont\color{popmuted}CE COURS SIGNÉ CONTIENT}\par\vspace{7pt}
  \noindent\begin{minipage}[t]{0.235\linewidth}\poptile{popblue}{Cours}{complet, illustré, pas à pas}\end{minipage}\hfill
  \begin{minipage}[t]{0.235\linewidth}\poptile{poporange}{Méthodes}{et exercices résolus}\end{minipage}\hfill
  \begin{minipage}[t]{0.235\linewidth}\poptile{poppink}{Exercices}{type examen + corrigés}\end{minipage}\hfill
  \begin{minipage}[t]{0.235\linewidth}\poptile{popgreen}{Fiche bilan}{l'essentiel à retenir}\end{minipage}\par}

% Sommaire
\newtcolorbox{sommaire}{enhanced,colback=white,colframe=popink,boxrule=2.2pt,arc=10pt,
  drop shadow=popink,left=18pt,right=18pt,top=13pt,bottom=13pt,before skip=6pt,after skip=20pt,
  before upper={\poppill{Sommaire}{poppurple}{white}\par\vspace{9pt}\popsemi\fontsize{10.6}{14.5}\selectfont\color{popink}}}

% Signature de fin de cours — poussée tout en bas de la dernière page
\newcommand{\findecours}{%
  \par\vfill\nopagebreak
  \begin{center}\popsquiggle{poporange}\end{center}\vspace{4pt}
  \signatureonbuch
}
\AtBeginDocument{\def\llbracket{\mathopen{[\mkern-3.5mu[}}\def\rrbracket{\mathclose{]\mkern-3.5mu]}}} % intervalles d'entiers (glyphe ⟦ absent de la police)
% Glyphes absents de Fira Math : redéfinis à partir de glyphes disponibles
\AtBeginDocument{%
  \def\setminus{\mathbin{\backslash}}\let\smallsetminus\setminus
  \def\vdots{\mathord{\vbox{\baselineskip3.2pt\lineskiplimit0pt\kern4pt\hbox{.}\hbox{.}\hbox{.}}}}%
  \def\ddots{\mathinner{\mkern1mu\raise6.5pt\hbox{.}\mkern2mu\raise3.6pt\hbox{.}\mkern2mu\raise0.7pt\hbox{.}\mkern1mu}}%
  \def\triangle{\mathord{\scalebox{0.9}{$\symup{\Delta}$}}}%
  \def\nabla{\mathord{\raisebox{\depth}{\scalebox{1}[-1]{$\symup{\Delta}$}}}}%
  \def\checkmark{\popcheck{popdark}}%
  \def\star{\mathbin{\ast}}\def\bigstar{\mathord{\ast}}%
  \def\diamond{\mathbin{\scalebox{0.6}{$\Diamond$}}}%
  \def\bigcup{\mathop{\mathchoice{\raisebox{-0.3em}{\scalebox{1.7}{$\cup$}}}{\scalebox{1.25}{$\cup$}}{\cup}{\cup}}\displaylimits}%
  \def\bigcap{\mathop{\mathchoice{\raisebox{-0.3em}{\scalebox{1.7}{$\cap$}}}{\scalebox{1.25}{$\cap$}}{\cap}{\cap}}\displaylimits}%
  \def\lesseqqgtr{\mathrel{\lessgtr}}\def\bigoplus{\mathop{\oplus}\displaylimits}%
}
\providecommand{\ssi}{si et seulement si} % abréviation fréquente
\AtBeginDocument{\providecommand{\wideparen}{\overparen}} % arc de cercle

%% FIGFIX-BEGIN v5 — correctifs globaux des figures (docs/onbuch-generation/tools/figfix)
\usepackage{adjustbox}\usepackage{etoolbox}\tcbuselibrary{raster}
% toute figure TikZ/pgfplots plus large que la ligne est réduite à la largeur disponible
\BeforeBeginEnvironment{tikzpicture}{\begin{adjustbox}{max width=\linewidth}}
\AfterEndEnvironment{tikzpicture}{\end{adjustbox}}
% légendes pgfplots lisibles par-dessus les courbes
\pgfplotsset{every axis/.append style={legend style={fill=white,fill opacity=0.92,text opacity=1,draw=black!15,inner sep=3pt}}}
% étiquettes d'arêtes (syntaxe edge["..."]) avec fond blanc
\tikzset{every edge quotes/.append style={fill=white,inner sep=1.2pt}}
%% FIGFIX-END
\begin{document}

\pagegarde{Géométrie analytique dans l'espace}{Mathématiques}{1ère E}{Géométrie dans l'espace}{N19}{10 h}{Cette leçon introduit les repères de l'espace et les coordonnées, puis traduit algébriquement droites et plans : représentations paramétriques, équations cartésiennes, positions relatives, vecteurs normaux et distances. Elle prolonge naturellement la géométrie vectorielle du plan et prépare les outils indispensables de la Terminale (produit scalaire dans l'espace, orthogonalité).
\vspace{4pt}
\begin{multicols}{2}\small
\begin{itemize}
\item À la fin de cette leçon, je sais reconnaître un repère de l'espace et y placer un point ou lire ses coordonnées
\item À la fin de cette leçon, je sais déterminer les coordonnées d'un vecteur dans une base de l'espace
\item À la fin de cette leçon, je sais reconnaître et écrire une représentation paramétrique d'une droite de l'espace et en extraire un point et un vecteur directeur
\item À la fin de cette leçon, je sais démontrer que deux droites de l'espace sont parallèles, sécantes, coplanaires ou non coplanaires
\item À la fin de cette leçon, je sais reconnaître et écrire une représentation paramétrique d'un plan et en extraire un point et deux vecteurs directeurs non colinéaires
\item À la fin de cette leçon, je sais déterminer une équation cartésienne d'un plan et un vecteur normal à ce plan
\end{itemize}\end{multicols}}

\begin{sommaire}
\begin{multicols}{2}
\begin{enumerate}
\item Repères et coordonnées dans l'espace
\item Représentations d'une droite de l'espace
\item Positions relatives de deux droites
\item Représentations paramétriques d'un plan
\item Équation cartésienne d'un plan et vecteur normal
\item Projetés orthogonaux et distances
\item Activité d'intégration
\item Méthodes et exercices résolus
\item Exercices
\item Corrigés des exercices
\item Fiche bilan
\end{enumerate}
\end{multicols}
\end{sommaire}

\begin{objectifs}
\begin{itemize}
\item À la fin de cette leçon, je sais reconnaître un repère de l'espace et y placer un point ou lire ses coordonnées
\item À la fin de cette leçon, je sais déterminer les coordonnées d'un vecteur dans une base de l'espace
\item À la fin de cette leçon, je sais reconnaître et écrire une représentation paramétrique d'une droite de l'espace et en extraire un point et un vecteur directeur
\item À la fin de cette leçon, je sais démontrer que deux droites de l'espace sont parallèles, sécantes, coplanaires ou non coplanaires
\item À la fin de cette leçon, je sais reconnaître et écrire une représentation paramétrique d'un plan et en extraire un point et deux vecteurs directeurs non colinéaires
\item À la fin de cette leçon, je sais déterminer une équation cartésienne d'un plan et un vecteur normal à ce plan
\item À la fin de cette leçon, je sais déterminer la position relative de deux plans (parallèles ou sécants)
\item À la fin de cette leçon, je sais déterminer le projeté orthogonal d'un point sur une droite ou sur un plan et calculer la distance correspondante
\end{itemize}
\end{objectifs}

\begin{prerequis}
\begin{itemize}
\item Vecteurs du plan : coordonnées, colinéarité, déterminant (Seconde)
\item Vecteurs de l'espace et règle du parallélépipède, caractérisation vectorielle des droites et plans (début de Première)
\item Repérage cartésien dans le plan et équations de droites du plan (Seconde)
\item Résolution de systèmes linéaires à deux ou trois inconnues
\item Positions relatives des droites et plans dans l'espace (approche géométrique de Seconde)
\end{itemize}
\end{prerequis}

\newpage

\coursec{Repères et coordonnées dans l'espace}

\textbf{Situation d'accroche.} À Garoua, un technicien de la SODECOTON doit installer une antenne parabolique au sommet d'un pylône. Avant de tendre le câble entre deux pylônes, il doit vérifier que ce câble ne croise pas la ligne électrique voisine, et calculer la distance minimale entre un ouvrier en hauteur et une paroi inclinée du bâtiment voisin. Comment décrire précisément, par des \emph{nombres}, la position d'un point dans l'espace ? Comment traduire un câble (une droite) ou un panneau (un plan) par des équations ?

Dans le plan, tu sais déjà qu'un repère $(O\,;\vec{\imath},\vec{\jmath})$ permet de repérer chaque point par deux nombres. Toute cette leçon repose sur la même idée, enrichie d'une dimension : ajouter une troisième direction pour sortir du plan. C'est la \cle{géométrie analytique de l'espace} : traduire les objets géométriques (points, droites, plans) en coordonnées et en équations, afin de remplacer des raisonnements de dessin par des calculs fiables. Dans cette première section, nous construisons l'outil de base : les repères de l'espace et les coordonnées.

\courssub{Base de l'espace et vecteurs coplanaires}

\textbf{Intuition.} Dans le plan, deux vecteurs non colinéaires suffisent à « engendrer » toutes les directions : tout vecteur du plan s'écrit comme combinaison de $\vec{\imath}$ et $\vec{\jmath}$. Dans l'espace, deux vecteurs ne suffisent plus : un vecteur vertical, par exemple, échappe au plan engendré par deux directions horizontales. Il faut une \emph{troisième} direction, qui ne soit pas « enfermée » dans le plan des deux premières. C'est exactement l'idée de la non-coplanarité.

\begin{definition}[Vecteurs coplanaires]
Trois vecteurs $\vec{u}$, $\vec{v}$ et $\vec{w}$ de l'espace sont dits \cle{coplanaires} s'il existe un point $O$ et trois points $A$, $B$, $C$ avec $\vec{u}=\overrightarrow{OA}$, $\vec{v}=\overrightarrow{OB}$, $\vec{w}=\overrightarrow{OC}$ tels que $O$, $A$, $B$, $C$ appartiennent à un même plan.

De façon équivalente : $\vec{u}$, $\vec{v}$, $\vec{w}$ sont coplanaires si et seulement si l'un d'eux est \cle{combinaison linéaire} des deux autres, c'est-à-dire s'il existe des réels $\alpha$ et $\beta$ tels que, par exemple,
\[ \vec{w} = \alpha\,\vec{u} + \beta\,\vec{v}. \]
\end{definition}

\begin{definition}[Base de l'espace]
Une \cle{base de l'espace} est un triplet $(\vec{\imath},\vec{\jmath},\vec{k})$ de trois vecteurs \textbf{non coplanaires}.

Pour toute base $(\vec{\imath},\vec{\jmath},\vec{k})$ et tout vecteur $\vec{u}$ de l'espace, il existe un \textbf{unique} triplet de réels $(x\,;y\,;z)$ tel que
\[ \vec{u} = x\,\vec{\imath} + y\,\vec{\jmath} + z\,\vec{k}. \]
Les réels $x$, $y$, $z$ sont les \cle{coordonnées} de $\vec{u}$ dans la base $(\vec{\imath},\vec{\jmath},\vec{k})$, et on note $\vec{u}\,(x\,;y\,;z)$ ou $\vec{u}\begin{pmatrix} x \\ y \\ z \end{pmatrix}$.
\end{definition}

\begin{exemplebox}[Tester la coplanarité]
Dans une base $(\vec{\imath},\vec{\jmath},\vec{k})$, on considère $\vec{u}\,(1\,;0\,;2)$, $\vec{v}\,(0\,;1\,;-1)$ et $\vec{w}\,(2\,;3\,;1)$. Cherchons $\alpha$ et $\beta$ tels que $\vec{w}=\alpha\vec{u}+\beta\vec{v}$ :
\[ \begin{cases} \alpha = 2 \\ \beta = 3 \\ 2\alpha - \beta = 1 \end{cases} \]
Les deux premières équations donnent $\alpha=2$, $\beta=3$, et la troisième est vérifiée : $2\times 2 - 3 = 1$. Donc $\vec{w}=2\vec{u}+3\vec{v}$ : les trois vecteurs sont \textbf{coplanaires} et ne forment \textbf{pas} une base de l'espace.
\end{exemplebox}

\courssub{Repère de l'espace et coordonnées d'un point}

\begin{definition}[Repère de l'espace]
Un \cle{repère de l'espace} est un quadruplet $(O\,;\vec{\imath},\vec{\jmath},\vec{k})$ où $O$ est un point de l'espace (l'\cle{origine}) et $(\vec{\imath},\vec{\jmath},\vec{k})$ une base de l'espace.

Pour tout point $M$ de l'espace, il existe un unique triplet $(x\,;y\,;z)$ de réels tel que
\[ \overrightarrow{OM} = x\,\vec{\imath} + y\,\vec{\jmath} + z\,\vec{k}. \]
On dit que $(x\,;y\,;z)$ sont les \cle{coordonnées} de $M$ dans le repère $(O\,;\vec{\imath},\vec{\jmath},\vec{k})$ : $x$ est l'\cle{abscisse}, $y$ l'\cle{ordonnée}, $z$ la \cle{cote}. On note $M\,(x\,;y\,;z)$.
\end{definition}

On distingue trois types de repères, du plus général au plus particulier :

\begin{definition}[Repère quelconque, orthogonal, orthonormal]
Soit $(O\,;\vec{\imath},\vec{\jmath},\vec{k})$ un repère de l'espace.
\begin{itemize}
\item C'est un repère \cle{quelconque} si $(\vec{\imath},\vec{\jmath},\vec{k})$ est une base sans autre condition ;
\item c'est un repère \cle{orthogonal} si les vecteurs $\vec{\imath}$, $\vec{\jmath}$, $\vec{k}$ sont deux à deux orthogonaux (les trois axes sont perpendiculaires deux à deux) ;
\item c'est un repère \cle{orthonormal} (ou orthonormé) s'il est orthogonal et si de plus $\|\vec{\imath}\|=\|\vec{\jmath}\|=\|\vec{k}\|=1$ (même unité de longueur sur les trois axes).
\end{itemize}
\end{definition}

\begin{attention}[Repère quelconque ou orthonormal ?]
L'erreur la plus fréquente consiste à utiliser des formules valables \emph{uniquement} en repère orthonormal (distance, et plus tard produit scalaire, orthogonalité) dans un repère quelconque. Retiens dès maintenant : \textbf{la formule de la distance n'est valable que dans un repère orthonormal}. Avant tout calcul de longueur, vérifie dans l'énoncé la mention « repère orthonormal ». En revanche, les coordonnées de $\overrightarrow{AB}$, du milieu, les critères de colinéarité et de coplanarité sont valables dans \emph{tout} repère.
\end{attention}

\textbf{Lecture géométrique des coordonnées.} Dire que $\overrightarrow{OM}=x\vec{\imath}+y\vec{\jmath}+z\vec{k}$, c'est dire que pour aller de $O$ à $M$, on avance de $x$ dans la direction $\vec{\imath}$, de $y$ dans la direction $\vec{\jmath}$, puis on « monte » de $z$ dans la direction $\vec{k}$. Le point $M$ est ainsi le sommet opposé à $O$ d'un \cle{parallélépipède} construit sur les trois axes : les trois coordonnées sont les longueurs (algébriques) de ses arêtes.

\begin{popfigure}
\begin{center}
\begin{tikzpicture}[scale=1.15, line join=round, >=Stealth]
% Vecteurs de base en perspective cavalière
\coordinate (O) at (0,0);
\coordinate (I) at (2.2,0);        % i : horizontal
\coordinate (J) at (1.05,0.75);    % j : fuite
\coordinate (K) at (0,1.9);        % k : vertical
% Point M(2 ; 3 ; 1.5) : coefficients 2,3,1.5
\coordinate (A) at ($(O)+2*(I)$);            % 2i
\coordinate (B) at ($(A)+3*(J)$);            % 2i+3j
\coordinate (M) at ($(B)+1.5*(K)$);          % +1.5k
\coordinate (C) at ($(O)+3*(J)$);            % 3j
\coordinate (D) at ($(O)+1.5*(K)$);          % 1.5k
\coordinate (E) at ($(A)+1.5*(K)$);          % 2i+1.5k
\coordinate (F) at ($(C)+1.5*(K)$);          % 3j+1.5k
\coordinate (G) at ($(A)+3*(J)$);            % = B
% Parallélépipède (arêtes)
\draw[popline] (O)--(A)--(B)--(C)--cycle;
\draw[popline] (D)--(E)--(M)--(F)--cycle;
\draw[popline] (A)--(E);
\draw[popline] (B)--(M);
\draw[popline] (C)--(F);
\draw[popline] (O)--(D);
% Axes
\draw[->, thick, popblue] (O)--($(I)+(0.6,0)$) node[below] {$\vec{\imath}$};
\draw[->, thick, popgreen] (O)--($(J)+(0.5,0.35)$) node[right] {$\vec{\jmath}$};
\draw[->, thick, poporange] (O)--($(K)+(0,0.6)$) node[left] {$\vec{k}$};
% Chemin de construction
\draw[->, very thick, poppurple] (O)--(A) node[midway, below] {$x=2$};
\draw[->, very thick, poppurple] (A)--(B) node[midway, below right] {$y=3$};
\draw[->, very thick, poppurple] (B)--(M) node[midway, right] {$z=1{,}5$};
% Vecteur OM
\draw[->, very thick, popink] (O)--(M) node[midway, above left] {$\overrightarrow{OM}$};
% Points
\fill[popdark] (O) circle (1.6pt) node[below left] {$O$};
\fill[popink] (M) circle (2pt) node[above right] {$M\,(2\,;3\,;1{,}5)$};
\fill[popmuted] (B) circle (1.4pt) node[below] {$M'$};
\end{tikzpicture}
\end{center}
\legende{Repère $(O\,;\vec{\imath},\vec{\jmath},\vec{k})$ en perspective cavalière. Le point $M\,(2\,;3\,;1{,}5)$ est construit par le parallélépipède : on reporte $x=2$ selon $\vec{\imath}$, $y=3$ selon $\vec{\jmath}$ (point $M'$, projeté de $M$ sur le plan de base), puis on monte de $z=1{,}5$ selon $\vec{k}$.}
\end{popfigure}

\begin{methode}[Placer un point dans un repère en perspective cavalière]
Pour placer $M\,(x\,;y\,;z)$ dans un repère $(O\,;\vec{\imath},\vec{\jmath},\vec{k})$ dessiné en perspective cavalière :
\begin{enumerate}
\item Dans le \textbf{plan de base} $(O\,;\vec{\imath},\vec{\jmath})$ (le plan « horizontal » du dessin), place le point $M'\,(x\,;y\,;0)$ : reporte $x$ sur l'axe $(O\,;\vec{\imath})$, puis $y$ parallèlement à $\vec{\jmath}$ ;
\item à partir de $M'$, trace la parallèle à l'axe $(O\,;\vec{k})$ (la « verticale » du dessin) et reporte $z$ : si $z>0$ on monte, si $z<0$ on descend ;
\item le point obtenu est $M$. Matérialise le parallélépipède en pointillés pour rendre la lecture facile.
\end{enumerate}
\end{methode}

\courssub{Coordonnées d'un vecteur et opérations}

\begin{propriete}[Coordonnées de $\overrightarrow{AB}$]
Dans un repère $(O\,;\vec{\imath},\vec{\jmath},\vec{k})$, soient $A\,(x_A\,;y_A\,;z_A)$ et $B\,(x_B\,;y_B\,;z_B)$. Alors
\[ \overrightarrow{AB}\begin{pmatrix} x_B - x_A \\ y_B - y_A \\ z_B - z_A \end{pmatrix}. \]
\end{propriete}

\begin{experience}[Démonstration guidée]
On part de la relation de Chasles, valable dans l'espace comme dans le plan :
\begin{enumerate}
\item Écris $\overrightarrow{AB} = \overrightarrow{AO} + \overrightarrow{OB} = \overrightarrow{OB} - \overrightarrow{OA}$ ;
\item remplace : $\overrightarrow{OA} = x_A\vec{\imath}+y_A\vec{\jmath}+z_A\vec{k}$ et $\overrightarrow{OB} = x_B\vec{\imath}+y_B\vec{\jmath}+z_B\vec{k}$ ;
\item soustrais composante par composante : $\overrightarrow{AB} = (x_B-x_A)\vec{\imath} + (y_B-y_A)\vec{\jmath} + (z_B-z_A)\vec{k}$ ;
\item par \textbf{unicité} des coordonnées dans une base, on obtient le résultat annoncé. $\square$
\end{enumerate}
\end{experience}

\begin{propriete}[Opérations sur les coordonnées]
Dans une base $(\vec{\imath},\vec{\jmath},\vec{k})$, soient $\vec{u}\,(x\,;y\,;z)$, $\vec{v}\,(x'\,;y'\,;z')$ et $\lambda$ un réel. Alors :
\[ \vec{u}+\vec{v}\,\begin{pmatrix} x+x' \\ y+y' \\ z+z' \end{pmatrix}, \qquad \lambda\vec{u}\,\begin{pmatrix} \lambda x \\ \lambda y \\ \lambda z \end{pmatrix}. \]
En particulier, $\vec{u}=\vec{v}$ si et seulement si $x=x'$, $y=y'$ et $z=z'$ ; et $\vec{u}=\vec{0}$ si et seulement si $x=y=z=0$.
\end{propriete}

\begin{propriete}[Coordonnées du milieu d'un segment]
Dans un repère quelconque, le milieu $I$ du segment $[AB]$, avec $A\,(x_A\,;y_A\,;z_A)$ et $B\,(x_B\,;y_B\,;z_B)$, a pour coordonnées
\[ I\left(\dfrac{x_A+x_B}{2}\,;\,\dfrac{y_A+y_B}{2}\,;\,\dfrac{z_A+z_B}{2}\right). \]
\end{propriete}

En effet, $I$ milieu de $[AB]$ équivaut à $\overrightarrow{OI}=\dfrac{1}{2}\left(\overrightarrow{OA}+\overrightarrow{OB}\right)$, et il suffit d'appliquer les opérations sur les coordonnées.

\courssub{Colinéarité et coplanarité en coordonnées}

\textbf{Intuition.} Deux vecteurs colinéaires ont « la même direction » : l'un est un agrandissement ou une réduction de l'autre. En coordonnées, cela se traduit par la \emph{proportionnalité} des trois composantes.

\begin{propriete}[Critère de colinéarité]
Dans une base de l'espace, deux vecteurs $\vec{u}\,(x\,;y\,;z)$ et $\vec{v}\,(x'\,;y'\,;z')$ sont \cle{colinéaires} si et seulement si leurs coordonnées sont proportionnelles, c'est-à-dire s'il existe un réel $k$ tel que
\[ x' = kx, \qquad y' = ky, \qquad z' = kz. \]
\end{propriete}

\begin{experience}[Démonstration guidée du critère de colinéarité]
Supposons $\vec{u}\neq\vec{0}$.
\begin{enumerate}
\item \textbf{Sens direct.} Par définition, $\vec{u}$ et $\vec{v}$ colinéaires signifie qu'il existe un réel $k$ tel que $\vec{v}=k\vec{u}$. En coordonnées : $k\vec{u}\,(kx\,;ky\,;kz)$, donc $x'=kx$, $y'=ky$, $z'=kz$ par unicité des coordonnées.
\item \textbf{Réciproque.} S'il existe $k$ tel que $x'=kx$, $y'=ky$, $z'=kz$, alors $\vec{v}=kx\vec{\imath}+ky\vec{\jmath}+kz\vec{k}=k(x\vec{\imath}+y\vec{\jmath}+z\vec{k})=k\vec{u}$, donc $\vec{u}$ et $\vec{v}$ sont colinéaires. $\square$
\end{enumerate}
\end{experience}

\begin{attention}[Le « produit en croix » du plan ne suffit plus]
Dans le plan, un seul déterminant $xy'-x'y$ testait la colinéarité. Dans l'espace, il faut comparer \textbf{trois} rapports $\dfrac{x'}{x}$, $\dfrac{y'}{y}$, $\dfrac{z'}{z}$ (lorsque les dénominateurs sont non nuls) et vérifier qu'ils sont \emph{tous égaux}. Si un seul rapport diffère, les vecteurs ne sont pas colinéaires. Attention aussi aux coordonnées nulles : si $x=0$ et $x'=0$, la condition sur la première composante est automatiquement satisfaite, quel que soit $k$.
\end{attention}

\begin{propriete}[Critère de coplanarité]
Trois vecteurs $\vec{u}$, $\vec{v}$, $\vec{w}$ sont coplanaires si et seulement si l'un d'eux est combinaison linéaire des deux autres, c'est-à-dire s'il existe des réels $\alpha$ et $\beta$ tels que $\vec{w}=\alpha\vec{u}+\beta\vec{v}$ (ou une relation analogue par permutation). En pratique, on résout le système de trois équations à deux inconnues obtenu en égalant les coordonnées : s'il admet une solution, les vecteurs sont coplanaires ; sinon, ils forment une base de l'espace.
\end{propriete}

\courssub{Distance entre deux points en repère orthonormal}

\begin{propriete}[Distance dans un repère orthonormal — admise]
Soit $(O\,;\vec{\imath},\vec{\jmath},\vec{k})$ un repère \textbf{orthonormal} de l'espace. Pour $A\,(x_A\,;y_A\,;z_A)$ et $B\,(x_B\,;y_B\,;z_B)$ :
\[ AB = \sqrt{(x_B-x_A)^2 + (y_B-y_A)^2 + (z_B-z_A)^2}. \]
Cette formule est \textbf{admise} à ce stade : elle sera démontrée en Terminale à l'aide du produit scalaire (c'est la généralisation du théorème de Pythagore à l'espace, appliqué deux fois : d'abord dans le plan de base, puis dans un plan vertical).
\end{propriete}

\begin{aretenir}[Ce qui est valable dans tout repère, ce qui ne l'est pas]
\begin{itemize}
\item \textbf{Dans tout repère} : coordonnées de $\overrightarrow{AB}$, du milieu de $[AB]$, opérations sur les vecteurs, critères de colinéarité et de coplanarité ;
\item \textbf{uniquement en repère orthonormal} : la formule de la distance $AB=\sqrt{(x_B-x_A)^2+(y_B-y_A)^2+(z_B-z_A)^2}$.
\end{itemize}
\end{aretenir}

\begin{exoresolu}[Un exemple complet en repère orthonormal]
L'espace est muni d'un repère orthonormal $(O\,;\vec{\imath},\vec{\jmath},\vec{k})$. On donne $A\,(1\,;-2\,;3)$ et $B\,(4\,;0\,;-1)$. Déterminer les coordonnées de $\overrightarrow{AB}$, les coordonnées du milieu $I$ de $[AB]$, et la longueur $AB$.
\tcblower
\textbf{Coordonnées de $\overrightarrow{AB}$.} On applique la formule « coordonnées de $B$ moins coordonnées de $A$ » :
\[ \overrightarrow{AB}\begin{pmatrix} 4-1 \\ 0-(-2) \\ -1-3 \end{pmatrix} = \begin{pmatrix} 3 \\ 2 \\ -4 \end{pmatrix}. \]
\textbf{Milieu de $[AB]$.} On prend la demi-somme des coordonnées :
\[ I\left(\dfrac{1+4}{2}\,;\,\dfrac{-2+0}{2}\,;\,\dfrac{3+(-1)}{2}\right) \quad\text{donc}\quad I\left(\dfrac{5}{2}\,;\,-1\,;\,1\right). \]
\textbf{Longueur $AB$.} Le repère est orthonormal, la formule de la distance est donc licite :
\begin{align*}
AB &= \sqrt{(4-1)^2+(0-(-2))^2+(-1-3)^2} \\
&= \sqrt{3^2+2^2+(-4)^2} \\
&= \sqrt{9+4+16} = \sqrt{29}.
\end{align*}
Comme $25<29<36$, on a $5<\sqrt{29}<6$ : $AB \approx 5{,}4$ unités. Si l'unité est le mètre, le câble du technicien mesure environ $\SI{5,4}{m}$.
\end{exoresolu}

\begin{exercice}[À toi de jouer]
Dans un repère orthonormal $(O\,;\vec{\imath},\vec{\jmath},\vec{k})$, on donne $P\,(2\,;1\,;-3)$, $Q\,(-1\,;3\,;0)$ et les vecteurs $\vec{u}\,(2\,;-1\,;4)$, $\vec{v}\,(-3\,;1{,}5\,;-6)$, $\vec{w}\,(1\,;0\,;2)$.
\begin{enumerate}
\item Déterminer les coordonnées de $\overrightarrow{PQ}$, du milieu $J$ de $[PQ]$, et calculer $PQ$.
\item Les vecteurs $\vec{u}$ et $\vec{v}$ sont-ils colinéaires ?
\item Les vecteurs $\vec{u}$, $\vec{v}$, $\vec{w}$ sont-ils coplanaires ?
\end{enumerate}
\end{exercice}

\begin{corrige}[Exercice 1]
\textbf{1.} $\overrightarrow{PQ}\begin{pmatrix} -1-2 \\ 3-1 \\ 0-(-3) \end{pmatrix} = \begin{pmatrix} -3 \\ 2 \\ 3 \end{pmatrix}$. Le milieu : $J\left(\dfrac{2-1}{2}\,;\,\dfrac{1+3}{2}\,;\,\dfrac{-3+0}{2}\right)$, soit $J\left(\dfrac{1}{2}\,;\,2\,;\,-\dfrac{3}{2}\right)$. Distance : $PQ=\sqrt{(-3)^2+2^2+3^2}=\sqrt{9+4+9}=\sqrt{22}$.

\textbf{2.} On compare les rapports : $\dfrac{-3}{2}=-1{,}5$ ; $\dfrac{1{,}5}{-1}=-1{,}5$ ; $\dfrac{-6}{4}=-1{,}5$. Les trois rapports sont égaux : $\vec{v}=-1{,}5\,\vec{u}$, donc $\vec{u}$ et $\vec{v}$ sont \textbf{colinéaires}.

\textbf{3.} Comme $\vec{v}=-1{,}5\,\vec{u}$, on a $\vec{v}=(-1{,}5)\vec{u}+0\,\vec{w}$ : $\vec{v}$ est combinaison linéaire de $\vec{u}$ et $\vec{w}$, donc les trois vecteurs sont \textbf{coplanaires}.
\end{corrige}

\begin{savaistu}[Descartes et la géométrie qui calcule]
C'est René Descartes qui, en 1637 dans \emph{La Géométrie}, a eu l'idée révolutionnaire de repérer les points par des nombres, transformant la géométrie en calcul algébrique. L'extension à l'espace, développée au XVIII\textsuperscript{e} siècle, est aujourd'hui omniprésente : le GPS de ton téléphone te localise par trois coordonnées (latitude, longitude, altitude), et les logiciels de conception utilisés par les ingénieurs de la SODECOTON ou d'ENEO représentent chaque pylône, chaque câble, chaque panneau par des coordonnées et des équations — exactement les outils que tu viens de découvrir.
\end{savaistu}

\begin{aretenir}[L'essentiel de la section]
\begin{itemize}
\item Trois vecteurs sont \cle{coplanaires} si l'un est combinaison linéaire des deux autres ; une \cle{base} de l'espace est un triplet de vecteurs non coplanaires ;
\item un \cle{repère} $(O\,;\vec{\imath},\vec{\jmath},\vec{k})$ associe à tout point $M$ un unique triplet $(x\,;y\,;z)$ tel que $\overrightarrow{OM}=x\vec{\imath}+y\vec{\jmath}+z\vec{k}$ ;
\item $\overrightarrow{AB}\,(x_B-x_A\,;\,y_B-y_A\,;\,z_B-z_A)$ et le milieu de $[AB]$ a pour coordonnées les demi-sommes ;
\item deux vecteurs sont colinéaires si et seulement si leurs coordonnées sont proportionnelles ;
\item en repère \textbf{orthonormal} seulement : $AB=\sqrt{(x_B-x_A)^2+(y_B-y_A)^2+(z_B-z_A)^2}$.
\end{itemize}
\end{aretenir}

\coursec{Représentations d'une droite de l'espace}

Dans la section précédente, tu as appris à repérer un point de l'espace par ses trois coordonnées $(x\,;y\,;z)$ dans un repère $(O\,;\vec{\imath},\vec{\jmath},\vec{k})$, et à calculer les coordonnées d'un vecteur. Nous allons maintenant utiliser ces outils pour décrire les objets géométriques les plus simples de l'espace : les droites. L'idée directrice est la même qu'en Seconde dans le plan, mais avec une coordonnée de plus : une droite sera décrite par des \emph{équations} portant sur les coordonnées de ses points. C'est toute la puissance de la géométrie analytique : traduire la géométrie en calcul.

\courssub{Caractérisation vectorielle d'une droite}

\textbf{Intuition.} Imagine le câble d'une ligne électrique qui part d'un poteau $A$ et file tout droit dans une direction fixe. Pour décrire ce câble, il suffit de deux informations : un point de départ $A$, et une direction, donnée par un vecteur $\vec{u}$ non nul. Tout point $M$ du câble est alors atteint en partant de $A$ et en avançant d'un certain nombre de fois le vecteur $\vec{u}$ (en avant ou en arrière).

\begin{definition}[Droite de l'espace définie par un point et un vecteur directeur]
Soit $A$ un point de l'espace et $\vec{u}$ un vecteur \cle{non nul}. La \cle{droite} passant par $A$ et de \cle{vecteur directeur} $\vec{u}$, notée $(A\,;\vec{u})$, est l'ensemble des points $M$ de l'espace tels que $\overrightarrow{AM}$ et $\vec{u}$ soient \cle{colinéaires} :
\[ M \in (A\,;\vec{u}) \iff \text{il existe un réel } t \text{ tel que } \overrightarrow{AM} = t\,\vec{u}. \]
\end{definition}

\begin{attention}[Le vecteur directeur doit être non nul]
La condition $\vec{u} \neq \vec{0}$ est indispensable : si $\vec{u} = \vec{0}$, la relation $\overrightarrow{AM} = t\vec{u}$ donnerait $\overrightarrow{AM} = \vec{0}$, c'est-à-dire $M = A$ pour tout $t$ : on décrirait un seul point, pas une droite. De même, une droite est définie par \emph{deux points distincts} : si $A = B$, il n'y a pas de droite $(AB)$.
\end{attention}

\begin{popfigure}
\begin{tikzpicture}[scale=1.05,>=Stealth]
  % Repère
  \draw[->,popmuted] (0,0) -- (1.1,0) node[below left] {$\vec{\imath}$};
  \draw[->,popmuted] (0,0) -- (0,1.1) node[left] {$\vec{k}$};
  \draw[->,popmuted] (0,0) -- (-0.65,-0.55) node[below] {$\vec{\jmath}$};
  \fill[popdark] (0,0) circle (1.4pt) node[above left] {$O$};
  % Droite
  \draw[popblue,very thick] (-0.6,0.4) -- (7.2,3.6) node[right] {$(d)$};
  % Point A
  \fill[poporange] (1.6,1.2) circle (2pt) node[below left] {$A$};
  % Vecteur directeur u
  \draw[->,poporange,very thick] (1.6,1.2) -- (3.2,1.8) node[midway,below] {$\vec{u}$};
  % Point M(t)
  \fill[poppurple] (5.6,3.2) circle (2pt) node[above] {$M(t)$};
  % Flèche AM = t u
  \draw[->,poppurple,thick,dashed] (1.6,1.2) -- (5.6,3.2) node[midway,above left] {$\overrightarrow{AM}=t\,\vec{u}$};
\end{tikzpicture}
\legende{Une droite $(d)$ de l'espace est entièrement déterminée par un point $A$ et un vecteur directeur $\vec{u}$. Chaque point $M$ de la droite correspond à une valeur du paramètre réel $t$.}
\end{popfigure}

\courssub{Représentation paramétrique d'une droite}

Plaçons-nous dans un repère $(O\,;\vec{\imath},\vec{\jmath},\vec{k})$ de l'espace. Soit $A(x_A\,;y_A\,;z_A)$ un point et $\vec{u}(a\,;b\,;c)$ un vecteur non nul. Un point $M(x\,;y\,;z)$ appartient à la droite $(A\,;\vec{u})$ si et seulement s'il existe un réel $t$ tel que $\overrightarrow{AM} = t\vec{u}$. Or $\overrightarrow{AM}$ a pour coordonnées $(x - x_A\,;y - y_A\,;z - z_A)$. L'égalité vectorielle se traduit donc, coordonnée par coordonnée, par trois égalités.

\begin{definition}[Représentation paramétrique d'une droite de l'espace]
Dans un repère de l'espace, la droite passant par $A(x_A\,;y_A\,;z_A)$ et de vecteur directeur $\vec{u}(a\,;b\,;c)$ (avec $\vec{u} \neq \vec{0}$, c'est-à-dire $a$, $b$, $c$ non tous nuls) admet la \cle{représentation paramétrique} :
\[ \begin{cases} x = x_A + at \\ y = y_A + bt \\ z = z_A + ct \end{cases}, \quad t \in \mathbb{R}. \]
Le réel $t$ est appelé le \cle{paramètre}. À chaque valeur de $t$ correspond un unique point de la droite, et réciproquement.
\end{definition}

\begin{exemplebox}[Une première représentation paramétrique]
Soit la droite $(d)$ passant par $A(2\,;-1\,;0)$ et dirigée par $\vec{u}(1\,;3\,;-2)$. Une représentation paramétrique de $(d)$ est :
\[ \begin{cases} x = 2 + t \\ y = -1 + 3t \\ z = -2t \end{cases}, \quad t \in \mathbb{R}. \]
Pour $t = 0$, on retrouve le point $A$. Pour $t = 2$, on obtient le point $(4\,;5\,;-4)$. Pour $t = -1$, on obtient $(1\,;-4\,;2)$. Chaque valeur de $t$ « fabrique » un point de la droite.
\end{exemplebox}

\begin{methode}[Lecture inverse : reconnaître une représentation paramétrique de droite]
Face à un système de la forme
\[ \begin{cases} x = 3 - 2t \\ y = 1 + t \\ z = -4 + 5t \end{cases}, \quad t \in \mathbb{R}, \]
on procède ainsi :
\begin{enumerate}
\item On vérifie que $x$, $y$ et $z$ sont des expressions \emph{affines} du \emph{même} paramètre $t$ (degré 1 en $t$) ;
\item On vérifie que les coefficients de $t$ ne sont pas tous nuls (ici $-2$, $1$, $5$) : sinon ce serait un point ;
\item Le \textbf{point} de la droite se lit dans les \emph{constantes} : $A(3\,;1\,;-4)$ (c'est le point obtenu pour $t = 0$) ;
\item Le \textbf{vecteur directeur} se lit dans les \emph{coefficients de $t$} : $\vec{u}(-2\,;1\,;5)$.
\end{enumerate}
\end{methode}

\begin{attention}[Attention aux paramètres « déguisés »]
Le paramètre peut s'appeler $k$, $\lambda$, $s$... et le point « visible » n'est pas toujours celui qu'on croit. Par exemple, le système
\[ \begin{cases} x = 5 + 2k \\ y = 3 - k \\ z = 1 + 4k \end{cases} \]
décrit la droite passant par $A(5\,;3\,;1)$ dirigée par $\vec{u}(2\,;-1\,;4)$. Le système $\begin{cases} x = 3 + 2s \\ y = 4 - s \\ z = -3 + 4s \end{cases}$ décrit \emph{la même droite} : le vecteur directeur est le même, et le point $(3\,;4\,;-3)$ appartient bien à la première droite (il correspond à $k = -1$). Apprends à lire les coefficients et à \emph{vérifier} l'appartenance des points, pas à mémoriser des formes.
\end{attention}

\begin{propriete}[Une droite admet une infinité de représentations paramétriques]
Une même droite $(d)$ admet une infinité de représentations paramétriques : on peut changer le point de référence (n'importe quel point de $(d)$ convient) et remplacer le vecteur directeur par n'importe quel vecteur non nul qui lui est colinéaire.
\end{propriete}

\begin{experience}[Démonstration guidée]
Soit $(d)$ la droite de représentation paramétrique $\begin{cases} x = x_A + at \\ y = y_A + bt \\ z = z_A + ct \end{cases}$, $t \in \mathbb{R}$, avec $A(x_A\,;y_A\,;z_A)$ et $\vec{u}(a\,;b\,;c)$ non nul.
\begin{enumerate}
\item \textbf{Changement de vecteur directeur.} Soit $\vec{v} = \lambda \vec{u}$ avec $\lambda \neq 0$. Montre que le système $\begin{cases} x = x_A + \lambda a\, s \\ y = y_A + \lambda b\, s \\ z = z_A + \lambda c\, s \end{cases}$, $s \in \mathbb{R}$, décrit exactement les mêmes points. \emph{Indication : pose $t = \lambda s$ ; quand $s$ décrit $\mathbb{R}$, $t$ décrit aussi $\mathbb{R}$ car $\lambda \neq 0$.}
\item \textbf{Changement de point.} Soit $B$ un point de $(d)$, obtenu pour la valeur $t_0$ du paramètre : $B$ a pour coordonnées $(x_A + a t_0\,;y_A + b t_0\,;z_A + c t_0)$. Montre que le système $\begin{cases} x = x_B + a\,s \\ y = y_B + b\,s \\ z = z_B + c\,s \end{cases}$ décrit les mêmes points. \emph{Indication : pose $t = t_0 + s$.}
\item \textbf{Conclusion.} En combinant les deux transformations, explique pourquoi $(d)$ admet une infinité de représentations paramétriques.
\end{enumerate}
\end{experience}

\begin{exemplebox}[Deux écritures pour une même droite]
La droite $(d)$ : $\begin{cases} x = 2 + t \\ y = -1 + 3t \\ z = -2t \end{cases}$ peut aussi s'écrire $\begin{cases} x = 4 + s \\ y = 5 + 3s \\ z = -4 - 2s \end{cases}$ (on a pris le point obtenu pour $t = 2$), ou encore $\begin{cases} x = 2 - 2k \\ y = -1 - 6k \\ z = 4k \end{cases}$ (vecteur directeur $(-2\,;-6\,;4) = -2\vec{u}$). Ces trois systèmes décrivent la \emph{même} droite.
\end{exemplebox}

\courssub{Droite définie par deux points distincts}

\begin{propriete}[Vecteur directeur $\overrightarrow{AB}$]
Soient $A$ et $B$ deux points \emph{distincts} de l'espace. La droite $(AB)$ est la droite passant par $A$ et de vecteur directeur $\overrightarrow{AB}$. Si $A(x_A\,;y_A\,;z_A)$ et $B(x_B\,;y_B\,;z_B)$, alors $\overrightarrow{AB}(x_B - x_A\,;y_B - y_A\,;z_B - z_A)$ et une représentation paramétrique de $(AB)$ est :
\[ \begin{cases} x = x_A + (x_B - x_A)\,t \\ y = y_A + (y_B - y_A)\,t \\ z = z_A + (z_B - z_A)\,t \end{cases}, \quad t \in \mathbb{R}. \]
Pour $t = 0$ on obtient $A$, pour $t = 1$ on obtient $B$. Remarque au passage : les points du \emph{segment} $[AB]$ correspondent exactement aux valeurs de $t$ dans $[0\,;1]$.
\end{propriete}

\begin{exemplebox}[Droite passant par deux points]
Soient $A(1\,;0\,;2)$ et $B(3\,;2\,;-1)$. On calcule $\overrightarrow{AB}(2\,;2\,;-3)$. Une représentation paramétrique de $(AB)$ est :
\[ \begin{cases} x = 1 + 2t \\ y = 2t \\ z = 2 - 3t \end{cases}, \quad t \in \mathbb{R}. \]
Vérifie : pour $t = 1$, on retombe bien sur $B(3\,;2\,;-1)$.
\end{exemplebox}

\courssub{Appartenance d'un point à une droite}

\begin{methode}[Vérifier qu'un point appartient à une droite]
Pour savoir si un point $M(x_M\,;y_M\,;z_M)$ appartient à la droite de représentation paramétrique $\begin{cases} x = x_A + at \\ y = y_A + bt \\ z = z_A + ct \end{cases}$ :
\begin{enumerate}
\item On remplace $x$, $y$, $z$ par les coordonnées de $M$ dans les \textbf{trois} équations ;
\item On résout \textbf{une} équation (la plus simple) pour trouver une valeur candidate $t_0$ de $t$ ;
\item On vérifie que \textbf{cette même valeur} $t_0$ satisfait les \textbf{deux autres} équations ;
\item Si oui : $M \in (d)$ (et $t_0$ est le paramètre de $M$). Si l'une des deux vérifications échoue : $M \notin (d)$.
\end{enumerate}
\end{methode}

\begin{exoresolu}[Le point $C$ appartient-il à la droite $(d)$ ?]
On considère la droite $(d)$ passant par $A(2\,;-1\,;0)$ et dirigée par $\vec{u}(1\,;3\,;-2)$. Le point $C(4\,;5\,;-4)$ appartient-il à $(d)$ ?
\tcblower
Une représentation paramétrique de $(d)$ est $\begin{cases} x = 2 + t \\ y = -1 + 3t \\ z = -2t \end{cases}$, $t \in \mathbb{R}$.

$C \in (d)$ si et seulement s'il existe un réel $t$ tel que :
\[ \begin{cases} 4 = 2 + t \\ 5 = -1 + 3t \\ -4 = -2t \end{cases} \]
La première équation donne $t = 2$. Vérifions avec cette même valeur :
\begin{itemize}
\item deuxième équation : $-1 + 3 \times 2 = 5$ \checkmark ;
\item troisième équation : $-2 \times 2 = -4$ \checkmark.
\end{itemize}
Les trois équations sont satisfaites pour $t = 2$, donc $C \in (d)$ : c'est le point de paramètre $t = 2$.

\emph{Contre-exemple :} soit $D(5\,;8\,;-4)$. La première équation donnerait $t = 3$, la deuxième $t = 3$, mais la troisième donne $t = 2$. Les valeurs ne coïncident pas : $D \notin (d)$.
\end{exoresolu}

\begin{attention}[Le piège classique : tester avec des valeurs de $t$ différentes]
L'erreur la plus fréquente consiste à conclure trop vite. Deux rédactions fautives :
\begin{itemize}
\item Trouver $t = 2$ avec l'équation en $x$, constater que $t = 2$ convient pour $y$, et conclure « le point appartient » \emph{sans tester} l'équation en $z$ ;
\item Résoudre les trois équations séparément, obtenir trois valeurs de $t$, et conclure « le point appartient car chaque équation a une solution ».
\end{itemize}
La bonne exigence est : \textbf{un seul et même réel $t$ doit satisfaire simultanément les trois équations}. Rédige toujours : « la première équation donne $t = \ldots$ ; reportons dans les deux autres... ».
\end{attention}

\courssub{Système d'équations cartésiennes d'une droite}

\textbf{Intuition.} Une représentation paramétrique décrit une droite « en la parcourant » grâce au paramètre $t$. On peut aussi décrire une droite « de l'extérieur », comme l'intersection de deux plans qui se coupent : pense à l'arête d'une case au village, là où deux murs se rencontrent, ou à la ligne d'intersection du toit et du mur. Cette idée mène aux \emph{équations cartésiennes} d'une droite.

\begin{methode}[Passer d'une représentation paramétrique à un système d'équations cartésiennes]
On part de $\begin{cases} x = x_A + at \\ y = y_A + bt \\ z = z_A + ct \end{cases}$ et on \emph{élimine} le paramètre $t$ :
\begin{enumerate}
\item On exprime $t$ dans deux des trois équations (celles où le coefficient de $t$ est non nul) ;
\item On égalise les expressions obtenues (double égalité) ;
\item On réécrit le résultat sous forme d'un système de deux équations cartésiennes en $x$, $y$, $z$ : chacune est l'équation d'un plan, et la droite est l'intersection de ces deux plans sécants.
\end{enumerate}
\end{methode}

\begin{exemplebox}[De la forme paramétrique à la forme cartésienne]
Reprenons $(d)$ : $\begin{cases} x = 2 + t \\ y = -1 + 3t \\ z = -2t \end{cases}$.
La première équation donne $t = x - 2$ ; la deuxième donne $t = \dfrac{y + 1}{3}$ ; la troisième donne $t = -\dfrac{z}{2}$. D'où la double égalité :
\[ x - 2 = \frac{y+1}{3} = -\frac{z}{2}. \]
On peut la réécrire comme système de deux équations, par exemple :
\[ \begin{cases} 3(x - 2) = y + 1 \\ 2(x - 2) = -z \end{cases} \iff \begin{cases} 3x - y - 7 = 0 \\ 2x + z - 4 = 0 \end{cases}. \]
Chacune de ces équations est celle d'un plan (nous y reviendrons en section 5) ; la droite $(d)$ est l'intersection de ces deux plans.
\end{exemplebox}

\begin{popfigure}
\begin{tikzpicture}[scale=1.0,>=Stealth]
  % Plan P1
  \fill[popblue!18] (-0.5,-0.4) -- (3.5,-1.4) -- (4.5,1.6) -- (0.5,2.6) -- cycle;
  \draw[popblue] (-0.5,-0.4) -- (3.5,-1.4) -- (4.5,1.6) -- (0.5,2.6) -- cycle;
  \node[popblue] at (3.9,2.0) {$(P_1)$};
  % Plan P2
  \fill[poporange!18] (0.2,-1.2) -- (4.2,-0.2) -- (3.2,2.8) -- (-0.8,1.8) -- cycle;
  \draw[poporange] (0.2,-1.2) -- (4.2,-0.2) -- (3.2,2.8) -- (-0.8,1.8) -- cycle;
  \node[poporange] at (-0.6,2.15) {$(P_2)$};
  % Droite d'intersection
  \draw[poppurple,very thick] (0.6,-0.9) -- (3.4,2.3);
  \node[poppurple] at (3.55,2.5) {$(d)$};
\end{tikzpicture}
\legende{Une droite $(d)$ de l'espace peut être vue comme l'intersection de deux plans sécants $(P_1)$ et $(P_2)$ : c'est le sens géométrique d'un système de deux équations cartésiennes.}
\end{popfigure}

\begin{attention}[Une droite n'a pas UNE équation cartésienne, mais UN SYSTÈME]
Dans le plan, une droite a une équation cartésienne unique (à un facteur près) : $ax + by + c = 0$. Dans l'espace, une seule équation $ax + by + cz + d = 0$ définit un \emph{plan}, pas une droite ! Une droite de l'espace est toujours donnée par un \textbf{système de deux équations} de plans (et il y a une infinité de choix possibles pour ces deux plans sécants). Écrire « l'équation de la droite est $3x - y - 7 = 0$ » est une erreur grave au Probatoire.
\end{attention}

\courssub{Cas particuliers : droites remarquables}

Certains coefficients nuls dans le vecteur directeur signalent des droites particulières, qu'il faut savoir reconnaître immédiatement.

\begin{aretenir}[Droites parallèles aux axes ou aux plans de coordonnées]
Soit $(d)$ de vecteur directeur $\vec{u}(a\,;b\,;c)$.
\begin{itemize}
\item Si $b = c = 0$ (donc $\vec{u}$ colinéaire à $\vec{\imath}$) : $(d)$ est \textbf{parallèle à l'axe des abscisses} $(Ox)$. Sa représentation s'écrit $\begin{cases} x = x_A + at \\ y = y_A \\ z = z_A \end{cases}$ : les coordonnées $y$ et $z$ sont constantes.
\item Si $a = c = 0$ : $(d)$ est parallèle à $(Oy)$ ; si $a = b = 0$ : $(d)$ est parallèle à $(Oz)$.
\item Si un seul coefficient est nul, par exemple $c = 0$ : $(d)$ est \textbf{parallèle au plan} $(xOy)$ (la cote $z$ reste constante : $z = z_A$).
\end{itemize}
\end{aretenir}

\begin{exemplebox}[Reconnaître des droites remarquables]
\begin{itemize}
\item $\begin{cases} x = 3 + 2t \\ y = 1 \\ z = -4 \end{cases}$ : droite parallèle à $(Ox)$, passant par $(3\,;1\,;-4)$.
\item $\begin{cases} x = 5 \\ y = 2 - t \\ z = 5 \end{cases}$ : droite parallèle à $(Oy)$.
\item $\begin{cases} x = 1 + t \\ y = 3 + 2t \\ z = 7 \end{cases}$ : droite parallèle au plan $(xOy)$, contenue dans le plan « horizontal » d'équation $z = 7$.
\end{itemize}
\end{exemplebox}

\begin{exercice}[Pour t'entraîner]
\begin{enumerate}
\item Détermine une représentation paramétrique de la droite $(\Delta)$ passant par $E(0\,;2\,;-3)$ et dirigée par $\vec{v}(4\,;-1\,;2)$.
\item La droite $(\Delta)$ coupe-t-elle le plan $(xOy)$ ? Si oui, en quel point ? \emph{Indication : un point de $(xOy)$ a une cote nulle.}
\item Donne un système d'équations cartésiennes de $(\Delta)$.
\item Les points $F(8\,;0\,;1)$ et $G(-4\,;3\,;-5)$ appartiennent-ils à $(\Delta)$ ?
\end{enumerate}
\end{exercice}

\begin{corrige}[Exercice]
\begin{enumerate}
\item $\begin{cases} x = 4t \\ y = 2 - t \\ z = -3 + 2t \end{cases}$, $t \in \mathbb{R}$.
\item On cherche $t$ tel que $z = 0$ : $-3 + 2t = 0$ donne $t = \dfrac{3}{2}$. Alors $x = 4 \times \dfrac{3}{2} = 6$ et $y = 2 - \dfrac{3}{2} = \dfrac{1}{2}$. La droite coupe $(xOy)$ au point $\left(6\,;\dfrac{1}{2}\,;0\right)$.
\item $t = \dfrac{x}{4} = 2 - y = \dfrac{z + 3}{2}$, soit par exemple $\begin{cases} x + 4y - 8 = 0 \\ x - 2z - 6 = 0 \end{cases}$.
\item Pour $F(8\,;0\,;1)$ : la première équation donne $t = 2$ ; alors $y = 2 - 2 = 0$ \checkmark\ et $z = -3 + 4 = 1$ \checkmark. Donc $F \in (\Delta)$. Pour $G(-4\,;3\,;-5)$ : la première équation donne $t = -1$ ; alors $y = 2 - (-1) = 3$ \checkmark\ et $z = -3 - 2 = -5$ \checkmark. Les trois équations sont satisfaites pour $t = -1$ : $G \in (\Delta)$.
\end{enumerate}
\end{corrige}

\begin{savaistu}[Des lignes droites dans le Grand Mbam]
Les ingénieurs qui tracent les lignes à haute tension entre la centrale de Song Loulou et Douala, ou qui positionnent les antennes relais de MTN et Orange le long de la Nationale, raisonnent exactement comme dans cette leçon : une ligne est définie par un point de départ et une direction, et chaque poteau correspond à une valeur du paramètre $t$. Les logiciels de cartographie et de conception (AutoCAD, systèmes GPS) stockent les droites de l'espace sous forme paramétrique : c'est la représentation la plus naturelle pour « avancer le long d'une ligne ».
\end{savaistu}

Tu sais désormais décrire une droite de l'espace de deux façons complémentaires : par une représentation paramétrique (un point, un vecteur directeur, un paramètre) et par un système d'équations cartésiennes (intersection de deux plans). Dans la prochaine section, nous utiliserons ces outils pour comparer deux droites entre elles : sécantes, parallèles ou non coplanaires — une situation sans équivalent dans le plan, et l'une des grandes nouveautés de la géométrie de l'espace.

\coursec{Positions relatives de deux droites}

Dans la section précédente, tu as appris à décrire une droite de l'espace à l'aide d'une représentation paramétrique : un point de passage et un vecteur directeur suffisent à la caractériser entièrement. Une question naturelle se pose alors : si l'on se donne \emph{deux} droites de l'espace, comment décrire la manière dont elles sont disposées l'une par rapport à l'autre ?

En géométrie plane, la réponse est bien connue : deux droites sont soit sécantes, soit parallèles (strictement parallèles ou confondues). Dans l'espace, une quatrième possibilité apparaît, qui déroute souvent les élèves : deux droites peuvent n'avoir \emph{aucun} point commun sans être parallèles. C'est toute la richesse de la géométrie dans l'espace, et c'est l'objet de cette section.

\courssub{Les quatre configurations possibles}

\begin{definition}[Positions relatives de deux droites de l'espace]
Soient $(d)$ et $(d')$ deux droites de l'espace. Une et une seule des quatre situations suivantes se produit :
\begin{itemize}
\item $(d)$ et $(d')$ sont \cle{confondues} : elles ont tous leurs points communs ;
\item $(d)$ et $(d')$ sont \cle{strictement parallèles} : elles sont coplanaires et n'ont aucun point commun ;
\item $(d)$ et $(d')$ sont \cle{sécantes} : elles ont exactement un point commun ;
\item $(d)$ et $(d')$ sont \cle{non coplanaires} : il n'existe aucun plan qui les contienne toutes les deux ; elles n'ont alors aucun point commun et ne sont pas parallèles.
\end{itemize}
Les trois premiers cas correspondent à des droites \cle{coplanaires} (contenues dans un même plan).
\end{definition}

\begin{popfigure}
\begin{tikzpicture}[scale=0.9]
% Panneau 1 : parallèles strictement
\begin{scope}[shift={(0,0)}]
\draw[popblue, thick] (-1.3,0) -- (1.3,0);
\draw[poporange, thick] (-1.3,0.8) -- (1.3,0.8);
\node at (0,-0.7) {\small parallèles};
\end{scope}
% Panneau 2 : confondues
\begin{scope}[shift={(4.2,0)}]
\draw[popblue, thick] (-1.3,0.4) -- (1.3,0.4);
\draw[poporange, thick, dashed] (-1.1,0.4) -- (1.1,0.4);
\node at (0,-0.7) {\small confondues};
\end{scope}
% Panneau 3 : sécantes
\begin{scope}[shift={(8.4,0)}]
\draw[popblue, thick] (-1.2,-0.3) -- (1.2,0.9);
\draw[poporange, thick] (-1.2,0.9) -- (1.2,-0.3);
\fill[popdark] (0,0.3) circle (1.5pt);
\node at (0,-0.7) {\small sécantes};
\end{scope}
% Panneau 4 : non coplanaires (cube)
\begin{scope}[shift={(12.6,0.2)}, scale=0.85]
\coordinate (A) at (0,0); \coordinate (B) at (1.4,0);
\coordinate (C) at (1.4,1.4); \coordinate (D) at (0,1.4);
\coordinate (E) at (0.5,0.5); \coordinate (F) at (1.9,0.5);
\coordinate (G) at (1.9,1.9); \coordinate (H) at (0.5,1.9);
\draw[popmuted] (A)--(B)--(C)--(D)--cycle;
\draw[popmuted] (B)--(F)--(G)--(C);
\draw[popmuted, dashed] (A)--(E)--(H)--(D);
\draw[popmuted, dashed] (E)--(F);
\draw[popmuted] (G)--(H);
\draw[popblue, very thick] (A)--(B);
\draw[poporange, very thick] (H)--(G);
\node at (0.95,-0.7) {\small non coplanaires};
\end{scope}
\end{tikzpicture}
\legende{Les quatre positions relatives possibles de deux droites de l'espace. Pour les droites non coplanaires, on a représenté deux arêtes d'un cube : elles ne se coupent pas et ne sont pas parallèles.}
\end{popfigure}

\begin{savaistu}[Des câbles qui ne se croisent pas... sans être parallèles]
Dans un immeuble en construction à Douala, observe deux câbles électriques : l'un court le long du plafond du couloir du rez-de-chaussée, l'autre traverse en diagonale le mur de l'escalier. Ils ne se croisent nulle part... et pourtant ils ne sont pas parallèles ! C'est la situation typique des droites \emph{non coplanaires}, totalement impossible en géométrie plane. C'est pourquoi les techniciens qui posent les installations doivent raisonner en trois dimensions, et pas seulement sur le plan du bâtiment.
\end{savaistu}

\courssub{Critère vectoriel de parallélisme}

L'outil essentiel de cette leçon est le \cle{vecteur directeur}. Rappelons que si $(d)$ passe par $A$ et est dirigée par $\vec{u}$, alors $M \in (d) \iff \overrightarrow{AM}$ et $\vec{u}$ sont colinéaires.

\begin{propriete}[Caractérisation du parallélisme]
Soient $(d)$ et $(d')$ deux droites de vecteurs directeurs respectifs $\vec{u}$ et $\vec{u'}$.
\begin{itemize}
\item $(d)$ et $(d')$ sont parallèles (strictement ou confondues) si et seulement si $\vec{u}$ et $\vec{u'}$ sont \cle{colinéaires}.
\item Si de plus un point de $(d)$ appartient à $(d')$, alors les droites sont \cle{confondues} ; sinon elles sont \cle{strictement parallèles}.
\end{itemize}
\end{propriete}

\begin{experience}[Démonstration guidée]
On raisonne par double implication.
\begin{enumerate}
\item \textbf{Sens direct.} Supposons $(d)$ et $(d')$ parallèles ou confondues. Par définition, elles sont coplanaires et ont la même direction : tout vecteur directeur de l'une dirige donc l'autre, ce qui signifie exactement que $\vec{u}$ et $\vec{u'}$ sont colinéaires.
\item \textbf{Sens réciproque.} Supposons $\vec{u'} = k\,\vec{u}$ avec $k \in \mathbb{R}^*$. Soit $A \in (d)$ et $A' \in (d')$.
\begin{itemize}
\item Si $A' \in (d)$, alors $(d')$ est la droite passant par $A'$ dirigée par $\vec{u}$ : c'est exactement $(d)$, donc les droites sont confondues.
\item Si $A' \notin (d)$, les deux droites sont contenues dans le plan défini par $A$, $\vec{u}$ et $\overrightarrow{AA'}$ ; dans ce plan, deux droites de même direction sont parallèles, et elles sont distinctes puisque $A' \notin (d)$ : elles sont strictement parallèles.
\end{itemize}
\end{enumerate}
La distinction entre « strictement parallèles » et « confondues » se ramène donc à un simple \textbf{test d'appartenance d'un point}.
\end{experience}

\begin{exemplebox}[Parallèles ou confondues ?]
Dans un repère $(O\,;\vec{\imath},\vec{\jmath},\vec{k})$, on considère :
\[ (d) : \begin{cases} x = 1 + 2t \\ y = 3 - t \\ z = -2 + 4t \end{cases} \quad \text{et} \quad (d') : \begin{cases} x = 5 - 4s \\ y = 1 + 2s \\ z = 6 - 8s \end{cases} \]
$(d)$ est dirigée par $\vec{u}(2\,;-1\,;4)$ et $(d')$ par $\vec{u'}(-4\,;2\,;-8)$. On remarque que $\vec{u'} = -2\,\vec{u}$ : les vecteurs sont colinéaires, donc les droites sont parallèles ou confondues.

Le point $A(1\,;3\,;-2)$ appartient à $(d)$. Appartient-il à $(d')$ ? On cherche $s$ tel que $5 - 4s = 1$, soit $s = 1$ ; alors $y = 1 + 2 = 3$ et $z = 6 - 8 = -2$. Les trois coordonnées conviennent : $A \in (d')$. Les droites sont donc \textbf{confondues}.
\end{exemplebox}

\courssub{Droites sécantes et droites non coplanaires}

Lorsque les vecteurs directeurs ne sont \emph{pas} colinéaires, les droites ne peuvent être ni parallèles ni confondues. Restent deux possibilités : sécantes ou non coplanaires. Pour trancher, on cherche un éventuel point d'intersection en résolvant un système.

\begin{propriete}[Critère d'intersection]
Soient $(d)$ et $(d')$ de vecteurs directeurs $\vec{u}$ et $\vec{u'}$ \textbf{non colinéaires}. On écrit leurs représentations paramétriques avec deux paramètres \emph{distincts} $t$ et $t'$, et l'on résout le système obtenu en égalant les coordonnées.
\begin{itemize}
\item Si le système admet un couple solution $(t\,;t')$, les droites sont \cle{sécantes} en un point dont on calcule les coordonnées.
\item Si le système n'admet aucune solution, les droites sont \cle{non coplanaires}.
\end{itemize}
\end{propriete}

\begin{experience}[Pourquoi ce critère fonctionne-t-il ?]
Un point $M$ appartient aux deux droites si et seulement s'il existe un réel $t$ (paramètre sur $(d)$) et un réel $t'$ (paramètre sur $(d')$) donnant les mêmes coordonnées. Égaler les trois coordonnées fournit un système de trois équations à deux inconnues $t$ et $t'$.
\begin{itemize}
\item S'il existe une solution, le point correspondant est commun aux deux droites. Il est \emph{unique} : sinon les droites, ayant deux points communs, seraient confondues, ce qui contredit la non-colinéarité de $\vec{u}$ et $\vec{u'}$.
\item S'il n'y a pas de solution, les droites sont disjointes et non parallèles. Or deux droites coplanaires disjointes sont nécessairement parallèles (résultat de géométrie plane) : les droites ne peuvent donc pas être coplanaires.
\end{itemize}
\end{experience}

\begin{attention}[Deux paramètres, deux lettres différentes !]
Le système d'intersection comporte \textbf{deux paramètres différents} : $t$ pour la première droite, $t'$ (ou $s$) pour la seconde. Ne jamais utiliser la même lettre ! Écrire le système avec un seul paramètre reviendrait à chercher des points atteints « au même instant » sur les deux droites, ce qui n'a aucun sens géométrique et conduit presque toujours à une fausse conclusion de non-intersection.
\end{attention}

\begin{attention}[Non colinéaires ne veut pas dire sécantes !]
C'est l'erreur la plus fréquente au Probatoire. Dans le plan, des vecteurs directeurs non colinéaires garantissent l'intersection. \textbf{Dans l'espace, c'est faux} : des droites de directions différentes peuvent être non coplanaires. Après avoir constaté que $\vec{u}$ et $\vec{u'}$ ne sont pas colinéaires, tu dois \emph{obligatoirement} résoudre le système d'intersection avant de conclure.
\end{attention}

\courssub{Stratégie générale d'étude}

\begin{methode}[Étudier la position relative de deux droites]
On se donne $(d)$ et $(d')$ par leurs représentations paramétriques.
\begin{enumerate}
\item \textbf{Lire} un vecteur directeur de chaque droite : $\vec{u}$ et $\vec{u'}$.
\item \textbf{Tester la colinéarité} : les coordonnées de $\vec{u}$ et $\vec{u'}$ sont-elles proportionnelles ?
\begin{itemize}
\item Si \textbf{oui} : tester si un point de $(d)$ appartient à $(d')$. Si oui, les droites sont \emph{confondues} ; si non, elles sont \emph{strictement parallèles}.
\item Si \textbf{non} : passer à l'étape 3.
\end{itemize}
\item \textbf{Résoudre le système d'intersection} à deux paramètres $t$ et $t'$ : on détermine $t$ et $t'$ à l'aide de deux équations, puis on \emph{vérifie} dans la troisième.
\begin{itemize}
\item Si la vérification fonctionne : droites \emph{sécantes}, et on calcule le point d'intersection.
\item Sinon : droites \emph{non coplanaires}.
\end{itemize}
\end{enumerate}
\end{methode}

\begin{popfigure}
\begin{tikzpicture}[
node distance=0.55cm,
startstop/.style={rectangle, rounded corners, draw=popdark, fill=popblueL, text width=4.6cm, align=center, font=\small},
decision/.style={diamond, draw=popdark, fill=popgoldL, text width=3.4cm, align=center, font=\small, aspect=2.2},
result/.style={rectangle, draw=popdark, fill=popgreenL, text width=3.2cm, align=center, font=\small},
arrow/.style={-{Stealth}, thick, popdark}]
\node[startstop] (start) {Lire $\vec{u}$ et $\vec{u'}$, vecteurs directeurs de $(d)$ et $(d')$};
\node[decision, below=of start] (col) {$\vec{u}$ et $\vec{u'}$ colinéaires ?};
\node[decision, below left=1.1cm and -0.6cm of col] (point) {Un point de $(d)$ est-il sur $(d')$ ?};
\node[decision, below right=1.1cm and -0.9cm of col] (sys) {Le système d'intersection a-t-il une solution ?};
\node[result, below left=0.9cm and -0.5cm of point] (conf) {Droites confondues};
\node[result, below right=0.9cm and -0.5cm of point] (par) {Strictement parallèles};
\node[result, below left=0.9cm and -0.6cm of sys] (sec) {Sécantes en un point};
\node[result, below right=0.9cm and -0.6cm of sys] (nc) {Non coplanaires};
\draw[arrow] (start) -- (col);
\draw[arrow] (col) -| node[above left, font=\small]{oui} (point);
\draw[arrow] (col) -| node[above right, font=\small]{non} (sys);
\draw[arrow] (point) -| node[above left, font=\small]{oui} (conf);
\draw[arrow] (point) -| node[above right, font=\small]{non} (par);
\draw[arrow] (sys) -| node[above left, font=\small]{oui} (sec);
\draw[arrow] (sys) -| node[above right, font=\small]{non} (nc);
\end{tikzpicture}
\legende{Arbre de décision : méthode complète pour déterminer la position relative de deux droites de l'espace.}
\end{popfigure}

\begin{exoresolu}[Deux études complètes : cas non coplanaires]
Dans un repère $(O\,;\vec{\imath},\vec{\jmath},\vec{k})$ de l'espace, on considère les droites :
\[ (d_1) : \begin{cases} x = 2 + t \\ y = -1 + 2t \\ z = 3 - t \end{cases} \quad (d_2) : \begin{cases} x = 1 + 2t' \\ y = 5 - t' \\ z = 1 + t' \end{cases} \quad (d_3) : \begin{cases} x = 3 + s \\ y = 2s \\ z = -1 + s \end{cases} \]
\textbf{1.} Étudier la position relative de $(d_1)$ et $(d_2)$.

\textbf{2.} Étudier la position relative de $(d_1)$ et $(d_3)$.
\tcblower
\textbf{1. Position de $(d_1)$ et $(d_2)$.}

\emph{Étape 1 : vecteurs directeurs.} $(d_1)$ est dirigée par $\vec{u_1}(1\,;2\,;-1)$ et $(d_2)$ par $\vec{u_2}(2\,;-1\,;1)$. Ces vecteurs ne sont pas colinéaires : en effet, $\dfrac{2}{1} = 2$ alors que $\dfrac{-1}{2} = -\dfrac{1}{2}$ ; les coordonnées ne sont pas proportionnelles.

\emph{Étape 2 : système d'intersection.} On cherche $t$ et $t'$ tels que :
\[ \begin{cases} 2 + t = 1 + 2t' \\ -1 + 2t = 5 - t' \\ 3 - t = 1 + t' \end{cases} \iff \begin{cases} t - 2t' = -1 \\ 2t + t' = 6 \\ -t - t' = -2 \end{cases} \]
On utilise les deux premières équations. La première donne $t = 2t' - 1$ ; en reportant dans la deuxième :
\[ 2(2t' - 1) + t' = 6 \iff 5t' = 8 \iff t' = \dfrac{8}{5}, \quad \text{d'où} \quad t = \dfrac{16}{5} - 1 = \dfrac{11}{5}. \]
\emph{Vérification dans la troisième équation :} $-t - t' = -\dfrac{11}{5} - \dfrac{8}{5} = -\dfrac{19}{5} \neq -2$. La troisième équation n'est pas satisfaite : le système n'a \textbf{aucune solution}. Les droites $(d_1)$ et $(d_2)$ sont \textbf{non coplanaires}.

\textbf{2. Position de $(d_1)$ et $(d_3)$.}

\emph{Étape 1.} $(d_3)$ est dirigée par $\vec{u_3}(1\,;2\,;1)$. Or $\vec{u_1}(1\,;2\,;-1)$ et $\vec{u_3}$ ne sont pas colinéaires : $\dfrac{1}{1} = \dfrac{2}{2} = 1$ mais $\dfrac{1}{-1} = -1 \neq 1$.

\emph{Étape 2 : système.}
\[ \begin{cases} 2 + t = 3 + s \\ -1 + 2t = 2s \\ 3 - t = -1 + s \end{cases} \iff \begin{cases} t - s = 1 \\ 2t - 2s = 1 \\ -t - s = -4 \end{cases} \]
La première équation donne $t = s + 1$. En reportant dans la troisième : $-(s+1) - s = -4$, soit $-2s = -3$, d'où $s = \dfrac{3}{2}$ et $t = \dfrac{5}{2}$.

\emph{Vérification dans la deuxième équation :} $2t - 2s = 5 - 3 = 2 \neq 1$. Le système n'a pas de solution : $(d_1)$ et $(d_3)$ sont \textbf{non coplanaires}.

\emph{Remarque méthodologique.} Dans les deux cas, deux des trois équations fournissent bien un couple $(t\,;t')$, mais la troisième le rejette. C'est exactement ce qui distingue l'espace du plan : la compatibilité doit porter sur les \textbf{trois} équations simultanément. Rédige toujours la résolution étape par étape, et conclus seulement après la vérification finale.
\end{exoresolu}

\begin{exoresolu}[Un couple de droites sécantes]
On considère les droites :
\[ (d) : \begin{cases} x = 1 + t \\ y = 2 - t \\ z = 3 + 2t \end{cases} \quad \text{et} \quad (d') : \begin{cases} x = 4 - t' \\ y = -1 + t' \\ z = 1 + 3t' \end{cases} \]
Montrer que $(d)$ et $(d')$ sont sécantes et déterminer les coordonnées de leur point d'intersection $I$.
\tcblower
\emph{Étape 1.} $\vec{u}(1\,;-1\,;2)$ dirige $(d)$ et $\vec{u'}(-1\,;1\,;3)$ dirige $(d')$. Sont-ils colinéaires ? Il faudrait un réel $k$ tel que $-1 = k$, $1 = -k$ et $3 = 2k$ : les deux premières conditions donnent $k = -1$, mais alors $2k = -2 \neq 3$. Les vecteurs ne sont \textbf{pas colinéaires}.

\emph{Étape 2.} On résout :
\[ \begin{cases} 1 + t = 4 - t' \\ 2 - t = -1 + t' \\ 3 + 2t = 1 + 3t' \end{cases} \iff \begin{cases} t + t' = 3 \\ -t - t' = -3 \\ 2t - 3t' = -2 \end{cases} \]
Les deux premières équations sont équivalentes à $t + t' = 3$, soit $t' = 3 - t$. En reportant dans la troisième : $2t - 3(3 - t) = -2$, c'est-à-dire $5t - 9 = -2$, d'où $t = \dfrac{7}{5}$ et $t' = 3 - \dfrac{7}{5} = \dfrac{8}{5}$.

Le système admet une solution : les droites sont \textbf{sécantes}. Le point d'intersection s'obtient en reportant $t = \dfrac{7}{5}$ dans $(d)$ :
\[ I\left(1 + \dfrac{7}{5}\,;\ 2 - \dfrac{7}{5}\,;\ 3 + \dfrac{14}{5}\right) \quad \text{soit} \quad I\left(\dfrac{12}{5}\,;\ \dfrac{3}{5}\,;\ \dfrac{29}{5}\right). \]
\emph{Contrôle} avec $(d')$ et $t' = \dfrac{8}{5}$ : $x = 4 - \dfrac{8}{5} = \dfrac{12}{5}$, $y = -1 + \dfrac{8}{5} = \dfrac{3}{5}$, $z = 1 + \dfrac{24}{5} = \dfrac{29}{5}$. Les coordonnées coincident : le résultat est confirmé.
\end{exoresolu}

\courssub{Lien avec la coplanarité : deux droites définissent un plan}

\begin{aretenir}[Droites et plans]
Deux droites \cle{sécantes} ou deux droites \cle{parallèles} (strictement ou confondues) sont toujours coplanaires : elles définissent un \textbf{plan unique}.
\begin{itemize}
\item Si $(d)$ et $(d')$ sont sécantes en $I$, le plan qu'elles définissent est dirigé par le couple de vecteurs $(\vec{u}, \vec{u'})$, non colinéaires, et passe par $I$.
\item Si $(d)$ et $(d')$ sont strictement parallèles, avec $A \in (d)$ et $A' \in (d')$, le plan qu'elles définissent passe par $A$ et est dirigé par $\vec{u}$ et $\overrightarrow{AA'}$.
\end{itemize}
À l'inverse, deux droites \cle{non coplanaires} ne sont contenues dans aucun plan commun : c'est la seule configuration « purement spatiale ».
\end{aretenir}

Ce résultat sera précieux dans la suite de la leçon : déterminer une représentation paramétrique ou une équation cartésienne du plan défini par deux droites sécantes ou parallèles est un exercice classique du Probatoire, qui relie directement cette section à la suivante.

\begin{exercice}[À toi de jouer]
Dans un repère de l'espace, on donne :
\[ (d) : \begin{cases} x = 3 - t \\ y = 1 + 2t \\ z = -2 + t \end{cases} \qquad (d') : \begin{cases} x = -1 + 2t' \\ y = 5 - 4t' \\ z = -4 - 2t' \end{cases} \qquad (d'') : \begin{cases} x = 2 + s \\ y = -3 + s \\ z = 1 - s \end{cases} \]
Déterminer la position relative de $(d)$ et $(d')$, puis celle de $(d)$ et $(d'')$.
\end{exercice}

\begin{corrige}[Exercice 1]
\textbf{Position de $(d)$ et $(d')$.} $\vec{u}(-1\,;2\,;1)$ dirige $(d)$ et $\vec{u'}(2\,;-4\,;-2)$ dirige $(d')$. On a $\vec{u'} = -2\,\vec{u}$ : les vecteurs sont colinéaires, donc les droites sont parallèles ou confondues. Le point $A(3\,;1\,;-2) \in (d)$ appartient-il à $(d')$ ? Il faudrait $-1 + 2t' = 3$, soit $t' = 2$ ; alors $y = 5 - 8 = -3 \neq 1$. Donc $A \notin (d')$ : les droites sont \textbf{strictement parallèles}.

\textbf{Position de $(d)$ et $(d'')$.} $\vec{u''}(1\,;1\,;-1)$ n'est pas colinéaire à $\vec{u}(-1\,;2\,;1)$ : en effet $\dfrac{1}{-1} = -1$ alors que $\dfrac{1}{2} \neq -1$, les coordonnées ne sont pas proportionnelles. On résout le système d'intersection :
\[ \begin{cases} 3 - t = 2 + s \\ 1 + 2t = -3 + s \\ -2 + t = 1 - s \end{cases} \iff \begin{cases} t + s = 1 \\ 2t - s = -4 \\ t + s = 3 \end{cases} \]
La première équation impose $t + s = 1$ et la troisième $t + s = 3$ : contradiction. Le système n'a pas de solution : les droites $(d)$ et $(d'')$ sont \textbf{non coplanaires}.
\end{corrige}

\coursec{Représentations paramétriques d'un plan}

Dans la section précédente, nous avons appris à comparer deux droites de l'espace : sécantes, parallèles ou non coplanaires. Nous avons vu qu'une droite est entièrement décrite par \emph{un point} et \emph{un vecteur directeur}, et qu'un seul paramètre $t$ suffit pour parcourir tous ses points. Mais l'espace ne se résume pas à des fils : il contient aussi des \emph{surfaces planes} — le plateau d'une table au marché de Mokolo, la tôle d'un toit, la surface d'un terrain de football. Comment décrire algébriquement un tel objet ? C'est toute la question de cette section : il faudra désormais \emph{deux} directions et donc \emph{deux} paramètres.

\courssub{Caractérisation vectorielle d'un plan}

\textbf{Intuition.} Pose une feuille de papier sur ta table. Pour atteindre n'importe quel point de cette feuille depuis un point fixe $A$ de la feuille, tu peux te déplacer d'une certaine quantité dans une première direction (par exemple le long du bord gauche), puis d'une autre quantité dans une seconde direction (par exemple le long du bord du bas). Deux directions non parallèles suffisent : c'est exactement l'idée du plan vectoriel.

\begin{definition}[Plan défini par un point et deux vecteurs non colinéaires]
Soit $A$ un point de l'espace et $\vec{u}$, $\vec{v}$ deux vecteurs \cle{non colinéaires}. Le \cle{plan} passant par $A$ et dirigé par $\vec{u}$ et $\vec{v}$, noté $(A\,;\,\vec{u},\vec{v})$, est l'ensemble des points $M$ de l'espace tels que :
\[ \overrightarrow{AM} = t\,\vec{u} + s\,\vec{v}, \quad \text{avec } (t,s) \in \mathbb{R}^2. \]
Les vecteurs $\vec{u}$ et $\vec{v}$ sont appelés \cle{vecteurs directeurs} du plan ; le couple $(\vec{u},\vec{v})$ est un \cle{couple de vecteurs directeurs} du plan.
\end{definition}

\begin{attention}[La condition « non colinéaires » est vitale]
Si $\vec{u}$ et $\vec{v}$ sont \textbf{colinéaires} (par exemple $\vec{v} = 2\vec{u}$), alors $t\vec{u} + s\vec{v} = (t + 2s)\vec{u}$ ne décrit que des multiples de $\vec{u}$ : l'ensemble obtenu est une \textbf{droite}, pas un plan. Avant toute conclusion, il faut donc \emph{toujours} vérifier que les deux vecteurs directeurs ne sont pas proportionnels.
\end{attention}

\begin{popfigure}
\begin{tikzpicture}[scale=1.1, >=Stealth]
  % Plan dessiné comme parallélogramme
  \coordinate (O) at (0,0);
  \coordinate (U) at (4.2,0.6);
  \coordinate (V) at (1.4,2.2);
  \coordinate (UV) at ($(U)+(V)$);
  \fill[popblueL, opacity=0.55] (O) -- (U) -- (UV) -- (V) -- cycle;
  \draw[popblue, thick] (O) -- (U) -- (UV) -- (V) -- cycle;
  % Point A
  \coordinate (A) at (1.1,0.75);
  \fill[popdark] (A) circle (1.6pt) node[below left, popdark] {$A$};
  % Vecteurs u et v depuis A
  \coordinate (u) at ($(A)+(1.6,0.23)$);
  \coordinate (v) at ($(A)+(0.55,0.85)$);
  \draw[->, poporange, very thick] (A) -- (u) node[midway, below, poporange] {$\vec{u}$};
  \draw[->, popgreen, very thick] (A) -- (v) node[midway, left, popgreen] {$\vec{v}$};
  % Point M = A + 2.2u + 1.5v (coordonnées réduites)
  \coordinate (M) at ($(A)+(2.2,0.31)+(0.83,1.28)$);
  \draw[->, poppurple, thick, dashed] (A) -- ($(A)+(2.2,0.31)$) node[midway, below, poppurple] {\small $t\vec{u}$};
  \draw[->, poppurple, thick, dashed] ($(A)+(2.2,0.31)$) -- (M) node[midway, right, poppurple] {\small $s\vec{v}$};
  \fill[poppurple] (M) circle (1.6pt) node[above right, poppurple] {$M$};
  \node[popblue] at (4.6,2.4) {$\mathcal{P}$};
\end{tikzpicture}
\legende{Le plan $\mathcal{P} = (A\,;\,\vec{u},\vec{v})$ : tout point $M$ du plan s'écrit $\overrightarrow{AM} = t\vec{u} + s\vec{v}$.}
\end{popfigure}

\begin{propriete}[Équivalence des trois modes de définition d'un plan]
Les trois façons suivantes de définir un plan sont équivalentes :
\begin{enumerate}
  \item un point $A$ et deux vecteurs $\vec{u}$, $\vec{v}$ non colinéaires ;
  \item trois points $A$, $B$, $C$ \textbf{non alignés} : le plan est alors $(A\,;\,\overrightarrow{AB},\overrightarrow{AC})$ ;
  \item deux points distincts $A$, $B$ et un vecteur $\vec{w}$ non colinéaire à $\overrightarrow{AB}$ : le plan est $(A\,;\,\overrightarrow{AB},\vec{w})$.
\end{enumerate}
\end{propriete}

\begin{experience}[Démonstration de l'équivalence]
Montrons que (2) équivaut à (1), le cas (3) se traitant de la même manière.
\begin{itemize}
  \item \textbf{De (2) vers (1).} Dire que $A$, $B$, $C$ ne sont pas alignés signifie exactement que $\overrightarrow{AB}$ et $\overrightarrow{AC}$ ne sont pas colinéaires (sinon $C$ serait sur la droite $(AB)$). Le couple $(\overrightarrow{AB},\overrightarrow{AC})$ est donc bien un couple de vecteurs directeurs, et le plan $(A\,;\,\overrightarrow{AB},\overrightarrow{AC})$ est bien défini.
  \item \textbf{De (1) vers (2).} Partons du plan $(A\,;\,\vec{u},\vec{v})$ avec $\vec{u}$ et $\vec{v}$ non colinéaires. Posons $B$ tel que $\overrightarrow{AB} = \vec{u}$ et $C$ tel que $\overrightarrow{AC} = \vec{v}$. Alors $B$ et $C$ appartiennent au plan (prendre $(t,s) = (1,0)$ puis $(0,1)$), et $A$, $B$, $C$ ne sont pas alignés puisque $\overrightarrow{AB} = \vec{u}$ et $\overrightarrow{AC} = \vec{v}$ ne sont pas colinéaires. Le plan est donc aussi « le plan $(ABC)$ ». \hfill $\square$
\end{itemize}
Retiens : \textbf{trois points non alignés définissent un unique plan}. C'est pourquoi une table à trois pieds ne bancale jamais, alors qu'une table à quatre pieds peut bancaler !
\end{experience}

\begin{popfigure}
\begin{tikzpicture}[scale=1.05, >=Stealth]
  \coordinate (O) at (0,0);
  \coordinate (U) at (4.4,0.7);
  \coordinate (V) at (1.5,2.3);
  \coordinate (UV) at ($(U)+(V)$);
  \fill[popgreenL, opacity=0.5] (O) -- (U) -- (UV) -- (V) -- cycle;
  \draw[popgreen, thick] (O) -- (U) -- (UV) -- (V) -- cycle;
  \coordinate (A) at (0.9,0.6);
  \coordinate (B) at (3.3,1.05);
  \coordinate (C) at (1.9,2.15);
  \draw[->, poporange, very thick] (A) -- (B) node[midway, below, poporange] {$\overrightarrow{AB}$};
  \draw[->, poppurple, very thick] (A) -- (C) node[midway, above left, poppurple] {$\overrightarrow{AC}$};
  \draw[popmuted, dashed] (B) -- (C);
  \fill[popdark] (A) circle (1.6pt) node[below left, popdark] {$A$};
  \fill[popdark] (B) circle (1.6pt) node[below right, popdark] {$B$};
  \fill[popdark] (C) circle (1.6pt) node[above, popdark] {$C$};
  \node[popgreen] at (4.9,2.5) {$(ABC)$};
\end{tikzpicture}
\legende{Trois points non alignés $A$, $B$, $C$ définissent le plan $(A\,;\,\overrightarrow{AB},\overrightarrow{AC})$.}
\end{popfigure}

\courssub{Représentation paramétrique d'un plan}

Plaçons-nous dans un repère $(O\,;\,\vec{\imath},\vec{\jmath},\vec{k})$ de l'espace. Traduisons la relation $\overrightarrow{AM} = t\vec{u} + s\vec{v}$ coordonnée par coordonnée.

\begin{definition}[Représentation paramétrique d'un plan]
Soit $\mathcal{P}$ le plan passant par $A(x_A\,;\,y_A\,;\,z_A)$ et dirigé par les vecteurs non colinéaires $\vec{u}(a\,;\,b\,;\,c)$ et $\vec{v}(a'\,;\,b'\,;\,c')$. Une \cle{représentation paramétrique} de $\mathcal{P}$ est :
\[ \begin{cases} x = x_A + at + a's \\ y = y_A + bt + b's \\ z = z_A + ct + c's \end{cases} \quad (t,s) \in \mathbb{R}^2. \]
Les réels $t$ et $s$ sont les \cle{paramètres}. À chaque couple $(t,s)$ correspond un unique point du plan, et réciproquement.
\end{definition}

\begin{aretenir}[Droite ou plan : combien de paramètres ?]
\begin{itemize}
  \item Une \textbf{droite} : \textbf{un} point, \textbf{un} vecteur directeur, \textbf{un} paramètre $t$.
  \item Un \textbf{plan} : \textbf{un} point, \textbf{deux} vecteurs directeurs non colinéaires, \textbf{deux} paramètres $t$ et $s$.
\end{itemize}
Le nombre de paramètres indique la « dimension » de l'objet : on dit qu'une droite est de dimension $1$ et un plan de dimension $2$.
\end{aretenir}

\begin{exemplebox}[Une représentation paramétrique de plan]
Soit le plan $\mathcal{P}$ passant par $A(1\,;\,0\,;\,2)$ et dirigé par $\vec{u}(1\,;\,1\,;\,0)$ et $\vec{v}(0\,;\,2\,;\,1)$.
\begin{itemize}
  \item \textbf{Non-colinéarité :} $\vec{u}$ et $\vec{v}$ sont-ils proportionnels ? Il faudrait un réel $k$ avec $0 = k \times 1$ (premières coordonnées), donc $k = 0$, mais alors $2 = k \times 1 = 0$ : absurde. Donc $\vec{u}$ et $\vec{v}$ sont \textbf{non colinéaires} : $\mathcal{P}$ est bien un plan.
  \item \textbf{Représentation paramétrique :}
  \[ \mathcal{P} : \begin{cases} x = 1 + t \\ y = t + 2s \\ z = 2 + s \end{cases} \quad (t,s) \in \mathbb{R}^2. \]
  \item Par exemple, pour $(t,s) = (1,1)$ on obtient le point $(2\,;\,3\,;\,3)$ ; pour $(t,s) = (0,0)$ on retrouve $A$.
\end{itemize}
\end{exemplebox}

\courssub{Reconnaître une représentation paramétrique de plan}

Face à un système à deux paramètres, deux réflexes : extraire les éléments géométriques, puis vérifier que l'on a bien affaire à un plan.

\begin{methode}[Reconnaître un plan donné par une représentation paramétrique]
On considère un système de la forme $x = \alpha + at + a's$, $y = \beta + bt + b's$, $z = \gamma + ct + c's$.
\begin{enumerate}
  \item \textbf{Lire un point} : les termes constants donnent $A(\alpha\,;\,\beta\,;\,\gamma)$ (c'est le point obtenu pour $t = s = 0$).
  \item \textbf{Lire les vecteurs directeurs} : les coefficients de $t$ donnent $\vec{u}(a\,;\,b\,;\,c)$, ceux de $s$ donnent $\vec{v}(a'\,;\,b'\,;\,c')$.
  \item \textbf{Tester la colinéarité} : chercher un réel $k$ tel que $\vec{v} = k\vec{u}$ (comparer les coordonnées deux à deux). Si un tel $k$ n'existe pas, les vecteurs sont non colinéaires et le système définit bien un \textbf{plan}.
  \item Si $\vec{u}$ et $\vec{v}$ sont colinéaires (et non tous deux nuls), le système définit une \textbf{droite} : on peut alors éliminer l'un des deux paramètres.
\end{enumerate}
\end{methode}

\begin{exemplebox}[Reconnaissance]
Considérons le système $\begin{cases} x = 3 - t + 2s \\ y = 1 + 2t \\ z = -2 + t - s \end{cases}$, $(t,s) \in \mathbb{R}^2$.
\begin{itemize}
  \item Point : $A(3\,;\,1\,;\,-2)$.
  \item Vecteurs : $\vec{u}(-1\,;\,2\,;\,1)$ (coefficients de $t$) et $\vec{v}(2\,;\,0\,;\,-1)$ (coefficients de $s$).
  \item Colinéarité : $\vec{v} = k\vec{u}$ exigerait $2 = -k$ et $0 = 2k$, soit $k = -2$ et $k = 0$ simultanément : impossible. Les vecteurs sont non colinéaires, le système définit bien un \textbf{plan}.
\end{itemize}
\end{exemplebox}

\begin{attention}[Le piège classique : deux paramètres ne font pas toujours un plan]
Considère le système $\begin{cases} x = 1 + t + 2s \\ y = 2 - t - 2s \\ z = 3t + 6s \end{cases}$. Ici $\vec{u}(1\,;\,-1\,;\,3)$ et $\vec{v}(2\,;\,-2\,;\,6) = 2\vec{u}$ : les vecteurs sont \textbf{colinéaires} ! En posant $\lambda = t + 2s$, le système se réécrit $x = 1 + \lambda$, $y = 2 - \lambda$, $z = 3\lambda$ : c'est la représentation paramétrique d'une \textbf{droite} passant par $(1\,;\,2\,;\,0)$ et dirigée par $\vec{u}$. Ne conclus jamais « c'est un plan » sous prétexte qu'il y a deux lettres $t$ et $s$ : vérifie d'abord la non-colinéarité.
\end{attention}

\courssub{Appartenance d'un point à un plan}

\begin{methode}[Tester si un point appartient à un plan donné paramétriquement]
Pour savoir si $M(x_M\,;\,y_M\,;\,z_M)$ appartient au plan $\mathcal{P} : \begin{cases} x = x_A + at + a's \\ y = y_A + bt + b's \\ z = z_A + ct + c's \end{cases}$ :
\begin{enumerate}
  \item Remplacer $x$, $y$, $z$ par les coordonnées de $M$ : on obtient un système de \textbf{trois équations} à \textbf{deux inconnues} $t$ et $s$.
  \item Résoudre avec deux des équations pour trouver $t$ et $s$.
  \item \textbf{Vérifier} que le couple $(t,s)$ trouvé satisfait la troisième équation. Si oui, $M \in \mathcal{P}$ ; sinon, $M \notin \mathcal{P}$.
\end{enumerate}
\end{methode}

\begin{exoresolu}[Plan par un point et deux vecteurs ; test d'appartenance]
Soit $\mathcal{P}$ le plan passant par $A(1\,;\,0\,;\,2)$, dirigé par $\vec{u}(1\,;\,1\,;\,0)$ et $\vec{v}(0\,;\,2\,;\,1)$.
\begin{enumerate}
  \item Donner une représentation paramétrique de $\mathcal{P}$.
  \item Le point $B(2\,;\,3\,;\,3)$ appartient-il à $\mathcal{P}$ ?
  \item Le point $C(0\,;\,1\,;\,0)$ appartient-il à $\mathcal{P}$ ?
\end{enumerate}
\tcblower
\textbf{1. Représentation paramétrique.} On a vérifié plus haut que $\vec{u}$ et $\vec{v}$ ne sont pas colinéaires. Donc :
\[ \mathcal{P} : \begin{cases} x = 1 + t \\ y = t + 2s \\ z = 2 + s \end{cases} \quad (t,s) \in \mathbb{R}^2. \]

\textbf{2. Test pour $B(2\,;\,3\,;\,3)$.} On remplace :
\[ \begin{cases} 2 = 1 + t \\ 3 = t + 2s \\ 3 = 2 + s \end{cases} \]
La première équation donne $t = 1$ ; la troisième donne $s = 1$. Vérifions dans la deuxième : $t + 2s = 1 + 2 = 3$ : c'est bien égal à $y_B = 3$. Le couple $(t,s) = (1,1)$ convient pour les trois équations, donc $B \in \mathcal{P}$.

\textbf{3. Test pour $C(0\,;\,1\,;\,0)$.} On remplace :
\[ \begin{cases} 0 = 1 + t \\ 1 = t + 2s \\ 0 = 2 + s \end{cases} \]
La première équation donne $t = -1$ ; la troisième donne $s = -2$. Vérifions dans la deuxième : $t + 2s = -1 - 4 = -5 \neq 1$. Le système est incompatible, donc $C \notin \mathcal{P}$.
\end{exoresolu}

\begin{exoresolu}[Plan défini par trois points non alignés]
Soient $A(1\,;\,1\,;\,0)$, $B(2\,;\,0\,;\,1)$ et $C(0\,;\,3\,;\,1)$.
\begin{enumerate}
  \item Justifier que $A$, $B$, $C$ définissent un plan.
  \item Donner une représentation paramétrique du plan $(ABC)$.
\end{enumerate}
\tcblower
\textbf{1.} On calcule $\overrightarrow{AB}(1\,;\,-1\,;\,1)$ et $\overrightarrow{AC}(-1\,;\,2\,;\,1)$. Ces vecteurs sont-ils colinéaires ? Il faudrait $k$ avec $-1 = k \times 1$ donc $k = -1$, mais alors $2 = k \times (-1) = 1$ : absurde. Ils sont non colinéaires, donc $A$, $B$, $C$ sont non alignés et définissent un unique plan.

\textbf{2.} Le plan $(ABC)$ est le plan $(A\,;\,\overrightarrow{AB},\overrightarrow{AC})$ :
\[ (ABC) : \begin{cases} x = 1 + t - s \\ y = 1 - t + 2s \\ z = t + s \end{cases} \quad (t,s) \in \mathbb{R}^2. \]
Vérification mentale : $(t,s) = (1,0)$ redonne $B$, $(t,s) = (0,1)$ redonne $C$.
\end{exoresolu}

\begin{savaistu}[Pourquoi « paramétrique » ?]
L'adjectif vient du mot grec \emph{parametros} : les paramètres $t$ et $s$ jouent le rôle de « coordonnées internes » au plan, exactement comme la longitude et la latitude repèrent un point à la surface de la Terre. Un géographe de l'Institut National de Cartographie du Cameroun qui localise un village par $(\text{longitude}, \text{latitude})$ fait, sans le dire, de la géométrie paramétrique : deux nombres suffisent pour décrire toute une surface. C'est aussi le principe utilisé en infographie 3D pour modéliser les carrosseries de voitures ou les paysages des jeux vidéo.
\end{savaistu}

\begin{exercice}[À toi de jouer]
Soit le système $\begin{cases} x = 2 + t - s \\ y = -1 + 3t \\ z = 4 - t + 2s \end{cases}$, $(t,s) \in \mathbb{R}^2$.
\begin{enumerate}
  \item Montrer que ce système définit un plan $\mathcal{P}$ dont on précisera un point et deux vecteurs directeurs.
  \item Les points $D(1\,;\,2\,;\,7)$ et $E(3\,;\,2\,;\,3)$ appartiennent-ils à $\mathcal{P}$ ?
  \item Donner une représentation paramétrique du plan passant par $F(0\,;\,1\,;\;2)$ et $G(1\,;\,0\,;\,3)$ et dirigé par $\vec{w}(2\,;\,1\,;\,-1)$, après avoir justifié qu'il est bien défini.
\end{enumerate}
\end{exercice}

\begin{corrige}[Exercice 1]
\textbf{1.} Point : $A(2\,;\,-1\,;\,4)$. Vecteurs : $\vec{u}(1\,;\,3\,;\,-1)$ et $\vec{v}(-1\,;\,0\,;\,2)$. Colinéarité : $\vec{v} = k\vec{u}$ exigerait $-1 = k$ et $0 = 3k$, impossible. Donc $\vec{u}$ et $\vec{v}$ sont non colinéaires : le système définit bien un plan.

\textbf{2.} Pour $D(1\,;\,2\,;\,7)$ : $1 = 2 + t - s$, $2 = -1 + 3t$, $7 = 4 - t + 2s$. La deuxième équation donne $t = 1$ ; la première donne alors $s = 2$. Vérification : $4 - 1 + 4 = 7$ : correct. Donc $D \in \mathcal{P}$. Pour $E(3\,;\,2\,;\,3)$ : $t = 1$ (deuxième équation), puis $3 = 2 + 1 - s$ donne $s = 0$ ; vérification : $4 - 1 + 0 = 3$ : correct. Donc $E \in \mathcal{P}$.

\textbf{3.} $\overrightarrow{FG}(1\,;\,-1\,;\,1)$. Est-il colinéaire à $\vec{w}(2\,;\,1\,;\,-1)$ ? Il faudrait $2 = k$ et $1 = -k$, impossible. Le plan est bien défini :
\[ \begin{cases} x = t + 2s \\ y = 1 - t + s \\ z = 2 + t - s \end{cases} \quad (t,s) \in \mathbb{R}^2. \]
\end{corrige}

\begin{aretenir}[L'essentiel de la section]
\begin{itemize}
  \item Un plan est défini par \textbf{un point et deux vecteurs non colinéaires}, ou de façon équivalente par \textbf{trois points non alignés}, ou par \textbf{deux points distincts et un vecteur non colinéaire à} $\overrightarrow{AB}$.
  \item Sa représentation paramétrique fait intervenir \textbf{deux paramètres} : $x = x_A + at + a's$, $y = y_A + bt + b's$, $z = z_A + ct + c's$.
  \item Les termes constants donnent le point, les coefficients de $t$ et de $s$ donnent les vecteurs directeurs.
  \item \textbf{Toujours} vérifier la non-colinéarité des vecteurs directeurs, sinon l'objet décrit est une droite.
  \item Pour tester l'appartenance d'un point, on résout un système de trois équations à deux inconnues $t$ et $s$, avec vérification dans l'équation restante.
\end{itemize}
Dans la prochaine section, nous verrons qu'un plan peut aussi être décrit par une unique \emph{équation cartésienne} $ax + by + cz + d = 0$, grâce à la notion de vecteur normal.
\end{aretenir}

\coursec{Équation cartésienne d'un plan et vecteur normal}

Dans la section précédente, nous avons décrit un plan par une \cle{représentation paramétrique} : un point $A$ du plan, deux vecteurs directeurs $\vec{u}$ et $\vec{v}$ non colinéaires, et deux paramètres réels $t$ et $s$ tels que $\overrightarrow{AM} = t\vec{u} + s\vec{v}$. Cette description est très utile, mais elle a un défaut : pour savoir si un point $M$ donné appartient au plan, il faut résoudre un système. Ne pourrait-on pas caractériser un plan par une \emph{seule} équation reliant directement les coordonnées $(x\,;y\,;z)$ de ses points, comme l'équation $ax + by + c = 0$ caractérise une droite dans le plan ? C'est exactement l'objet de cette section, et la clé sera une idée géométrique simple : un plan est entièrement déterminé par un point et une \emph{direction perpendiculaire} au plan.

\courssub{Vecteur normal à un plan}

\textbf{Intuition.} Pose une feuille de papier à plat sur ta table : elle modélise un plan. Plante maintenant un crayon verticalement sur la feuille. Le crayon est perpendiculaire à \emph{tous} les segments que tu peux tracer sur la feuille, dans toutes les directions. Ce crayon matérialise ce que les mathématiciens appellent un \cle{vecteur normal} au plan. Remarque qu'il n'y a qu'une seule direction possible pour ce crayon (à l'orientation près) : tous les vecteurs normaux d'un même plan sont colinéaires.

\begin{definition}[Vecteur normal à un plan]
Soit $\mathcal{P}$ un plan de l'espace. On appelle \cle{vecteur normal} au plan $\mathcal{P}$ tout vecteur $\vec{n}$ \textbf{non nul} qui est orthogonal à tout vecteur directeur de $\mathcal{P}$.

Si $\mathcal{P}$ est dirigé par les vecteurs non colinéaires $\vec{u}$ et $\vec{v}$, alors $\vec{n}$ est normal à $\mathcal{P}$ si et seulement si $\vec{n}$ est orthogonal à $\vec{u}$ \emph{et} à $\vec{v}$.
\end{definition}

\begin{attention}[Orthogonalité dans l'espace : notion admise en Première]
En Première, l'orthogonalité de deux vecteurs de l'espace est utilisée de manière intuitive (directions perpendiculaires) : elle sera formalisée rigoureusement en Terminale grâce au produit scalaire. Retiens néanmoins dès maintenant deux faits essentiels :
\begin{itemize}
\item un plan admet une infinité de vecteurs normaux, tous \textbf{colinéaires entre eux} ;
\item la donnée d'un point $A$ et d'un vecteur normal $\vec{n}$ détermine \textbf{un unique plan} : le plan passant par $A$ et perpendiculaire à la direction de $\vec{n}$.
\end{itemize}
\end{attention}

\courssub{Théorème fondamental : équation cartésienne d'un plan}

\begin{propriete}[Équation cartésienne d'un plan — théorème fondamental]
Dans un repère $(O\,;\vec{\imath},\vec{\jmath},\vec{k})$ de l'espace :
\begin{enumerate}
\item Tout plan $\mathcal{P}$ admet une équation, dite \cle{équation cartésienne}, de la forme
\[ ax + by + cz + d = 0, \]
où $a$, $b$, $c$, $d$ sont des réels avec $(a\,;b\,;c)$ \textbf{non tous nuls}. Le vecteur $\vec{n}(a\,;b\,;c)$ est alors un vecteur normal à $\mathcal{P}$.
\item Réciproquement (admise), tout ensemble de points $M(x\,;y\,;z)$ vérifiant une équation $ax + by + cz + d = 0$ avec $(a\,;b\,;c)$ non tous nuls est un plan, de vecteur normal $\vec{n}(a\,;b\,;c)$.
\end{enumerate}
La démonstration du sens direct est formatrice et exigible au Probatoire : elle est guidée ci-dessous.
\end{propriete}

\begin{experience}[Démonstration guidée du sens direct : de la coplanarité à l'équation]
Soit $\mathcal{P}$ le plan passant par $A(x_A\,;y_A\,;z_A)$ et dirigé par les vecteurs non colinéaires $\vec{u}(\alpha\,;\beta\,;\gamma)$ et $\vec{v}(\alpha'\,;\beta'\,;\gamma')$. Soit $M(x\,;y\,;z)$ un point de l'espace.

\textbf{Étape 1 — Traduction vectorielle.} $M \in \mathcal{P}$ si et seulement si les vecteurs $\overrightarrow{AM}$, $\vec{u}$ et $\vec{v}$ sont coplanaires, c'est-à-dire si et seulement si
\[ \det\left(\overrightarrow{AM},\ \vec{u},\ \vec{v}\right) = 0. \]

\textbf{Étape 2 — Calcul du déterminant.} On a $\overrightarrow{AM}(x - x_A\,;\,y - y_A\,;\,z - z_A)$. En développant le déterminant selon la première colonne :
\[
\det\left(\overrightarrow{AM},\vec{u},\vec{v}\right)
= (x - x_A)\begin{vmatrix} \beta & \beta' \\ \gamma & \gamma' \end{vmatrix}
- (y - y_A)\begin{vmatrix} \alpha & \alpha' \\ \gamma & \gamma' \end{vmatrix}
+ (z - z_A)\begin{vmatrix} \alpha & \alpha' \\ \beta & \beta' \end{vmatrix}.
\]

\textbf{Étape 3 — Forme obtenue.} Posons
\[ a = \beta\gamma' - \gamma\beta', \qquad b = \gamma\alpha' - \alpha\gamma', \qquad c = \alpha\beta' - \beta\alpha', \]
et $d = -ax_A - by_A - cz_A$. L'équation s'écrit alors $ax + by + cz + d = 0$.

\textbf{Étape 4 — Les coefficients ne sont pas tous nuls.} Si l'on avait $a = b = c = 0$, alors on aurait $\beta\gamma' = \gamma\beta'$, $\gamma\alpha' = \alpha\gamma'$ et $\alpha\beta' = \beta\alpha'$ : les coordonnées de $\vec{u}$ et $\vec{v}$ seraient proportionnelles, donc $\vec{u}$ et $\vec{v}$ seraient colinéaires, ce qui est exclu par hypothèse. De plus, le vecteur $\vec{n}(a\,;b\,;c)$ ainsi construit est orthogonal à $\vec{u}$ et à $\vec{v}$ : c'est bien un vecteur normal au plan. $\blacksquare$
\end{experience}

\begin{aretenir}[Lecture immédiate d'un vecteur normal]
Dès qu'un plan est donné par une équation cartésienne, tu lis \textbf{gratuitement} un vecteur normal dans les coefficients de $x$, $y$ et $z$ :
\[ \mathcal{P} : ax + by + cz + d = 0 \quad\Longrightarrow\quad \vec{n}(a\,;b\,;c) \perp \mathcal{P}. \]
Le coefficient $d$, lui, ne renseigne pas sur la direction du plan mais sur sa \emph{position} (il dépend du point de passage).
\end{aretenir}

\begin{exemplebox}[Lecture d'un vecteur normal]
Le plan $\mathcal{P}$ d'équation $2x + y + 2z - 4 = 0$ admet pour vecteur normal $\vec{n}(2\,;1\,;2)$. Le point $A(2\,;0\,;0)$ appartient-il à $\mathcal{P}$ ? On vérifie : $2(2) + 0 + 2(0) - 4 = 0$. Oui, $A \in \mathcal{P}$. De même $B(0\,;4\,;0)$ et $C(0\,;0\,;2)$ appartiennent à $\mathcal{P}$, car $2(0) + 4 + 2(0) - 4 = 0$ et $2(0) + 0 + 2(2) - 4 = 0$.
\end{exemplebox}

\begin{popfigure}
\begin{tikzpicture}[scale=1.05, line cap=round, line join=round]
% Plan 2x + y + 2z = 4 : trace un parallélogramme suggérant le plan
\coordinate (P1) at (-2.6,-1.3);
\coordinate (P2) at (2.8,-0.6);
\coordinate (P3) at (3.6,1.6);
\coordinate (P4) at (-1.8,0.9);
\fill[popblueL, opacity=0.55] (P1) -- (P2) -- (P3) -- (P4) -- cycle;
\draw[popblue, thick] (P1) -- (P2) -- (P3) -- (P4) -- cycle;
% Axes
\draw[->, popmuted] (-3,0) -- (4.2,0) node[below left] {$x$};
\draw[->, popmuted] (0,-1.8) -- (0,3.4) node[left] {$z$};
\draw[->, popmuted] (0,0) -- (-1.9,-1.5) node[below] {$y$};
\node[below left, popmuted] at (0,0) {$O$};
% Point A(2;0;0) du plan
\coordinate (A) at (2,0);
\fill[popdark] (A) circle (1.6pt) node[below right, popdark] {$A(2\,;0\,;0)$};
% Vecteur normal n(2;1;2) depuis A
\draw[->, very thick, poporange] (A) -- ++(1.1,1.7) node[above, poporange] {$\vec{n}(2\,;1\,;2)$};
\draw[poporange, dashed] (A) ++(0.28,0.43) -- ++(0.14,-0.09) -- ++(-0.28,-0.43);
\node[popblue] at (-2.2,0.4) {$\mathcal{P} : 2x + y + 2z - 4 = 0$};
\end{tikzpicture}
\legende{Le plan $\mathcal{P} : 2x + y + 2z - 4 = 0$ et son vecteur normal $\vec{n}(2\,;1\,;2)$ tracé depuis le point $A(2\,;0\,;0)$ du plan.}
\end{popfigure}

\courssub{Passage entre représentation paramétrique et équation cartésienne}

\begin{methode}[De la représentation paramétrique à l'équation cartésienne]
Soit un plan donné par $A(x_A\,;y_A\,;z_A)$ et deux vecteurs directeurs $\vec{u}$, $\vec{v}$ non colinéaires, c'est-à-dire par un système paramétrique en $t$ et $s$. Pour obtenir une équation cartésienne :
\begin{enumerate}
\item Écrire que $M(x\,;y\,;z) \in \mathcal{P} \iff \det(\overrightarrow{AM}, \vec{u}, \vec{v}) = 0$ ;
\item Développer le déterminant : on obtient directement $ax + by + cz + d = 0$ ;
\item \emph{Variante} : éliminer $t$ et $s$ entre les trois équations paramétriques (par substitution).
\end{enumerate}
\end{methode}

\begin{methode}[De l'équation cartésienne à la représentation paramétrique]
Soit $\mathcal{P} : ax + by + cz + d = 0$ avec, par exemple, $a \neq 0$.
\begin{enumerate}
\item Isoler une variable dont le coefficient est non nul : $x = -\dfrac{by + cz + d}{a}$ ;
\item Choisir les deux autres variables comme paramètres libres : $y = t$, $z = s$ ;
\item Écrire le système paramétrique et en déduire un point du plan (par exemple pour $t = s = 0$) et deux vecteurs directeurs.
\end{enumerate}
\end{methode}

\begin{exoresolu}[Paramétrique $\leftrightarrow$ cartésienne]
\textbf{Énoncé.} (1) Le plan $\mathcal{P}$ passe par $A(1\,;-1\,;2)$ et est dirigé par $\vec{u}(1\,;1\,;0)$ et $\vec{v}(0\,;1\,;1)$. Déterminer une équation cartésienne de $\mathcal{P}$. (2) Donner une représentation paramétrique du plan $\mathcal{Q} : x - 2y + z - 3 = 0$.
\tcblower
\textbf{Solution.} (1) $M(x\,;y\,;z) \in \mathcal{P} \iff \det(\overrightarrow{AM}, \vec{u}, \vec{v}) = 0$ avec $\overrightarrow{AM}(x-1\,;\,y+1\,;\,z-2)$. En développant selon la première colonne :
\[
\begin{vmatrix} x-1 & 1 & 0 \\ y+1 & 1 & 1 \\ z-2 & 0 & 1 \end{vmatrix}
= (x-1)\begin{vmatrix} 1 & 1 \\ 0 & 1 \end{vmatrix} - (y+1)\begin{vmatrix} 1 & 0 \\ 0 & 1 \end{vmatrix} + (z-2)\begin{vmatrix} 1 & 0 \\ 1 & 1 \end{vmatrix}
= (x-1) - (y+1) + (z-2).
\]
D'où $\mathcal{P} : x - y + z - 4 = 0$. Vérification : $A$ satisfait l'équation ($1 - (-1) + 2 - 4 = 0$), et le vecteur normal lu $\vec{n}(1\,;-1\,;1)$ est bien orthogonal à $\vec{u}$ et à $\vec{v}$.

(2) On isole $x = 2y - z + 3$ et on pose $y = t$, $z = s$ :
\[
\begin{cases} x = 3 + 2t - s \\ y = t \\ z = s \end{cases}, \quad (t,s) \in \mathbb{R}^2.
\]
Le plan $\mathcal{Q}$ passe par $B(3\,;0\,;0)$ (obtenu pour $t = s = 0$) et est dirigé par $\vec{u_1}(2\,;1\,;0)$ et $\vec{v_1}(-1\,;0\,;1)$.
\end{exoresolu}

\courssub{Plans particuliers}

\begin{propriete}[Plans remarquables]
Dans un repère de l'espace :
\begin{itemize}
\item un plan d'équation $x = k$ est \textbf{parallèle au plan $(O\,;\vec{\jmath},\vec{k})$} (vecteur normal $\vec{\imath}$) ; de même $y = k$ définit un plan parallèle à $(O\,;\vec{\imath},\vec{k})$ et $z = k$ un plan parallèle à $(O\,;\vec{\imath},\vec{\jmath})$ ;
\item un plan passe \textbf{par l'origine} $O$ si et seulement si $d = 0$ dans son équation $ax + by + cz + d = 0$ ;
\item un plan est \textbf{parallèle à un axe de coordonnées} si le coefficient de la variable correspondante est nul : par exemple $2x + 3y - 6 = 0$ (coefficient de $z$ nul) définit un plan parallèle à l'axe $(Oz)$.
\end{itemize}
\end{propriete}

\begin{exemplebox}[Plans de coordonnées à l'entrepôt]
Imagine un entrepôt de stockage d'ignames à Bafoussam : le sol est le plan $z = 0$, le mur Est est le plan $x = 0$, le mur Nord le plan $y = 0$. Une étagère horizontale à \SI{1,5}{m} du sol est le plan $z = 1{,}5$ : elle est parallèle au sol, et le vecteur $\vec{k}(0\,;0\,;1)$ (la verticale) est normal aux deux.
\end{exemplebox}

\courssub{Équation d'un plan passant par trois points non alignés}

\begin{methode}[Déterminer une équation cartésienne du plan (ABC)]
Soient $A$, $B$, $C$ trois points \textbf{non alignés} (vérifier d'abord que $\overrightarrow{AB}$ et $\overrightarrow{AC}$ ne sont pas colinéaires !). Deux stratégies :
\begin{enumerate}
\item \emph{Par le déterminant} : $M(x\,;y\,;z) \in (ABC) \iff \det(\overrightarrow{AM}, \overrightarrow{AB}, \overrightarrow{AC}) = 0$, puis développer ;
\item \emph{Par le système} : écrire que $A$, $B$, $C$ vérifient $ax + by + cz + d = 0$, ce qui donne un système de 3 équations reliant $(a, b, c, d)$ ; fixer une inconnue non nulle (par exemple $d$ ou $a$) et résoudre.
\end{enumerate}
Dans les deux cas, conclure en vérifiant que les trois points satisfont l'équation trouvée.
\end{methode}

\begin{exoresolu}[Équation du plan (ABC) et position relative]
\textbf{Énoncé.} On donne $A(1\,;0\,;0)$, $B(0\,;2\,;0)$ et $C(0\,;0\,;3)$.
\begin{enumerate}
\item Déterminer une équation cartésienne du plan $(ABC)$.
\item Étudier la position relative de $(ABC)$ et du plan $\mathcal{Q} : 6x + 3y + 2z - 12 = 0$.
\end{enumerate}
\tcblower
\textbf{Solution.} 1. $\overrightarrow{AB}(-1\,;2\,;0)$ et $\overrightarrow{AC}(-1\,;0\,;3)$ ne sont pas colinéaires (leurs coordonnées ne sont pas proportionnelles), donc $A$, $B$, $C$ définissent bien un plan. Cherchons son équation sous la forme $ax + by + cz + d = 0$ en écrivant que les trois points la vérifient :
\[
A \in (ABC) : a + d = 0, \qquad B \in (ABC) : 2b + d = 0, \qquad C \in (ABC) : 3c + d = 0.
\]
Choisissons $d = -6$ : alors $a = 6$, $b = 3$, $c = 2$. D'où
\[ (ABC) : 6x + 3y + 2z - 6 = 0. \]
Vérification : $A : 6 - 6 = 0$ ; $B : 6 - 6 = 0$ ; $C : 6 - 6 = 0$. Un vecteur normal est $\vec{n}(6\,;3\,;2)$.

2. Le plan $\mathcal{Q}$ a pour vecteur normal $\vec{n'}(6\,;3\,;2)$. Or $\vec{n'} = \vec{n}$ : les vecteurs normaux sont colinéaires, donc les plans sont parallèles. Sont-ils confondus ? Testons $A(1\,;0\,;0)$ dans l'équation de $\mathcal{Q}$ : $6(1) - 12 = -6 \neq 0$, donc $A \notin \mathcal{Q}$. Les plans $(ABC)$ et $\mathcal{Q}$ sont \textbf{strictement parallèles}.
\end{exoresolu}

\courssub{Positions relatives de deux plans}

\begin{propriete}[Parallélisme et intersection de deux plans]
Soient $\mathcal{P} : ax + by + cz + d = 0$ et $\mathcal{P}' : a'x + b'y + c'z + d' = 0$, de vecteurs normaux respectifs $\vec{n}(a\,;b\,;c)$ et $\vec{n'}(a'\,;b'\,;c')$.
\begin{itemize}
\item $\mathcal{P}$ et $\mathcal{P}'$ sont \textbf{parallèles} si et seulement si $\vec{n}$ et $\vec{n'}$ sont \textbf{colinéaires}. Dans ce cas :
  \begin{itemize}
  \item si les quadruplets $(a,b,c,d)$ et $(a',b',c',d')$ sont proportionnels, les plans sont \textbf{confondus} ;
  \item sinon ils sont \textbf{strictement parallèles} (aucun point commun).
  \end{itemize}
\item $\mathcal{P}$ et $\mathcal{P}'$ sont \textbf{sécants} si et seulement si $\vec{n}$ et $\vec{n'}$ ne sont \textbf{pas colinéaires} ; leur intersection est alors une \textbf{droite}.
\end{itemize}
\end{propriete}

\begin{experience}[Démonstration guidée : pourquoi des normaux colinéaires impliquent des plans parallèles]
\textbf{Sens direct.} Supposons $\vec{n'} = \lambda \vec{n}$ avec $\lambda \neq 0$. Alors $a' = \lambda a$, $b' = \lambda b$, $c' = \lambda c$, et l'équation de $\mathcal{P}'$ s'écrit $\lambda(ax + by + cz) + d' = 0$, soit, en divisant par $\lambda \neq 0$ : $ax + by + cz + \dfrac{d'}{\lambda} = 0$. Posons $d'' = \dfrac{d'}{\lambda}$.
\begin{itemize}
\item Si $d'' = d$, les deux équations sont équivalentes : $\mathcal{P} = \mathcal{P}'$ (plans confondus).
\item Si $d'' \neq d$, un point $(x\,;y\,;z)$ ne peut vérifier simultanément $ax + by + cz = -d$ et $ax + by + cz = -d''$ : l'intersection est vide, les plans sont strictement parallèles.
\end{itemize}
\textbf{Réciproque (admise).} Si $\vec{n}$ et $\vec{n'}$ ne sont pas colinéaires, le système des deux équations admet des solutions dépendant d'un paramètre libre : l'intersection est une droite, dont la direction est orthogonale à la fois à $\vec{n}$ et à $\vec{n'}$. $\blacksquare$
\end{experience}

\begin{attention}[Piège classique : proportionnalité des coefficients]
Des coefficients $(a\,;b\,;c)$ proportionnels donnent des plans \textbf{parallèles}, mais ce n'est pas fini : il faut encore comparer le coefficient constant $d$ !

Exemple : $\mathcal{P} : x - y + 2z - 3 = 0$ et $\mathcal{P}' : 2x - 2y + 4z - 6 = 0$. On a $(2\,;-2\,;4) = 2(1\,;-1\,;2)$ \emph{et} $-6 = 2 \times (-3)$ : les plans sont \textbf{confondus}. En revanche, avec $\mathcal{P}'' : 2x - 2y + 4z + 1 = 0$, les plans sont \textbf{strictement parallèles}. Rédaction attendue au Probatoire : exhiber un point de $\mathcal{P}$ et montrer qu'il n'appartient pas à $\mathcal{P}''$.
\end{attention}

\begin{popfigure}
\begin{tikzpicture}[scale=0.92, line cap=round]
% Deux plans sécants
\begin{scope}[xshift=-4.2cm]
\coordinate (A1) at (-1.9,-1.1); \coordinate (A2) at (1.7,-0.5);
\coordinate (A3) at (2.3,1.3); \coordinate (A4) at (-1.3,0.7);
\fill[popblueL, opacity=0.6] (A1)--(A2)--(A3)--(A4)--cycle;
\draw[popblue] (A1)--(A2)--(A3)--(A4)--cycle;
\coordinate (B1) at (-1.5,-1.4); \coordinate (B2) at (1.9,-1.0);
\coordinate (B3) at (1.3,1.5); \coordinate (B4) at (-2.1,1.1);
\fill[poporangeL, opacity=0.55] (B1)--(B2)--(B3)--(B4)--cycle;
\draw[poporange] (B1)--(B2)--(B3)--(B4)--cycle;
\draw[very thick, popdark] (-1.75,-0.12) -- (2.05,0.28) node[right, popdark] {$\Delta$};
\node[popblue] at (-1.55,1.05) {$\mathcal{P}$};
\node[poporange] at (1.5,-1.35) {$\mathcal{P}'$};
\node[align=center] at (0,-2.1) {\small Plans sécants : intersection = droite $\Delta$};
\end{scope}
% Deux plans parallèles
\begin{scope}[xshift=4.2cm]
\coordinate (C1) at (-1.9,-0.4); \coordinate (C2) at (1.7,0.2);
\coordinate (C3) at (2.3,2.0); \coordinate (C4) at (-1.3,1.4);
\fill[popblueL, opacity=0.6] (C1)--(C2)--(C3)--(C4)--cycle;
\draw[popblue] (C1)--(C2)--(C3)--(C4)--cycle;
\coordinate (D1) at (-1.9,-1.6); \coordinate (D2) at (1.7,-1.0);
\coordinate (D3) at (2.3,0.8); \coordinate (D4) at (-1.3,0.2);
\fill[popgreenL, opacity=0.6] (D1)--(D2)--(D3)--(D4)--cycle;
\draw[popgreen] (D1)--(D2)--(D3)--(D4)--cycle;
\draw[->, thick, popdark] (0.2,1.2) -- (0.2,0.0);
\draw[->, thick, popdark] (0.2,0.0) -- (0.2,-1.0) node[below, popdark] {$\vec{n}$};
\node[popblue] at (-1.6,1.7) {$\mathcal{P}$};
\node[popgreen] at (-1.6,-1.35) {$\mathcal{P}'$};
\node[align=center] at (0.2,-2.1) {\small Plans strictement parallèles :\\ \small vecteurs normaux colinéaires};
\end{scope}
\end{tikzpicture}
\legende{À gauche : deux plans sécants et leur droite d'intersection $\Delta$. À droite : deux plans strictement parallèles, un même vecteur normal $\vec{n}$ leur est orthogonal.}
\end{popfigure}

\courssub{Intersection de deux plans sécants}

\begin{methode}[Déterminer la droite d'intersection de deux plans sécants]
Soient $\mathcal{P} : ax + by + cz + d = 0$ et $\mathcal{P}' : a'x + b'y + c'z + d' = 0$ sécants.
\begin{enumerate}
\item Vérifier que les vecteurs normaux ne sont pas colinéaires (sinon les plans sont parallèles et il n'y a pas de droite d'intersection) ;
\item Résoudre le système des deux équations en choisissant une variable comme paramètre (par exemple $z = t$), après avoir vérifié que ce choix est licite ;
\item Exprimer les deux autres variables en fonction de $t$ : on obtient une représentation paramétrique de la droite $\Delta = \mathcal{P} \cap \mathcal{P}'$ ;
\item En déduire un point de $\Delta$ (prendre $t = 0$) et un vecteur directeur.
\end{enumerate}
\end{methode}

\begin{exoresolu}[Droite d'intersection de deux plans]
\textbf{Énoncé.} Soient $\mathcal{P} : x + y - z + 1 = 0$ et $\mathcal{P}' : 2x - y + z - 4 = 0$. Montrer que $\mathcal{P}$ et $\mathcal{P}'$ sont sécants et déterminer une représentation paramétrique de leur droite d'intersection $\Delta$.
\tcblower
\textbf{Solution.} Les vecteurs normaux sont $\vec{n}(1\,;1\,;-1)$ et $\vec{n'}(2\,;-1\,;1)$. Ils ne sont pas colinéaires car $\dfrac{2}{1} = 2$ alors que $\dfrac{-1}{1} = -1 \neq 2$ : les plans sont donc sécants et leur intersection est une droite $\Delta$.

Posons $z = t$ et résolvons le système :
\[
\begin{cases} x + y = t - 1 \\ 2x - y = 4 - t \end{cases}
\]
En additionnant membre à membre : $3x = 3$, donc $x = 1$. Puis $y = t - 1 - x = t - 2$. D'où :
\[
\Delta : \begin{cases} x = 1 \\ y = -2 + t \\ z = t \end{cases}, \quad t \in \mathbb{R}.
\]
La droite $\Delta$ passe par le point $I(1\,;-2\,;0)$ et est dirigée par $\vec{w}(0\,;1\,;1)$. Vérification : $I \in \mathcal{P}$ car $1 + (-2) - 0 + 1 = 0$, et $I \in \mathcal{P}'$ car $2(1) - (-2) + 0 - 4 = 0$.
\end{exoresolu}

\begin{exercice}[À toi de jouer]
On considère les plans $\mathcal{P}_1 : x - 2y + z - 1 = 0$, $\mathcal{P}_2 : 3x - 6y + 3z - 3 = 0$ et $\mathcal{P}_3 : 2x - 4y + 2z + 5 = 0$.
\begin{enumerate}
\item Déterminer la position relative de $\mathcal{P}_1$ et $\mathcal{P}_2$, puis de $\mathcal{P}_1$ et $\mathcal{P}_3$.
\item Déterminer une équation cartésienne du plan passant par $A(0\,;1\,;-1)$, $B(2\,;0\,;1)$ et $C(1\,;1\,;0)$.
\end{enumerate}
\end{exercice}

\begin{corrige}[Exercice 1]
1. $\mathcal{P}_2$ a pour coefficients $(3\,;-6\,;3\,;-3) = 3 \times (1\,;-2\,;1\,;-1)$ : les quadruplets sont proportionnels, donc $\mathcal{P}_1$ et $\mathcal{P}_2$ sont \textbf{confondus}. Pour $\mathcal{P}_3$ : $(2\,;-4\,;2) = 2(1\,;-2\,;1)$ mais $5 \neq 2 \times (-1) = -2$ : les plans sont \textbf{strictement parallèles} (par exemple $M(1\,;0\,;0) \in \mathcal{P}_1$ car $1 - 1 = 0$, mais $2(1) + 5 = 7 \neq 0$ donc $M \notin \mathcal{P}_3$).

2. $\overrightarrow{AB}(2\,;-1\,;2)$ et $\overrightarrow{AC}(1\,;0\,;1)$ ne sont pas colinéaires (coordonnées non proportionnelles), donc $A$, $B$, $C$ définissent un plan. Le système donné par l'appartenance de $A$, $B$, $C$ au plan d'équation $ax + by + cz + d = 0$ s'écrit : $b - c + d = 0$, $2a + c + d = 0$, $a + b + d = 0$. En choisissant $d = -1$ : $b - c = 1$, $2a + c = 1$, $a + b = 1$. De la première équation $b = 1 + c$ ; la troisième donne $a = 1 - b = -c$ ; la deuxième donne alors $-2c + c = 1$, soit $c = -1$, puis $a = 1$ et $b = 0$. D'où le plan : $x - z - 1 = 0$. Vérification : $A : 0 - (-1) - 1 = 0$ ; $B : 2 - 1 - 1 = 0$ ; $C : 1 - 0 - 1 = 0$.
\end{corrige}

\begin{savaistu}[Des plans aux réseaux mobiles]
Quand un ingénieur d'Orange ou de MTN planifie la couverture d'une ville comme Douala, il modélise les fronts d'onde émis par les antennes relais par des plans de l'espace, et les zones d'interférence entre deux antennes correspondent précisément à l'\emph{intersection de deux plans} : une droite le long de laquelle le signal doit être réajusté. Les équations cartésiennes $ax + by + cz + d = 0$ sont aussi à la base des logiciels de cartographie 3D utilisés pour l'électrification rurale : chaque façade de bâtiment, chaque versant de colline y est stocké sous forme d'équation de plan.
\end{savaistu}

\begin{aretenir}[L'essentiel de la section]
\begin{itemize}
\item Un plan est caractérisé par une équation cartésienne $ax + by + cz + d = 0$ avec $(a\,;b\,;c) \neq (0\,;0\,;0)$ ; le vecteur $\vec{n}(a\,;b\,;c)$ est \textbf{normal} au plan.
\item Pour trouver l'équation du plan $(ABC)$ : $\det(\overrightarrow{AM}, \overrightarrow{AB}, \overrightarrow{AC}) = 0$, ou un système obtenu en écrivant que $A$, $B$, $C$ vérifient l'équation.
\item Deux plans sont \textbf{parallèles} $\iff$ leurs vecteurs normaux sont colinéaires ; comparer ensuite le coefficient constant $d$ pour distinguer \textbf{confondus} et \textbf{strictement parallèles}.
\item Deux plans \textbf{sécants} se coupent selon une \textbf{droite}, obtenue en résolvant le système des deux équations avec un paramètre libre.
\end{itemize}
\end{aretenir}

\coursec{Projetés orthogonaux et distances}

Dans la section précédente, nous avons appris à décrire un plan par une équation cartésienne $ax + by + cz + d = 0$ et à lire sur cette équation un \cle{vecteur normal} $\vec{n}\begin{pmatrix} a \\ b \\ c \end{pmatrix}$. Ce vecteur normal est précisément l'outil qui va nous permettre de répondre à une question très concrète : \emph{quelle est la plus courte distance entre un point et un plan ?} Un technicien qui pose un panneau solaire veut savoir à quelle distance du toit se trouve le câble ; un ingénieur de la SNH veut connaître la profondeur d'un forage par rapport à une couche géologique. Dans les deux cas, il s'agit de \emph{projeter orthogonalement} un point sur un plan ou sur une droite.

Dans toute cette section, l'espace est muni d'un repère \textbf{orthonormé} $(O\,;\,\vec{i},\vec{j},\vec{k})$. Nous utiliserons librement le fait suivant, \textbf{admis} en Première (il sera justifié en Terminale par le produit scalaire) : deux vecteurs $\vec{u}\begin{pmatrix} x \\ y \\ z \end{pmatrix}$ et $\vec{v}\begin{pmatrix} x' \\ y' \\ z' \end{pmatrix}$ sont orthogonaux si et seulement si
\[ x\,x' + y\,y' + z\,z' = 0. \]

\courssub{Projeté orthogonal d'un point sur une droite}

\textbf{Intuition.} Imagine que tu te tiens au bord d'une route rectiligne (la droite $d$) et que tu veux la traverser au plus court. Tu marches perpendiculairement à la route : le point où tu arrives est ton \emph{projeté orthogonal} sur la droite.

\begin{definition}[Projeté orthogonal sur une droite]
Soit $d$ une droite de l'espace et $M$ un point. Le \cle{projeté orthogonal} de $M$ sur $d$ est le point $H$ de $d$ tel que la droite $(MH)$ soit orthogonale à $d$.
\begin{itemize}
\item Si $M \in d$, alors $H = M$.
\item Si $M \notin d$, $H$ est l'unique point de $d$ vérifiant $\overrightarrow{MH} \perp \vec{u}$, où $\vec{u}$ est un vecteur directeur de $d$.
\end{itemize}
L'existence et l'unicité de $H$ sont \textbf{admises} (elles seront démontrées en Terminale à l'aide du produit scalaire).
\end{definition}

\begin{definition}[Distance d'un point à une droite]
La \cle{distance} du point $M$ à la droite $d$, notée $d(M, d)$, est la longueur $MH$, où $H$ est le projeté orthogonal de $M$ sur $d$ :
\[ d(M, d) = MH. \]
C'est la \textbf{plus petite} des distances $MK$ lorsque $K$ décrit $d$ : en effet, pour tout point $K$ de $d$ distinct de $H$, le triangle $MHK$ est rectangle en $H$, donc $MK^2 = MH^2 + HK^2 > MH^2$.
\end{definition}

\begin{methode}[Calculer la distance d'un point à une droite]
Soit $d$ passant par $A$ et dirigée par $\vec{u}$, et $M$ un point.
\begin{enumerate}
\item \textbf{Paramétrer} le point courant de $d$ : $H(t) = A + t\,\vec{u}$, c'est-à-dire écrire les coordonnées de $H$ en fonction du paramètre $t \in \mathbb{R}$.
\item \textbf{Traduire l'orthogonalité} : écrire que $\overrightarrow{MH} \perp \vec{u}$, c'est-à-dire que la somme des produits des coordonnées correspondantes est nulle (condition admise, rappelée en introduction). On obtient une équation du premier degré en $t$, que l'on résout.
\item \textbf{Calculer} $MH$ avec les coordonnées de $H$ trouvées :
\[ MH = \sqrt{(x_H - x_M)^2 + (y_H - y_M)^2 + (z_H - z_M)^2}. \]
\end{enumerate}
\end{methode}

\begin{exoresolu}[Distance d'un point à une droite]
Dans un repère orthonormé, on donne $M(2\,;\,1\,;\,3)$ et la droite $d$ passant par $A(0\,;\,0\,;\,1)$ et dirigée par $\vec{u}\begin{pmatrix} 1 \\ 1 \\ 0 \end{pmatrix}$. Déterminer le projeté orthogonal $H$ de $M$ sur $d$ et la distance $d(M, d)$.
\tcblower
\textbf{Solution.} \textbf{Étape 1.} Un point courant de $d$ s'écrit $H(t) = (t\,;\,t\,;\,1)$, $t \in \mathbb{R}$. Alors
\[ \overrightarrow{MH}\begin{pmatrix} t - 2 \\ t - 1 \\ 1 - 3 \end{pmatrix} = \begin{pmatrix} t - 2 \\ t - 1 \\ -2 \end{pmatrix}. \]

\textbf{Étape 2.} La condition d'orthogonalité $\overrightarrow{MH} \perp \vec{u}$ s'écrit :
\[ (t-2)\times 1 + (t-1)\times 1 + (-2)\times 0 = 0 \iff 2t - 3 = 0 \iff t = \dfrac{3}{2}. \]
Donc $H\left(\dfrac{3}{2}\,;\,\dfrac{3}{2}\,;\,1\right)$.

\textbf{Étape 3.} On calcule ensuite :
\[ MH = \sqrt{\left(\tfrac{3}{2}-2\right)^2 + \left(\tfrac{3}{2}-1\right)^2 + (1-3)^2} = \sqrt{\tfrac{1}{4} + \tfrac{1}{4} + 4} = \sqrt{\tfrac{9}{2}} = \dfrac{3}{\sqrt{2}} = \dfrac{3\sqrt{2}}{2}. \]
Ainsi $d(M, d) = \dfrac{3\sqrt{2}}{2} \approx 2{,}12$.
\end{exoresolu}

\courssub{Projeté orthogonal d'un point sur un plan}

\textbf{Intuition.} Pose une lampe au-dessus d'une table : l'ombre du filament sur la table, lorsque la lumière arrive perpendiculairement, c'est le projeté orthogonal. Géométriquement, on « descend » du point $M$ vers le plan $P$ en suivant la direction du vecteur normal $\vec{n}$.

\begin{definition}[Projeté orthogonal sur un plan]
Soit $P$ un plan de vecteur normal $\vec{n}$ et $M$ un point de l'espace. Le \cle{projeté orthogonal} de $M$ sur $P$ est le point $H$, intersection du plan $P$ et de la droite passant par $M$ et dirigée par $\vec{n}$.
\begin{itemize}
\item Si $M \in P$, alors $H = M$.
\item Si $M \notin P$, la droite $(M\,;\,\vec{n})$ est perpendiculaire à $P$ et le coupe en un unique point $H$.
\end{itemize}
\end{definition}

\begin{propriete}[Existence et unicité du projeté orthogonal sur un plan]
Soit $P$ un plan, $\vec{n}$ un vecteur normal à $P$ et $M$ un point. La droite passant par $M$ et dirigée par $\vec{n}$ coupe $P$ en un \textbf{unique} point $H$ : le projeté orthogonal de $M$ sur $P$. De plus, pour tout point $K$ de $P$ distinct de $H$, on a $MK > MH$.
\par\emph{Cette propriété est admise en Première ; elle sera démontrée en Terminale grâce au produit scalaire. L'idée : comme $(MH)$ est perpendiculaire à $P$, elle est orthogonale à toute droite de $P$ passant par $H$, donc le triangle $MHK$ est rectangle en $H$, d'où $MK^2 = MH^2 + HK^2 > MH^2$.}
\end{propriete}

\begin{definition}[Distance d'un point à un plan]
La \cle{distance} du point $M$ au plan $P$, notée $d(M, P)$, est la longueur $MH$, où $H$ est le projeté orthogonal de $M$ sur $P$ :
\[ d(M, P) = MH = \min_{K \in P} MK. \]
\end{definition}

\begin{popfigure}
\begin{tikzpicture}[scale=1.05]
% Plan P vu en perspective (parallélogramme)
\fill[popblueL,opacity=0.55] (0,0) -- (5.2,0) -- (6.4,1.8) -- (1.2,1.8) -- cycle;
\draw[popblue,thick] (0,0) -- (5.2,0) -- (6.4,1.8) -- (1.2,1.8) -- cycle;
\node[popblue] at (5.9,0.35) {$P$};
% Projeté H
\coordinate (H) at (3.1,0.9);
% Point M au-dessus
\coordinate (M) at (3.1,3.6);
\draw[popink,very thick] (M) -- (H) node[midway,right] {$MH$};
% vecteur normal
\draw[-{Stealth},poporange,very thick] (M) -- (3.1,4.6) node[above] {$\vec{n}$};
% angle droit
\draw[popdark] (3.1,0.9) ++(0,0.28) -- ++(0.22,0.11) -- ++(0,-0.28);
\fill[popink] (M) circle (1.6pt) node[left] {$M$};
\fill[popgreen] (H) circle (1.6pt) node[below right] {$H$};
% un autre point K du plan
\coordinate (K) at (4.6,1.35);
\fill[popmuted] (K) circle (1.4pt) node[right] {$K$};
\draw[popmuted,dashed] (M) -- (K) node[midway,above right] {$MK > MH$};
\end{tikzpicture}
\legende{Le projeté orthogonal $H$ de $M$ sur $P$ : la perpendiculaire $(M\,;\,\vec{n})$ perce le plan en $H$, et $MH$ est la plus courte distance de $M$ au plan.}
\end{popfigure}

\begin{methode}[Calculer la distance d'un point à un plan en trois étapes]
Soit $P : ax + by + cz + d = 0$ et $M(x_M\,;\,y_M\,;\,z_M)$.
\begin{enumerate}
\item \textbf{Vecteur normal} : lire $\vec{n}\begin{pmatrix} a \\ b \\ c \end{pmatrix}$ sur l'équation de $P$.
\item \textbf{Paramétrage de la perpendiculaire} : écrire la représentation paramétrique de la droite $(M\,;\,\vec{n})$ :
\[ \begin{cases} x = x_M + at \\ y = y_M + bt \\ z = z_M + ct \end{cases} \quad (t \in \mathbb{R}). \]
\item \textbf{Intersection} : reporter ces expressions dans l'équation de $P$, résoudre l'équation en $t$, en déduire $H$, puis calculer $d(M, P) = MH$.
\end{enumerate}
\textbf{Astuce de calcul.} Comme $\overrightarrow{MH} = t\,\vec{n}$, on a $MH = |t| \times \|\vec{n}\| = |t|\sqrt{a^2+b^2+c^2}$ : on peut donc obtenir la distance sans calculer explicitement les coordonnées de $H$.
\end{methode}

\begin{exemplebox}[Distance de $M(1\,;\,2\,;\,3)$ au plan $P : x - 2y + 2z - 4 = 0$]
\textbf{Étape 1.} Un vecteur normal est $\vec{n}\begin{pmatrix} 1 \\ -2 \\ 2 \end{pmatrix}$.

\textbf{Étape 2.} La perpendiculaire à $P$ passant par $M$ a pour représentation paramétrique :
\[ \begin{cases} x = 1 + t \\ y = 2 - 2t \\ z = 3 + 2t \end{cases} \quad (t \in \mathbb{R}). \]

\textbf{Étape 3.} On injecte dans l'équation de $P$ :
\begin{align*}
(1+t) - 2(2-2t) + 2(3+2t) - 4 &= 0 \\
1 + t - 4 + 4t + 6 + 4t - 4 &= 0 \\
9t - 1 &= 0 \iff t = \tfrac{1}{9}.
\end{align*}
Le projeté orthogonal est $H\left(1+\tfrac{1}{9}\,;\,2-\tfrac{2}{9}\,;\,3+\tfrac{2}{9}\right)$, c'est-à-dire $H\left(\dfrac{10}{9}\,;\,\dfrac{16}{9}\,;\,\dfrac{29}{9}\right)$. Alors
\[ \overrightarrow{MH}\begin{pmatrix} 1/9 \\ -2/9 \\ 2/9 \end{pmatrix}, \qquad MH = \sqrt{\tfrac{1}{81} + \tfrac{4}{81} + \tfrac{4}{81}} = \sqrt{\tfrac{9}{81}} = \dfrac{3}{9} = \dfrac{1}{3}. \]
Donc $d(M, P) = \dfrac{1}{3}$. (Vérification par l'astuce : $MH = |t|\,\|\vec{n}\| = \dfrac{1}{9}\times\sqrt{1+4+4} = \dfrac{1}{9}\times 3 = \dfrac{1}{3}$.)

\textbf{Vérification (complément de Terminale).} La formule $d(M,P) = \dfrac{|ax_M + by_M + cz_M + d|}{\sqrt{a^2+b^2+c^2}}$ donne :
\[ d(M,P) = \dfrac{|1 - 4 + 6 - 4|}{\sqrt{1+4+4}} = \dfrac{|-1|}{3} = \dfrac{1}{3}. \checkmark \]
\end{exemplebox}

\begin{aretenir}[Formule de la distance point-plan (hors exigibilité immédiate)]
Si $P : ax + by + cz + d = 0$ et $M(x_M\,;\,y_M\,;\,z_M)$, alors
\[ d(M, P) = \dfrac{|a\,x_M + b\,y_M + c\,z_M + d|}{\sqrt{a^2 + b^2 + c^2}}. \]
Cette formule sera \textbf{démontrée en Terminale} à l'aide du produit scalaire. En Première, tu dois savoir retrouver le résultat par la méthode de la perpendiculaire en trois étapes ; la formule te sert de vérification rapide.
\end{aretenir}

\begin{attention}[La distance n'est pas n'importe quelle distance]
La distance d'un point $M$ à un plan $P$ n'est \textbf{pas} la distance de $M$ à un point quelconque de $P$. C'est le \cle{minimum} des distances $MK$ pour $K \in P$, atteint \textbf{uniquement} au projeté orthogonal $H$. Erreurs fréquentes :
\begin{itemize}
\item prendre un point $K$ « commode » du plan (par exemple l'intersection avec un axe) et calculer $MK$ : le résultat est toujours \emph{supérieur} à la vraie distance ;
\item oublier que la perpendiculaire doit être dirigée par un vecteur \emph{normal} à $P$, pas par un vecteur quelconque du plan ;
\item dans la formule de vérification, oublier la valeur absolue au numérateur ou la racine carrée au dénominateur ;
\item si $M \in P$, la distance est nulle : inutile de lancer les calculs, vérifier d'abord si les coordonnées de $M$ satisfont l'équation de $P$.
\end{itemize}
\end{attention}

\courssub{Applications : plans parallèles, hauteur et volume d'un tétraèdre}

\textbf{Distance entre deux plans strictement parallèles.} Soient $P$ et $P'$ deux plans strictement parallèles (même direction, équations non équivalentes). Tous les points de $P$ sont à la même distance de $P'$ : on définit
\[ d(P, P') = d(M, P') \quad \text{pour n'importe quel point } M \in P. \]
Il suffit donc de choisir un point facile de $P$ (par exemple en posant $y = z = 0$ dans son équation) et de calculer sa distance à $P'$.

\begin{exemplebox}[Distance entre deux plans parallèles]
Soient $P : x - 2y + 2z - 4 = 0$ et $P' : x - 2y + 2z + 2 = 0$. Ils ont le même vecteur normal $\vec{n}\begin{pmatrix} 1 \\ -2 \\ 2 \end{pmatrix}$, et leurs équations ne sont pas équivalentes (les constantes diffèrent : $-4 \neq 2$) : ils sont donc strictement parallèles. Le point $A(4\,;\,0\,;\,0)$ appartient à $P$ (car $4 - 0 + 0 - 4 = 0$). La perpendiculaire $(A\,;\,\vec{n})$ : $x = 4+t$, $y = -2t$, $z = 2t$ rencontre $P'$ lorsque
\[ (4+t) - 2(-2t) + 2(2t) + 2 = 0 \iff 9t + 6 = 0 \iff t = -\dfrac{2}{3}. \]
Grâce à l'astuce $MH = |t|\,\|\vec{n}\|$, la distance est
\[ d(P, P') = |t| \times \|\vec{n}\| = \dfrac{2}{3} \times \sqrt{1+4+4} = \dfrac{2}{3} \times 3 = 2. \]
\end{exemplebox}

\textbf{Hauteur et volume d'un tétraèdre.} Dans un tétraèdre $ABCD$, la \cle{hauteur} issue de $D$ est la distance de $D$ au plan $(ABC)$ : c'est $DH$, où $H$ est le projeté orthogonal de $D$ sur $(ABC)$.

\begin{methode}[Calculer le volume d'un tétraèdre]
\begin{enumerate}
\item Déterminer une équation cartésienne du plan de base, par exemple $(ABC)$, à l'aide d'un vecteur normal obtenu à partir de $\overrightarrow{AB}$ et $\overrightarrow{AC}$.
\item Calculer la hauteur $h = d(D, (ABC))$ par la méthode en trois étapes.
\item Calculer l'aire $\mathcal{A}$ du triangle de base $ABC$.
\item Conclure : $V = \dfrac{1}{3} \times \mathcal{A} \times h$.
\end{enumerate}
\end{methode}

\begin{popfigure}
\begin{tikzpicture}[scale=1.15]
% Base ABC
\coordinate (A) at (0,0);
\coordinate (B) at (4.6,0);
\coordinate (C) at (3.2,1.9);
\fill[popgreenL,opacity=0.6] (A) -- (B) -- (C) -- cycle;
\draw[popgreen,thick] (A) -- (B) -- (C) -- cycle;
% Sommet D et projeté H
\coordinate (H) at (2.5,0.75);
\coordinate (D) at (2.5,3.4);
\draw[popink,thick] (D) -- (A);
\draw[popink,thick] (D) -- (B);
\draw[popink,thick] (D) -- (C);
\draw[poporange,very thick,dashed] (D) -- (H) node[midway,right] {$h = d(D,(ABC))$};
\draw[popdark] (2.5,0.75) ++(0,0.24) -- ++(0.19,0.09) -- ++(0,-0.24);
\fill[popink] (A) circle (1.5pt) node[below left] {$A$};
\fill[popink] (B) circle (1.5pt) node[below right] {$B$};
\fill[popink] (C) circle (1.5pt) node[right] {$C$};
\fill[popink] (D) circle (1.5pt) node[above] {$D$};
\fill[poporange] (H) circle (1.5pt) node[below] {$H$};
\end{tikzpicture}
\legende{Tétraèdre $ABCD$ : la hauteur issue de $D$ est la distance de $D$ au plan de base $(ABC)$, et $V = \dfrac{1}{3}\,\mathcal{A}_{ABC}\times DH$.}
\end{popfigure}

\begin{exoresolu}[Volume d'un tétraèdre]
Dans un repère orthonormé, on donne $A(0\,;\,0\,;\,0)$, $B(2\,;\,0\,;\,0)$, $C(0\,;\,3\,;\,0)$ et $D(1\,;\,1\,;\,6)$. Calculer le volume du tétraèdre $ABCD$.
\tcblower
\textbf{Solution.} \textbf{Étape 1.} Les points $A$, $B$, $C$ ont tous une cote nulle : le plan $(ABC)$ a pour équation $z = 0$, de vecteur normal $\vec{n}\begin{pmatrix} 0 \\ 0 \\ 1 \end{pmatrix}$.

\textbf{Étape 2.} La perpendiculaire à $(ABC)$ passant par $D$ a pour représentation paramétrique $(x\,;\,y\,;\,z) = (1\,;\,1\,;\,6+t)$, $t \in \mathbb{R}$. Elle coupe le plan $z = 0$ pour $6 + t = 0$, soit $t = -6$, en $H(1\,;\,1\,;\,0)$. La hauteur est $DH = |{-6}|\times\|\vec{n}\| = 6$.

\textbf{Étape 3.} Le triangle $ABC$ est rectangle en $A$ : en effet $\overrightarrow{AB}\begin{pmatrix} 2 \\ 0 \\ 0 \end{pmatrix}$ et $\overrightarrow{AC}\begin{pmatrix} 0 \\ 3 \\ 0 \end{pmatrix}$ vérifient $2\times 0 + 0\times 3 + 0\times 0 = 0$, donc $\overrightarrow{AB} \perp \overrightarrow{AC}$. Ainsi
\[ \mathcal{A}_{ABC} = \dfrac{AB \times AC}{2} = \dfrac{2 \times 3}{2} = 3. \]

\textbf{Étape 4.} $V = \dfrac{1}{3} \times \mathcal{A}_{ABC} \times DH = \dfrac{1}{3} \times 3 \times 6 = 6$ unités de volume.
\end{exoresolu}

\begin{savaistu}[Des distances qui font tourner le monde]
Les systèmes \textbf{GPS} calculent en permanence des positions dans l'espace par triangulation : corriger les erreurs de trajectoire revient à mesurer des distances d'un point (le récepteur) à des plans ou des droites de référence. Au Cameroun, les ingénieurs de la \textbf{SNH} (Société Nationale des Hydrocarbures) et les équipes du bassin de \textbf{Logbaba} à Douala pilotent les forages pétroliers grâce à des logiciels qui calculent à chaque instant la distance entre le trépan (un point) et les couches géologiques visées (modélisées par des plans), ainsi que l'écart entre le puits réel et le puits théorique (distance d'un point à une droite). Les mathématiques de cette leçon sont littéralement au cœur de l'industrie pétrolière camerounaise !
\end{savaistu}

\begin{exercice}[Pour s'entraîner]
Dans un repère orthonormé, on donne $M(2\,;\,-1\,;\,4)$ et le plan $P : 2x + y - 2z + 3 = 0$.
\begin{enumerate}
\item Déterminer une représentation paramétrique de la perpendiculaire à $P$ passant par $M$.
\item En déduire le projeté orthogonal $H$ de $M$ sur $P$ et la distance $d(M, P)$.
\item Vérifier le résultat avec la formule de la distance (complément de Terminale).
\end{enumerate}
\end{exercice}

\begin{corrige}[Exercice]
\begin{enumerate}
\item Un vecteur normal est $\vec{n}\begin{pmatrix} 2 \\ 1 \\ -2 \end{pmatrix}$. La perpendiculaire à $P$ passant par $M$ a pour représentation paramétrique :
\[ \begin{cases} x = 2 + 2t \\ y = -1 + t \\ z = 4 - 2t \end{cases} \quad (t \in \mathbb{R}). \]
\item On injecte dans l'équation de $P$ :
\[ 2(2+2t) + (-1+t) - 2(4-2t) + 3 = 0 \iff 4 + 4t - 1 + t - 8 + 4t + 3 = 0 \iff 9t - 2 = 0, \]
donc $t = \dfrac{2}{9}$. Alors $H\left(2+\dfrac{4}{9}\,;\,-1+\dfrac{2}{9}\,;\,4-\dfrac{4}{9}\right)$, c'est-à-dire $H\left(\dfrac{22}{9}\,;\,-\dfrac{7}{9}\,;\,\dfrac{32}{9}\right)$, et
\[ d(M,P) = |t|\,\|\vec{n}\| = \dfrac{2}{9}\times\sqrt{4+1+4} = \dfrac{2}{9}\times 3 = \dfrac{2}{3}. \]
\item Formule : $d(M,P) = \dfrac{|2\times 2 + (-1) - 2\times 4 + 3|}{\sqrt{4+1+4}} = \dfrac{|4 - 1 - 8 + 3|}{3} = \dfrac{|-2|}{3} = \dfrac{2}{3}$. Les deux méthodes concordent.
\end{enumerate}
\end{corrige}

\newpage

\coursec{Activité d'intégration}

\courssub{Objectifs de l'activité}

Cette activité d'intégration te propose de jouer le rôle d'un technicien géomètre chargé de vérifier la sécurité d'une installation à Garoua. En travaillant en groupe, tu vas mobiliser \emph{tous} les savoir-faire de la leçon :

\begin{itemize}
\item placer des points et lire des coordonnées dans un repère de l'espace dessiné en perspective cavalière ;
\item déterminer une représentation paramétrique d'une droite définie par deux points ;
\item étudier la position relative de deux droites de l'espace (sécantes, parallèles ou non coplanaires) ;
\item déterminer une équation cartésienne d'un plan et un vecteur normal à ce plan ;
\item calculer la distance d'un point à un plan et la distance d'un point à une droite.
\end{itemize}

Au-delà des calculs, cette activité développe des savoir-être essentiels : la rigueur dans la rédaction d'un rapport technique, l'esprit d'équipe et le sens des responsabilités, car une erreur de calcul sur une ligne électrique peut coûter des vies.

\courssub{Situation}

La société \textbf{ENEO} et l'opérateur \textbf{CAMTEL} partagent un site technique à la périphérie de Garoua, dans la région du Nord. Sur ce site, deux pylônes supportent un câble téléphonique, une ligne électrique moyenne tension traverse la zone, et un panneau solaire alimente le relais de communication. Avant la mise en service, le bureau d'études doit vérifier que la ligne électrique ne croise pas le câble téléphonique et que les distances de sécurité sont respectées.

On modélise le site dans un repère orthonormal de l'espace $(O\,;\vec{\imath},\vec{\jmath},\vec{k})$, l'unité étant le mètre. L'axe $(O\,;\vec{\imath})$ pointe vers l'est, l'axe $(O\,;\vec{\jmath})$ vers le nord, et l'axe $(O\,;\vec{k})$ vers le haut (verticale du lieu).

\begin{itemize}
\item Les deux pylônes du câble téléphonique ont leurs sommets en $A(0\,;0\,;8)$ et $B(12\,;0\,;8)$. Le câble téléphonique est assimilé à la droite $(AB)$.
\item La ligne électrique est assimilée à la droite $(CD)$, passant par $C(2\,;5\,;10)$ et $D(14\,;5\,;6)$.
\item Le panneau solaire est assimilé au plan $\mathcal{P}$ passant par les trois points $E(0\,;10\,;0)$, $F(6\,;10\,;0)$ et $G(0\,;10\,;4)$.
\item Un technicien chargé de la maintenance se trouve au point $M(6\,;3\,;2)$.
\end{itemize}

Le bureau d'études te confie la mission suivante : produire un \textbf{rapport technique} répondant aux questions de sécurité posées par le chef de projet.

\courssub{Consignes}

Travail en groupes de quatre élèves. Chaque groupe rédigera un rapport technique soigné, comportant une figure et toutes les justifications.

\begin{enumerate}
\item \textbf{Repérage.} Sur une feuille, dessine un repère de l'espace en perspective cavalière (prends par exemple \SI{0,5}{cm} pour 1 m sur les axes $(Ox)$ et $(Oz)$, et \SI{0,25}{cm} avec un angle de $45^\circ$ pour l'axe $(Oy)$). Place les points $A$, $B$, $C$, $D$, $E$, $F$, $G$ et $M$.
\item \textbf{Les deux câbles.} Détermine une représentation paramétrique du câble téléphonique $(AB)$, puis une représentation paramétrique de la ligne électrique $(CD)$. Précise pour chacune un point et un vecteur directeur.
\item \textbf{Question de sécurité.} Les deux câbles se croisent-ils ? Étudie la position relative des droites $(AB)$ et $(CD)$ : sont-elles parallèles, sécantes ou non coplanaires ? Conclus sur la sécurité de l'installation.
\item \textbf{Le panneau solaire.} Justifie que les points $E$, $F$, $G$ définissent bien un plan. Détermine un vecteur normal $\vec{n}$ au plan $\mathcal{P}$, puis une équation cartésienne de $\mathcal{P}$.
\item \textbf{Hauteur du pylône par rapport au panneau.} Calcule la distance du sommet $A$ au plan $\mathcal{P}$ du panneau solaire. Interprète le résultat.
\item \textbf{Sécurité du technicien.} Calcule la distance du technicien $M(6\,;3\,;2)$ au câble téléphonique $(AB)$. La norme de sécurité exige une distance d'au moins $3$ mètres : le technicien est-il en sécurité ?
\item \textbf{Rédaction.} Présente ta restitution sous la forme d'un rapport technique : objet de la mission, méthodes utilisées, résultats, conclusions et recommandations.
\end{enumerate}

\courssub{Ressources à mobiliser}

\begin{aretenir}[Boîte à outils de la leçon]
\begin{itemize}
\item \textbf{Représentation paramétrique d'une droite} passant par $P_0(x_0\,;y_0\,;z_0)$ et de vecteur directeur $\vec{u}(a\,;b\,;c)$ :
\[
\begin{cases} x = x_0 + at \\ y = y_0 + bt \\ z = z_0 + ct \end{cases} \quad (t \in \mathbb{R}).
\]
\item \textbf{Positions de deux droites} de vecteurs directeurs $\vec{u}$ et $\vec{v}$ : si $\vec{u}$ et $\vec{v}$ sont colinéaires, les droites sont parallèles (confondues ou strictement parallèles) ; sinon, on cherche un point d'intersection en résolvant le système obtenu en égalisant les deux représentations paramétriques : s'il y a une solution, les droites sont sécantes ; sinon, elles sont non coplanaires.
\item \textbf{Équation cartésienne d'un plan} de vecteur normal $\vec{n}(a\,;b\,;c)$ : $ax + by + cz + d = 0$, où $d$ s'obtient en écrivant qu'un point connu du plan vérifie l'équation.
\item \textbf{Distance d'un point à un plan} : si $\mathcal{P} : ax + by + cz + d = 0$ et $M_0(x_0\,;y_0\,;z_0)$, alors
\[
d(M_0,\mathcal{P}) = \frac{|ax_0 + by_0 + cz_0 + d|}{\sqrt{a^2 + b^2 + c^2}}.
\]
\item \textbf{Distance d'un point à une droite} : $d(M_0, \Delta) = M_0H$, où $H$ est le projeté orthogonal de $M_0$ sur $\Delta$, obtenu en résolvant $\overrightarrow{P_0H} = t\,\vec{u}$ avec $\overrightarrow{M_0H} \cdot \vec{u} = 0$.
\end{itemize}
\end{aretenir}

\courssub{Démarche et corrigé commenté}

\begin{exoresolu}[Rapport technique — Mission de Garoua]
Reprendre l'ensemble des consignes 1 à 6 de la mission et y répondre de manière rigoureuse.
\tcblower
\textbf{1. Repérage.} Tous les points ont une ordonnée $y$ positive ou nulle : le site s'étend vers le nord. Les points $A$ et $B$ ont la même cote $z = 8$ et la même ordonnée $y = 0$ : le câble téléphonique est horizontal et orienté est-ouest. Les points $E$, $F$, $G$ ont tous la même ordonnée $y = 10$ : le panneau solaire est contenu dans un plan vertical parallèle au plan $(xOz)$. La figure en perspective cavalière (ci-dessous) confirme visuellement ces observations.

\textbf{2. Représentations paramétriques des câbles.}

Pour le câble téléphonique : $\overrightarrow{AB}(12-0\,;0-0\,;8-8) = (12\,;0\,;0)$. On peut simplifier et choisir le vecteur directeur $\vec{u}(1\,;0\,;0)$. Une représentation paramétrique de $(AB)$ est :
\[
\begin{cases} x = t \\ y = 0 \\ z = 8 \end{cases} \quad (t \in \mathbb{R}).
\]

Pour la ligne électrique : $\overrightarrow{CD}(14-2\,;5-5\,;6-10) = (12\,;0\,;-4)$, que l'on simplifie en $\vec{v}(3\,;0\,;-1)$. Une représentation paramétrique de $(CD)$ est :
\[
\begin{cases} x = 2 + 3s \\ y = 5 \\ z = 10 - s \end{cases} \quad (s \in \mathbb{R}).
\]

\textbf{3. Position relative des deux câbles.}

Les vecteurs $\vec{u}(1\,;0\,;0)$ et $\vec{v}(3\,;0\,;-1)$ ne sont pas colinéaires : il n'existe pas de réel $k$ tel que $\vec{v} = k\vec{u}$, car la troisième coordonnée donnerait $-1 = 0$. Les droites ne sont donc pas parallèles.

Cherchons un éventuel point d'intersection en égalisant les coordonnées :
\[
\begin{cases} t = 2 + 3s \\ 0 = 5 \\ 8 = 10 - s \end{cases}
\]
La deuxième équation $0 = 5$ est impossible : le système n'a aucune solution. Les droites $(AB)$ et $(CD)$ n'ont pas de point commun et ne sont pas parallèles : elles sont \textbf{non coplanaires}.

\emph{Conclusion de sécurité :} les deux câbles ne se croisent pas. Le câble téléphonique reste dans le plan $y = 0$ tandis que la ligne électrique reste dans le plan $y = 5$ : elles évoluent dans deux plans horizontaux distincts distants de $5$ mètres. L'installation est conforme sur ce point.

\textbf{4. Équation cartésienne du plan du panneau.}

$\overrightarrow{EF}(6\,;0\,;0)$ et $\overrightarrow{EG}(0\,;0\,;4)$ ne sont pas colinéaires, donc $E$, $F$, $G$ définissent bien un plan.

Un vecteur normal $\vec{n}$ est orthogonal à la fois à $\overrightarrow{EF}$ et $\overrightarrow{EG}$. Cherchons $\vec{n}(a\,;b\,;c)$ :
\[
\vec{n} \cdot \overrightarrow{EF} = 6a = 0 \implies a = 0 ; \qquad \vec{n} \cdot \overrightarrow{EG} = 4c = 0 \implies c = 0.
\]
On peut donc prendre $\vec{n}(0\,;1\,;0)$. Une équation de $\mathcal{P}$ est de la forme $y + d = 0$. Comme $E(0\,;10\,;0) \in \mathcal{P}$, on a $10 + d = 0$, soit $d = -10$. Finalement :
\[
\mathcal{P} : y - 10 = 0.
\]
Ce résultat était prévisible : les trois points ont tous pour ordonnée $y = 10$, le plan du panneau est le plan vertical d'équation $y = 10$.

\textbf{5. Distance du sommet $A$ au plan du panneau.}
\[
d(A,\mathcal{P}) = \frac{|y_A - 10|}{\sqrt{0^2 + 1^2 + 0^2}} = \frac{|0 - 10|}{1} = 10 \text{ m}.
\]
Le sommet du pylône se trouve à $10$ mètres du plan du panneau solaire.

\textbf{6. Distance du technicien au câble $(AB)$.}

Soit $H$ le projeté orthogonal de $M(6\,;3\,;2)$ sur $(AB)$. Comme $H \in (AB)$, il existe $t$ tel que $H(t\,;0\,;8)$. Alors $\overrightarrow{MH}(t-6\,;-3\,;6)$. La condition d'orthogonalité $\overrightarrow{MH} \cdot \vec{u} = 0$ donne :
\[
(t-6) \times 1 + (-3) \times 0 + 6 \times 0 = 0 \implies t = 6.
\]
Donc $H(6\,;0\,;8)$ et :
\[
d(M,(AB)) = MH = \sqrt{(6-6)^2 + (0-3)^2 + (8-2)^2} = \sqrt{0 + 9 + 36} = \sqrt{45} = 3\sqrt{5} \approx 6{,}7 \text{ m}.
\]
Comme $3\sqrt{5} \approx 6{,}7 > 3$, le technicien respecte largement la distance de sécurité de $3$ mètres par rapport au câble téléphonique.

\textbf{Recommandations finales du rapport :} les câbles ne se croisent pas (droites non coplanaires, séparées de $5$ m dans la direction nord-sud), le sommet du pylône est à $10$ m du panneau, et le poste de travail du technicien est conforme. Le site peut être mis en service.
\end{exoresolu}

\begin{popfigure}
\begin{tikzpicture}[
    x={(0.3cm,0cm)}, 
    y={(-0.15cm,-0.15cm)}, 
    z={(0cm,0.3cm)},
    scale=1,
    >=Stealth,
    point/.style={circle,fill,inner sep=1.5pt},
    axis/.style={thin,gray!60},
    cable/.style={thick,popblue},
    line/.style={thick,poporange},
    panel/.style={fill=popgreen!20,draw=popgreen,opacity=0.6},
    label/.style={font=\scriptsize,inner sep=1pt}
]
% Axes
\draw[axis,->] (0,0,0) -- (15,0,0) node[right,label] {$x$ (Est)};
\draw[axis,->] (0,0,0) -- (0,12,0) node[left,label] {$y$ (Nord)};
\draw[axis,->] (0,0,0) -- (0,0,12) node[above,label] {$z$ (Haut)};
\draw[axis] (0,0,0) -- (-1,0,0);
\draw[axis] (0,0,0) -- (0,-1,0);
\draw[axis] (0,0,0) -- (0,0,-1);

% Points
\coordinate[point,label=below left:$A$] (A) at (0,0,8);
\coordinate[point,label=below right:$B$] (B) at (12,0,8);
\coordinate[point,label=above left:$C$] (C) at (2,5,10);
\coordinate[point,label=above right:$D$] (D) at (14,5,6);
\coordinate[point,label=below:$E$] (E) at (0,10,0);
\coordinate[point,label=below:$F$] (F) at (6,10,0);
\coordinate[point,label=above:$G$] (G) at (0,10,4);
\coordinate[point,label=right:$M$] (M) at (6,3,2);

% Câble téléphonique (AB)
\draw[cable] (A) -- (B) node[midway,above,label,popblue] {Câble téléphonique};

% Ligne électrique (CD)
\draw[line] (C) -- (D) node[midway,above,label,poporange] {Ligne électrique};

% Panneau solaire (Plan P : rectangle E-F-G)
\draw[panel] (E) -- (F) -- (G) -- cycle;
\draw[panel,dashed] (F) -- (0,10,4) -- (G); % Compléter le rectangle visuellement si besoin, mais E,F,G suffisent
\node[label,popgreen] at (3,10,2) {Panneau $\mathcal{P}$};

% Projeté H pour M sur AB
\coordinate[point,label=above:$H$] (H) at (6,0,8);
\draw[dashed,thin,poppink] (M) -- (H) node[midway,right,label,poppink] {$MH$};

% Légendes points
\foreach \p/\pos in {A/below left,B/below right,C/above left,D/above right,E/below,F/below,G/above,M/right} {
    \node[point] at (\p) {};
    % Labels déjà mis dans coordinate
}
\end{tikzpicture}
\legende{Modélisation du site de Garoua en perspective cavalière (échelle 0,3 cm/m). Le câble téléphonique (bleu) est horizontal à $y=0$, la ligne électrique (orange) à $y=5$, le panneau solaire (vert) dans le plan $y=10$.}
\end{popfigure}

\courssub{Bilan de l'activité}

Cette activité a permis de mettre en évidence plusieurs idées essentielles de la leçon :

\begin{itemize}
\item \textbf{Le repérage est la clé de toute la modélisation} : une fois le repère orthonormal choisi, des objets concrets (câbles, panneau, technicien) deviennent des objets mathématiques (droites, plan, point) sur lesquels on peut calculer.
\item \textbf{Deux droites de l'espace ne se croisent pas forcément}, même si leurs supports semblent se superposer sur un dessin : contrairement au plan, l'espace offre la possibilité de droites \cle{non coplanaires}. Ici, la lecture des équations ($y = 0$ pour l'un, $y = 5$ pour l'autre) montre immédiatement l'impossibilité d'une intersection.
\item \textbf{Un vecteur normal détermine entièrement l'équation d'un plan} (à condition de connaître un point) : c'est l'outil le plus efficace pour passer de la géométrie au calcul.
\item \textbf{Les formules de distance} (point-plan, point-droite via le projeté orthogonal) donnent des réponses concrètes à des questions de sécurité : on ne se contente plus de dire « c'est loin », on \emph{prouve} que la distance vaut $10$ m ou $3\sqrt{5}$ m.
\item Enfin, la rédaction d'un rapport technique montre que les mathématiques sont un langage de communication professionnel : un résultat n'a de valeur que s'il est justifié, interprété et transmis clairement.
\end{itemize}

\begin{attention}[Ce que le correcteur attend dans ton rapport]
Toute affirmation de position relative doit être \emph{justifiée par le calcul} (colinéarité des vecteurs directeurs, résolution du système), jamais par la seule lecture de la figure. De même, une distance doit être accompagnée de sa formule, de l'application numérique et d'une unité. Une figure soignée est un plus, mais elle ne remplace pas une démonstration.
\end{attention}

\newpage

\coursec{Méthodes et exercices résolus}

\coursec{Géométrie analytique dans l'espace}

\courssub{Repères et coordonnées dans l'espace}

\begin{savaistu}[Descartes et la naissance de la géométrie analytique]
René Descartes (1596--1650) a unifié la géométrie et l'algèbre en introduisant les coordonnées. La légende raconte qu'alité, il observait une mouche voler dans sa chambre et a eu l'idée de repérer sa position par trois distances aux murs perpendiculaires. C'est l'acte de naissance du repère orthonormé de l'espace, outil indispensable aujourd'hui pour la navigation GPS (satellites MTN/Orange, localisation précise à Yaoundé ou Douala), la modélisation 3D en architecture ou la robotique.
\end{savaistu}

\begin{experience}[Construire un repère de l'espace avec un cube]
Matériel : un cube en fil de fer (ou une boîte rectangulaire), trois fils de couleur rouge, vert, bleu.
\begin{enumerate}
\item Choisis un sommet $A$ comme origine. Place les trois fils le long des trois arêtes issues de $A$ : $\vec{i}$ (rouge), $\vec{j}$ (vert), $\vec{k}$ (bleu).
\item Vérifie que tu ne peux pas placer le fil bleu dans le plan formé par le rouge et le vert : les trois vecteurs sont \cle{non coplanaires}.
\item Repère le sommet opposé $G$. Pour aller de $A$ à $G$, tu as suivi $\vec{i}$, puis $\vec{j}$, puis $\vec{k}$. Les coordonnées de $G$ sont donc $(1;1;1)$.
\item Place le point $I$ milieu de $[FG]$. Décompose $\overrightarrow{AI}$ : $\overrightarrow{AB}+\overrightarrow{BF}+\overrightarrow{FI} = \vec{i}+\vec{k}+\frac{1}{2}\vec{j}$. Coordonnées : $(1;\frac{1}{2};1)$.
\end{enumerate}
Cette manipulation montre qu'un repère permet de coder tout point ou vecteur par un triplet de nombres.
\end{experience}

\begin{popfigure}
\begin{tikzpicture}[scale=1.1, >=Stealth, thick]
% Cavalier perspective: x right, y 45deg up-left (factor 0.7), z vertical up
\coordinate (A) at (0,0);
\coordinate (B) at (3,0);
\coordinate (D) at (1.05,1.05);
\coordinate (E) at (0,3);
\coordinate (C) at ($(B)+(D)$);
\coordinate (F) at ($(B)+(E)$);
\coordinate (H) at ($(D)+(E)$);
\coordinate (G) at ($(C)+(E)$);
% Hidden edges dashed
\draw[dashed, gray] (C) -- (G) (F) -- (G) (G) -- (H) (B) -- (C) (D) -- (C);
% Visible edges
\draw (A) -- (B) node[midway, below] {$\vec{i}$};
\draw (A) -- (D) node[midway, left] {$\vec{j}$};
\draw (A) -- (E) node[midway, left] {$\vec{k}$};
\draw (E) -- (F) -- (B);
\draw (E) -- (H) -- (D);
\draw (F) -- (G) (H) -- (G);
% Vertices
\foreach \p/\pos in {A/below left, B/below right, C/right, D/left, E/above left, F/above right, G/above right, H/above left}
  \node[circle, fill, inner sep=1.5pt, label=\pos:$\p$] at (\p) {};
% Point I midpoint FG
\coordinate (I) at ($(F)!0.5!(G)$);
\node[circle, fill, inner sep=1.5pt, label=above:$I$] at (I) {};
\draw[poporange, thick, ->] (B) -- (I) node[midway, right] {$\overrightarrow{BI}$};
\draw[popblue, thick, ->] (A) -- (G) node[midway, above left] {$\overrightarrow{AG}$};
\end{tikzpicture}
\legende{Cube $ABCDEFGH$ muni du repère $(A;\vec{i},\vec{j},\vec{k})$ avec $\vec{i}=\overrightarrow{AB}$, $\vec{j}=\overrightarrow{AD}$, $\vec{k}=\overrightarrow{AE}$. $I$ milieu de $[FG]$.}
\end{popfigure}

\begin{methode}[Coordonnées d'un vecteur dans une base]
Pour déterminer les coordonnées d'un vecteur $\vec{u}$ dans une base $(\vec{i},\vec{j},\vec{k})$ de l'espace :
\begin{enumerate}
\item \textbf{Identifier} les vecteurs de la base : ils sont souvent donnés par les arêtes d'un cube, d'un parallélépipède ou d'un tétraèdre (par exemple $\vec{i}=\overrightarrow{AB}$, $\vec{j}=\overrightarrow{AD}$, $\vec{k}=\overrightarrow{AE}$).
\item \textbf{Décomposer} $\vec{u}$ à l'aide de la relation de Chasles en faisant apparaître uniquement $\vec{i}$, $\vec{j}$, $\vec{k}$ : écrire $\vec{u}=x\vec{i}+y\vec{j}+z\vec{k}$.
\item \textbf{Conclure} : $\vec{u}\begin{pmatrix} x \\ y \\ z \end{pmatrix}$ dans la base $(\vec{i},\vec{j},\vec{k})$.
\end{enumerate}
\textbf{Réflexes :} dans un cube $ABCDEFGH$, $\overrightarrow{AC}=\vec{i}+\vec{j}$ (diagonale d'une face), $\overrightarrow{AG}=\vec{i}+\vec{j}+\vec{k}$ (grande diagonale) ; le milieu $I$ de $[AB]$ donne $\overrightarrow{AI}=\dfrac{1}{2}\vec{i}$.
\end{methode}

\begin{exoresolu}[Coordonnées dans un cube]
Soit $ABCDEFGH$ un cube. On munit l'espace de la base $(\vec{i},\vec{j},\vec{k})$ avec $\vec{i}=\overrightarrow{AB}$, $\vec{j}=\overrightarrow{AD}$, $\vec{k}=\overrightarrow{AE}$. Soit $I$ le milieu de $[FG]$. Déterminer les coordonnées des vecteurs $\overrightarrow{AG}$ et $\overrightarrow{BI}$ dans cette base.
\tcblower
\textbf{Coordonnées de $\overrightarrow{AG}$.} Par la relation de Chasles, en passant par les sommets du cube :
\[ \overrightarrow{AG}=\overrightarrow{AB}+\overrightarrow{BC}+\overrightarrow{CG}. \]
Or $ABCD$ est un carré, donc $\overrightarrow{BC}=\overrightarrow{AD}=\vec{j}$, et $BCGF$ est un carré, donc $\overrightarrow{CG}=\overrightarrow{AE}=\vec{k}$. Ainsi :
\[ \overrightarrow{AG}=\vec{i}+\vec{j}+\vec{k}, \qquad \overrightarrow{AG}\begin{pmatrix} 1 \\ 1 \\ 1 \end{pmatrix}. \]

\textbf{Coordonnées de $\overrightarrow{BI}$.} On décompose en passant par des points connus :
\[ \overrightarrow{BI}=\overrightarrow{BF}+\overrightarrow{FI}. \]
D'une part $\overrightarrow{BF}=\overrightarrow{AE}=\vec{k}$ (arêtes parallèles du cube). D'autre part $I$ est le milieu de $[FG]$ et $\overrightarrow{FG}=\overrightarrow{AD}=\vec{j}$, donc $\overrightarrow{FI}=\dfrac{1}{2}\overrightarrow{FG}=\dfrac{1}{2}\vec{j}$. Par conséquent :
\[ \overrightarrow{BI}=\vec{k}+\frac{1}{2}\vec{j}=0\vec{i}+\frac{1}{2}\vec{j}+1\vec{k}, \qquad \overrightarrow{BI}\begin{pmatrix} 0 \\ \frac{1}{2} \\ 1 \end{pmatrix}. \]

\textbf{Conclusion :} $\overrightarrow{AG}\begin{pmatrix} 1 \\ 1 \\ 1 \end{pmatrix}$ et $\overrightarrow{BI}\begin{pmatrix} 0 \\ \frac{1}{2} \\ 1 \end{pmatrix}$ dans la base $(\vec{i},\vec{j},\vec{k})$.
\end{exoresolu}

\begin{aretenir}[Coordonnées des sommets d'un cube dans le repère $(A;\overrightarrow{AB},\overrightarrow{AD},\overrightarrow{AE})$]
Si le cube a pour arête la longueur unité :
\[
\begin{array}{c|cccccccc}
\text{Sommet} & A & B & C & D & E & F & G & H \\ \hline
\text{Coordonnées} & (0;0;0) & (1;0;0) & (1;1;0) & (0;1;0) & (0;0;1) & (1;0;1) & (1;1;1) & (0;1;1)
\end{array}
\]
Le centre $\Omega$ (milieu de $[AG]$) a pour coordonnées $\left(\dfrac{1}{2};\dfrac{1}{2};\dfrac{1}{2}\right)$.
\end{aretenir}

\begin{methode}[Repère de l'espace, coordonnées d'un point]
\begin{enumerate}
\item \textbf{Reconnaître un repère} $(O;\vec{i},\vec{j},\vec{k})$ : vérifier que les trois vecteurs sont \cle{non coplanaires}, c'est-à-dire qu'aucun n'est combinaison linéaire des deux autres. Dans un cube $ABCDEFGH$, $(A;\overrightarrow{AB},\overrightarrow{AD},\overrightarrow{AE})$ est un repère car les trois arêtes issues de $A$ ne sont pas dans un même plan.
\item \textbf{Lire les coordonnées} d'un point $M$ : décomposer $\overrightarrow{OM}=x\vec{i}+y\vec{j}+z\vec{k}$ ; alors $M(x;y;z)$.
\item \textbf{Placer un point} $M(x;y;z)$ : à partir de $O$, effectuer la translation de vecteur $x\vec{i}$, puis $y\vec{j}$, puis $z\vec{k}$ (l'ordre est indifférent).
\end{enumerate}
\textbf{Réflexe :} dans le repère $(A;\overrightarrow{AB},\overrightarrow{AD},\overrightarrow{AE})$ du cube, $A(0;0;0)$, $B(1;0;0)$, $D(0;1;0)$, $E(0;0;1)$, $C(1;1;0)$, $F(1;0;1)$, $H(0;1;1)$, $G(1;1;1)$.
\end{methode}

\begin{exoresolu}[Repère et coordonnées dans un cube]
Soit $ABCDEFGH$ un cube de centre $\Omega$. On considère le repère $(A;\overrightarrow{AB},\overrightarrow{AD},\overrightarrow{AE})$.
\begin{enumerate}
\item Justifier que $(A;\overrightarrow{AB},\overrightarrow{AD},\overrightarrow{AE})$ est un repère de l'espace.
\item Donner les coordonnées des points $C$, $F$ et $\Omega$.
\end{enumerate}
\tcblower
\textbf{1.} Les vecteurs $\overrightarrow{AB}$ et $\overrightarrow{AD}$ dirigent deux arêtes de la face $ABCD$, donc ils engendrent le plan $(ABC)$. Le vecteur $\overrightarrow{AE}$ n'est pas parallèle à ce plan (l'arête $[AE]$ est perpendiculaire à la face $ABCD$). Les trois vecteurs sont donc non coplanaires : $(A;\overrightarrow{AB},\overrightarrow{AD},\overrightarrow{AE})$ est bien un \cle{repère de l'espace}.

\textbf{2.} On décompose chaque vecteur position :
\begin{itemize}
\item $\overrightarrow{AC}=\overrightarrow{AB}+\overrightarrow{AD}$ (diagonale du carré $ABCD$), donc $C(1;1;0)$.
\item $\overrightarrow{AF}=\overrightarrow{AB}+\overrightarrow{AE}$ (diagonale du carré $ABFE$), donc $F(1;0;1)$.
\item $\Omega$ est le milieu de la grande diagonale $[AG]$, donc $\overrightarrow{A\Omega}=\dfrac{1}{2}\overrightarrow{AG}=\dfrac{1}{2}\overrightarrow{AB}+\dfrac{1}{2}\overrightarrow{AD}+\dfrac{1}{2}\overrightarrow{AE}$, d'où $\Omega\left(\dfrac{1}{2};\dfrac{1}{2};\dfrac{1}{2}\right)$.
\end{itemize}
\textbf{Conclusion :} $C(1;1;0)$, $F(1;0;1)$ et $\Omega\left(\frac{1}{2};\frac{1}{2};\frac{1}{2}\right)$.
\end{exoresolu}

\courssub{Représentations d'une droite de l'espace}

\begin{experience}[Découvrir la représentation paramétrique d'une droite]
Dans un repère, on donne le point $A(1;2;-1)$ et le vecteur $\vec{u}\begin{pmatrix} 2 \\ -1 \\ 3 \end{pmatrix}$.
\begin{enumerate}
\item Place $A$. Place $B$ tel que $\overrightarrow{AB}=\vec{u}$, puis $C$ tel que $\overrightarrow{AC}=2\vec{u}$, $D$ tel que $\overrightarrow{AD}=-\vec{u}$.
\item Que remarques-tu sur les points $A,B,C,D$ ? Ils sont alignés.
\item Si $M$ est un point quelconque de cette droite, il existe un réel $t$ tel que $\overrightarrow{AM}=t\vec{u}$.
\item En coordonnées : $\begin{cases} x_M = 1+2t \\ y_M = 2-t \\ z_M = -1+3t \end{cases}$. C'est la \cle{représentation paramétrique} de la droite.
\end{enumerate}
\end{experience}

\begin{methode}[Représentation paramétrique d'une droite]
Un système de la forme
\[ \begin{cases} x = x_A + \alpha t \\ y = y_A + \beta t \\ z = z_A + \gamma t \end{cases} \quad (t\in\mathbb{R}) \]
est une représentation paramétrique d'une droite \textbf{si et seulement si} $(\alpha,\beta,\gamma)\neq(0,0,0)$.
\begin{enumerate}
\item \textbf{Lire le point} : $A(x_A;y_A;z_A)$ (les constantes).
\item \textbf{Lire le vecteur directeur} : $\vec{u}\begin{pmatrix} \alpha \\ \beta \\ \gamma \end{pmatrix}$ (les coefficients du paramètre $t$).
\item \textbf{Pour trouver une représentation} d'une droite $(AB)$ : prendre $A$ comme point de passage et $\overrightarrow{AB}$ comme vecteur directeur, puis écrire le système.
\item \textbf{Pour tester si un point $M$ appartient à la droite} : vérifier qu'il existe un \emph{même} réel $t$ satisfaisant les trois équations.
\end{enumerate}
\end{methode}

\begin{popfigure}
\begin{tikzpicture}[scale=0.9, >=Stealth]
% Axes
\draw[->] (-0.5,0) -- (4,0) node[right] {$x$};
\draw[->] (0,-0.5) -- (0,3.5) node[above] {$y$};
\draw[->] (0,0) -- (-1.5,-1.5) node[below left] {$z$};
% Line d
\draw[popblue, thick] (-1,2.5) -- (3.5,-0.5);
\draw[popblue, thick, ->] (1,1) -- (1.5,0.75) node[above right] {$\vec{u}$};
\node[circle, fill, inner sep=1.5pt, label=above left:$A$] at (1,1) {};
\node[circle, fill, inner sep=1.5pt, label=right:$M$] at (2,0.5) {};
\draw[dashed] (1,1) -- (2,0.5) node[midway, right] {$t\vec{u}$};
% Parametric text
\node[popblue, align=left] at (4,2.5) {$(d):\begin{cases} x=x_A+\alpha t\\ y=y_A+\beta t\\ z=z_A+\gamma t \end{cases}$};
\end{tikzpicture}
\legende{Représentation paramétrique d'une droite : point $A$ et vecteur directeur $\vec{u}$.}
\end{popfigure}

\begin{exoresolu}[Droite définie par deux points]
L'espace est muni d'un repère $(O;\vec{i},\vec{j},\vec{k})$. On donne $A(1;-2;3)$ et $B(3;0;-1)$.
\begin{enumerate}
\item Déterminer une représentation paramétrique de la droite $(AB)$.
\item Le point $C(5;2;-5)$ appartient-il à $(AB)$ ? Et le point $D(0;-3;4)$ ?
\end{enumerate}
\tcblower
\textbf{1.} Un vecteur directeur de $(AB)$ est $\overrightarrow{AB}\begin{pmatrix} 3-1 \\ 0-(-2) \\ -1-3 \end{pmatrix}=\begin{pmatrix} 2 \\ 2 \\ -4 \end{pmatrix}$. On peut simplifier : $\vec{u}=\frac{1}{2}\overrightarrow{AB}\begin{pmatrix} 1 \\ 1 \\ -2 \end{pmatrix}$ dirige aussi $(AB)$. La droite passe par $A(1;-2;3)$, donc une représentation paramétrique est :
\[ (AB):\begin{cases} x = 1 + t \\ y = -2 + t \\ z = 3 - 2t \end{cases} \quad (t\in\mathbb{R}). \]

\textbf{2.} \emph{Point $C(5;2;-5)$ :} la première équation donne $t=5-1=4$. Vérifions : $y=-2+4=2$ \checkmark{} et $z=3-2\times 4=-5$ \checkmark. Le même $t=4$ convient aux trois équations, donc $C\in(AB)$.

\emph{Point $D(0;-3;4)$ :} la première équation donne $t=0-1=-1$. Alors $y=-2+(-1)=-3$ \checkmark{} mais $z=3-2\times(-1)=5\neq 4$. Il n'existe pas de réel $t$ unique, donc $D\notin(AB)$.

\textbf{Conclusion :} $C$ appartient à $(AB)$ (pour $t=4$), $D$ n'y appartient pas.
\end{exoresolu}

\begin{aretenir}[Représentation paramétrique d'une droite]
Une droite $d$ passant par $A(x_A;y_A;z_A)$ et de vecteur directeur $\vec{u}\begin{pmatrix} \alpha \\ \beta \\ \gamma \end{pmatrix}\neq\vec{0}$ admet pour représentation paramétrique :
\[ (d):\begin{cases} x = x_A + \alpha t \\ y = y_A + \beta t \\ z = z_A + \gamma t \end{cases} \quad (t\in\mathbb{R}). \]
\textbf{Appartenance :} $M(x_M;y_M;z_M)\in d \Leftrightarrow \exists t\in\mathbb{R},\ \begin{cases} x_M=x_A+\alpha t\\ y_M=y_A+\beta t\\ z_M=z_A+\gamma t \end{cases}$.
\end{aretenir}

\courssub{Positions relatives de deux droites}

\begin{popfigure}
\begin{tikzpicture}[scale=0.8, >=Stealth]
% Case 1: Parallel
\begin{scope}[xshift=0cm]
\draw[popgreen, thick] (0,0) -- (2,0.5);
\draw[popgreen, thick] (0,1.5) -- (2,2);
\draw[->, popgreen] (0.5,0.25) -- (1,0.375) node[above right] {$\vec{u}$};
\draw[->, popgreen] (0.5,1.75) -- (1,1.875) node[above right] {$\vec{u}$};
\node at (1, -0.5) {Parallèles};
\node at (1, -0.8) {($\vec{u}\parallel\vec{u}'$)};
\end{scope}
% Case 2: Intersecting
\begin{scope}[xshift=5cm]
\draw[poporange, thick] (0,0) -- (2,1.5);
\draw[poporange, thick] (0,1.5) -- (2,0);
\fill (1,0.75) circle (2pt) node[above right] {$I$};
\node at (1, -0.5) {Sécantes};
\node at (1, -0.8) {($\vec{u}\nparallel\vec{u}'$, intersection)};
\end{scope}
% Case 3: Skew (non coplanaires)
\begin{scope}[xshift=10cm]
% Line 1 horizontal
\draw[poppurple, thick] (0,0.5) -- (2,0.5);
\draw[->, poppurple] (0.5,0.5) -- (1,0.5) node[above] {$\vec{u}$};
% Line 2 "vertical" in perspective (going into depth)
\draw[poppurple, thick] (1,1.5) -- (1.5,0.2);
\draw[->, poppurple] (1.1,1.3) -- (1.2,1.0) node[right] {$\vec{u}'$};
% Indicate non-coplanarity
\draw[dashed, gray] (1,0.5) -- (1,1.5);
\node at (1, -0.5) {Non coplanaires};
\node at (1, -0.8) {($\vec{u}\nparallel\vec{u}'$, pas d'intersection)};
\end{scope}
\end{tikzpicture}
\legende{Les trois positions relatives de deux droites dans l'espace.}
\end{popfigure}

\begin{methode}[Position relative de deux droites]
Soient $d$ de vecteur directeur $\vec{u}$ passant par $A$, et $d'$ de vecteur directeur $\vec{u}'$ passant par $A'$.
\begin{enumerate}
\item \textbf{Tester la colinéarité} de $\vec{u}$ et $\vec{u}'$ (coordonnées proportionnelles).
\begin{itemize}
\item Si $\vec{u}$ et $\vec{u}'$ sont colinéaires : les droites sont \cle{parallèles}. Elles sont \textbf{confondues} si de plus $A\in d'$ (ou $A'\in d$), \textbf{strictement parallèles} sinon.
\end{itemize}
\item Si $\vec{u}$ et $\vec{u}'$ ne sont pas colinéaires : \textbf{résoudre le système} formé par les deux représentations paramétriques (avec deux paramètres distincts $t$ et $t'$).
\begin{itemize}
\item Si le système admet une solution : les droites sont \cle{sécantes} en un point que l'on calcule.
\item Si le système n'a pas de solution : les droites sont \cle{non coplanaires}.
\end{itemize}
\end{enumerate}
\textbf{Réflexe :} pour tester la colinéarité, comparer les rapports $\dfrac{\alpha'}{\alpha}$, $\dfrac{\beta'}{\beta}$, $\dfrac{\gamma'}{\gamma}$ (attention aux coordonnées nulles : un zéro au numérateur impose un zéro au dénominateur).
\end{methode}

\begin{exoresolu}[Positions relatives de droites]
Dans un repère de l'espace, on donne les droites :
\[ d:\begin{cases} x = 1 + t \\ y = 2 - t \\ z = 3 + 2t \end{cases}\ (t\in\mathbb{R}), \qquad d':\begin{cases} x = 2 + 2t' \\ y = 1 + t' \\ z = -1 - t' \end{cases}\ (t'\in\mathbb{R}), \qquad d'':\begin{cases} x = -3 - 2s \\ y = 4 + 2s \\ z = -1 - 4s \end{cases}\ (s\in\mathbb{R}). \]
Étudier la position relative de $d$ et $d''$, puis de $d$ et $d'$.
\tcblower
\textbf{Position de $d$ et $d''$.} Vecteurs directeurs : $\vec{u}\begin{pmatrix} 1 \\ -1 \\ 2 \end{pmatrix}$ pour $d$ et $\vec{w}\begin{pmatrix} -2 \\ 2 \\ -4 \end{pmatrix}$ pour $d''$. On constate que $\vec{w}=-2\vec{u}$ : les vecteurs sont colinéaires, donc $d$ et $d''$ sont parallèles. Sont-elles confondues ? Le point $A(1;2;3)$ de $d$ appartient-il à $d''$ ? Il faudrait $-3-2s=1$ soit $s=-2$ ; alors $y=4+2(-2)=0\neq 2$. Donc $A\notin d''$.

\textbf{Conclusion :} $d$ et $d''$ sont \textbf{strictement parallèles}.

\textbf{Position de $d$ et $d'$.} Vecteur directeur de $d'$ : $\vec{u}'\begin{pmatrix} 2 \\ 1 \\ -1 \end{pmatrix}$. Les rapports $\frac{2}{1}=2$, $\frac{1}{-1}=-1$ ne sont pas égaux : $\vec{u}$ et $\vec{u}'$ ne sont pas colinéaires. Résolvons le système d'intersection :
\[ \begin{cases} 1+t = 2+2t' \\ 2-t = 1+t' \\ 3+2t = -1-t' \end{cases} \]
Des deux premières équations : $t=1+2t'$ et $t=1-t'$, d'où $1+2t'=1-t'$, soit $3t'=0$, $t'=0$ puis $t=1$. Vérifions dans la troisième équation : $3+2\times 1=5$ et $-1-0=-1$ : $5\neq -1$. Le système est incompatible.

\textbf{Conclusion :} $d$ et $d'$ ne sont ni parallèles ni sécantes : elles sont \textbf{non coplanaires}.
\end{exoresolu}

\begin{aretenir}[Positions relatives de deux droites]
Soient $d$ (directeur $\vec{u}$, point $A$) et $d'$ (directeur $\vec{u}'$, point $A'$).
\begin{itemize}
\item $\vec{u}\parallel\vec{u}'$ et $A\in d'$ $\Rightarrow$ $d$ et $d'$ \textbf{confondues}.
\item $\vec{u}\parallel\vec{u}'$ et $A\notin d'$ $\Rightarrow$ $d$ et $d'$ \textbf{strictement parallèles} (coplanaires).
\item $\vec{u}\nparallel\vec{u}'$ et système compatible $\Rightarrow$ $d$ et $d'$ \textbf{sécantes} (coplanaires).
\item $\vec{u}\nparallel\vec{u}'$ et système incompatible $\Rightarrow$ $d$ et $d'$ \textbf{non coplanaires}.
\end{itemize}
\end{aretenir}

\courssub{Représentations paramétriques d'un plan}

\begin{methode}[Représentation paramétrique d'un plan]
Un système
\[ \begin{cases} x = x_A + \alpha t + \alpha' t' \\ y = y_A + \beta t + \beta' t' \\ z = z_A + \gamma t + \gamma' t' \end{cases} \quad (t,t'\in\mathbb{R}) \]
est une représentation paramétrique d'un plan \textbf{si et seulement si} les vecteurs $\vec{u}\begin{pmatrix} \alpha \\ \beta \\ \gamma \end{pmatrix}$ et $\vec{v}\begin{pmatrix} \alpha' \\ \beta' \\ \gamma' \end{pmatrix}$ sont \cle{non colinéaires}.
\begin{enumerate}
\item \textbf{Lire les éléments} : le plan passe par $A(x_A;y_A;z_A)$ et est dirigé par $\vec{u}$ et $\vec{v}$.
\item \textbf{Vérifier la non-colinéarité} : les coordonnées de $\vec{u}$ et $\vec{v}$ ne doivent pas être proportionnelles.
\item \textbf{Plan défini par trois points non alignés} $A$, $B$, $C$ : prendre $\vec{u}=\overrightarrow{AB}$ et $\vec{v}=\overrightarrow{AC}$ après avoir vérifié qu'ils ne sont pas colinéaires.
\end{enumerate}
\end{methode}

\begin{popfigure}
\begin{tikzpicture}[scale=1, >=Stealth]
% Plane as parallelogram
\coordinate (A) at (0,0);
\coordinate (B) at (3,0.5);
\coordinate (D) at (0.5,2.5);
\coordinate (C) at ($(B)+(D)-(A)$);
\draw[popblue, thick, fill=popblue!10] (A) -- (B) -- (C) -- (D) -- cycle;
% Vectors
\draw[->, thick, poporange] (A) -- (B) node[midway, below] {$\vec{u}$};
\draw[->, thick, popgreen] (A) -- (D) node[midway, left] {$\vec{v}$};
% Point M = A + t u + t' v
\coordinate (M) at ($(A)+0.6*(B)-0.6*(A)+0.4*(D)-0.4*(A)$);
\draw[dashed] (A) -- (M);
\draw[->, thick, poppink] (A) -- (M) node[midway, above right] {$\overrightarrow{AM}=t\vec{u}+t'\vec{v}$};
\node[circle, fill, inner sep=1.5pt, label=below left:$A$] at (A) {};
\node[circle, fill, inner sep=1.5pt, label=above right:$M$] at (M) {};
% Axes indication
\draw[->, gray] (A) -- (-0.5,-0.5) node[below left] {$x$};
\draw[->, gray] (A) -- (-0.5,0.5) node[above left] {$y$};
\draw[->, gray] (A) -- (0,-1) node[below] {$z$};
\end{tikzpicture}
\legende{Plan paramétrique : point $A$ et deux vecteurs directeurs non colinéaires $\vec{u}$, $\vec{v}$.}
\end{popfigure}

\begin{exoresolu}[Plan passant par trois points]
Dans un repère de l'espace, on donne $A(1;0;2)$, $B(2;1;0)$ et $C(0;3;1)$.
\begin{enumerate}
\item Montrer que $A$, $B$, $C$ définissent un plan.
\item Donner une représentation paramétrique du plan $(ABC)$.
\item Le point $M(3;-2;5)$ appartient-il à ce plan ?
\end{enumerate}
\tcblower
\textbf{1.} $\overrightarrow{AB}\begin{pmatrix} 1 \\ 1 \\ -2 \end{pmatrix}$ et $\overrightarrow{AC}\begin{pmatrix} -1 \\ 3 \\ -1 \end{pmatrix}$. Ces vecteurs sont-ils colinéaires ? Il faudrait un réel $k$ avec $-1=k\times 1$ soit $k=-1$, mais alors la deuxième coordonnée donnerait $3=-1\times 1=-1$, ce qui est faux. Donc $\overrightarrow{AB}$ et $\overrightarrow{AC}$ ne sont pas colinéaires : les points $A$, $B$, $C$ sont non alignés et définissent un unique plan.

\textbf{2.} Le plan $(ABC)$ passe par $A(1;0;2)$ et est dirigé par $\overrightarrow{AB}$ et $\overrightarrow{AC}$ :
\[ (ABC):\begin{cases} x = 1 + t - t' \\ y = t + 3t' \\ z = 2 - 2t - t' \end{cases} \quad (t,t'\in\mathbb{R}). \]

\textbf{3.} $M(3;-2;5)\in(ABC)$ s'il existe $t,t'$ vérifiant le système. La deuxième équation donne $t+3t'=-2$, et la première $t-t'=2$. Par soustraction : $4t'=-4$, soit $t'=-1$, puis $t=1$. Vérifions la troisième : $z=2-2(1)-(-1)=1\neq 5$. Le système est incompatible.

\textbf{Conclusion :} $M$ n'appartient pas au plan $(ABC)$.
\end{exoresolu}

\begin{aretenir}[Représentation paramétrique d'un plan]
Un plan $\mathcal{P}$ passant par $A(x_A;y_A;z_A)$ et dirigé par deux vecteurs non colinéaires $\vec{u}\begin{pmatrix} \alpha \\ \beta \\ \gamma \end{pmatrix}$ et $\vec{v}\begin{pmatrix} \alpha' \\ \beta' \\ \gamma' \end{pmatrix}$ admet pour représentation paramétrique :
\[ \mathcal{P}:\begin{cases} x = x_A + \alpha t + \alpha' t' \\ y = y_A + \beta t + \beta' t' \\ z = z_A + \gamma t + \gamma' t' \end{cases} \quad (t,t'\in\mathbb{R}). \]
\textbf{Appartenance :} $M\in\mathcal{P} \Leftrightarrow \exists (t,t')\in\mathbb{R}^2$ vérifiant le système.
\end{aretenir}

\courssub{Équation cartésienne d'un plan et vecteur normal}

\begin{savaistu}[Équation cartésienne et applications]
L'équation $ax+by+cz+d=0$ est l'analogue en 3D de l'équation $ax+by+c=0$ d'une droite dans le plan. Elle est cruciale en infographie (rendu 3D, détection de collisions), en télécommunications (modélisation des zones de couverture d'antennes relais MTN/Orange au Cameroun) et en architecture (calcul de surfaces de toits plans).
\end{savaistu}

\begin{methode}[Équation cartésienne d'un plan, vecteur normal]
\textbf{Formule clé :} le plan passant par $A(x_A;y_A;z_A)$ et de vecteur normal $\vec{n}\begin{pmatrix} a \\ b \\ c \end{pmatrix}$ (avec $\vec{n}\neq\vec{0}$) a pour équation :
\[ a(x-x_A)+b(y-y_A)+c(z-z_A)=0 \quad\Longleftrightarrow\quad ax+by+cz+d=0. \]
\begin{enumerate}
\item \textbf{Plan défini par un point et un vecteur normal} : appliquer directement la formule, puis développer et réduire.
\item \textbf{Plan défini par trois points} $A$, $B$, $C$ : écrire que $\vec{n}$ est normal au plan signifie $\vec{n}\cdot\overrightarrow{AB}=0$ et $\vec{n}\cdot\overrightarrow{AC}=0$ ; résoudre pour trouver $(a;b;c)$, puis utiliser le point $A$ pour trouver $d$.
\item \textbf{Lecture inverse} : dans $ax+by+cz+d=0$, un vecteur normal est $\vec{n}\begin{pmatrix} a \\ b \\ c \end{pmatrix}$.
\item \textbf{Appartenance} : $M\in\mathcal{P}\Leftrightarrow ax_M+by_M+cz_M+d=0$.
\end{enumerate}
\end{methode}

\begin{popfigure}
\begin{tikzpicture}[scale=1, >=Stealth]
% Plane
\coordinate (A) at (0,0);
\coordinate (B) at (3,0.3);
\coordinate (D) at (0.5,2.5);
\coordinate (C) at ($(B)+(D)-(A)$);
\draw[popblue, thick, fill=popblue!10] (A) -- (B) -- (C) -- (D) -- cycle;
% Normal vector
\draw[->, very thick, poporange] (1.5,1.2) -- (1.5,2.7) node[above] {$\vec{n}$};
\node[circle, fill, inner sep=1.5pt, label=below left:$A$] at (A) {};
% Point M outside
\coordinate (M) at (2,3.5);
\node[circle, fill, inner sep=1.5pt, label=above:$M$] at (M) {};
\draw[dashed] (M) -- (1.5,1.2) node[midway, right] {$\overrightarrow{AM}$};
% Right angle mark
\draw (1.3,1.4) -- (1.4,1.3) -- (1.5,1.4);
\end{tikzpicture}
\legende{Vecteur normal $\vec{n}$ orthogonal au plan : $\vec{n}\cdot\overrightarrow{AM}=0$ pour tout $M\in\mathcal{P}$.}
\end{popfigure}

\begin{exoresolu}[Équation cartésienne d'un plan]
Dans un repère orthonormé de l'espace, on donne $A(1;0;2)$, $B(2;1;0)$ et $C(0;3;1)$.
\begin{enumerate}
\item Déterminer un vecteur $\vec{n}$ normal au plan $(ABC)$.
\item En déduire une équation cartésienne du plan $(ABC)$.
\item Le point $D(4;-1;1)$ appartient-il à ce plan ?
\end{enumerate}
\tcblower
\textbf{1.} On a $\overrightarrow{AB}\begin{pmatrix} 1 \\ 1 \\ -2 \end{pmatrix}$ et $\overrightarrow{AC}\begin{pmatrix} -1 \\ 3 \\ -1 \end{pmatrix}$. Soit $\vec{n}\begin{pmatrix} a \\ b \\ c \end{pmatrix}$ un vecteur normal au plan. Alors :
\[ \vec{n}\cdot\overrightarrow{AB}=0 \Leftrightarrow a+b-2c=0 \qquad\text{et}\qquad \vec{n}\cdot\overrightarrow{AC}=0 \Leftrightarrow -a+3b-c=0. \]
Choisissons $c$ comme paramètre : de la première équation $a=2c-b$ ; en reportant dans la deuxième : $-(2c-b)+3b-c=0$, soit $4b-3c=0$, d'où $b=\dfrac{3c}{4}$. Prenons $c=4$ : alors $b=3$ et $a=2\times 4-3=5$. Un vecteur normal est :
\[ \vec{n}\begin{pmatrix} 5 \\ 3 \\ 4 \end{pmatrix}. \]
\emph{Vérification :} $5+3-8=0$ \checkmark{} et $-5+9-4=0$ \checkmark.

\textbf{2.} Le plan a une équation de la forme $5x+3y+4z+d=0$. Comme $A(1;0;2)\in(ABC)$ :
\[ 5(1)+3(0)+4(2)+d=0 \Leftrightarrow 13+d=0 \Leftrightarrow d=-13. \]
\[ \boxed{(ABC):\ 5x+3y+4z-13=0} \]
\emph{Vérification avec $B$ et $C$ :} $10+3+0-13=0$ \checkmark{} ; $0+9+4-13=0$ \checkmark.

\textbf{3.} Pour $D(4;-1;1)$ : $5(4)+3(-1)+4(1)-13=20-3+4-13=8\neq 0$.

\textbf{Conclusion :} $D$ n'appartient pas au plan $(ABC)$.
\end{exoresolu}

\begin{aretenir}[Équation cartésienne et vecteur normal]
\begin{itemize}
\item Plan $\mathcal{P}$ d'équation $ax+by+cz+d=0$ $\Rightarrow$ $\vec{n}\begin{pmatrix} a \\ b \\ c \end{pmatrix}$ est un vecteur normal à $\mathcal{P}$.
\item Plan passant par $A(x_A;y_A;z_A)$ de vecteur normal $\vec{n}(a;b;c)$ : $a(x-x_A)+b(y-y_A)+c(z-z_A)=0$.
\item Deux plans $\mathcal{P}_1: a_1x+b_1y+c_1z+d_1=0$ et $\mathcal{P}_2: a_2x+b_2y+c_2z+d_2=0$ :
  \begin{itemize}
  \item Parallèles $\Leftrightarrow \vec{n}_1\parallel\vec{n}_2$ (coefficients proportionnels).
  \item Perpendiculaires $\Leftrightarrow \vec{n}_1\cdot\vec{n}_2=0$ ($a_1a_2+b_1b_2+c_1c_2=0$).
  \item Sécants $\Leftrightarrow$ non parallèles (intersection = droite).
  \end{itemize}
\end{itemize}
\end{aretenir}

\courssub{Projetés orthogonaux et distances}

\begin{methode}[Projeté orthogonal et distance]
\textbf{Distance d'un point $M$ à un plan $\mathcal{P}:ax+by+cz+d=0$} (repère orthonormé) :
\begin{enumerate}
\item Écrire la droite $(d)$ passant par $M$ et de vecteur directeur $\vec{n}\begin{pmatrix} a \\ b \\ c \end{pmatrix}$ (normale au plan).
\item Le projeté orthogonal $H$ est l'intersection de $(d)$ et $\mathcal{P}$ : remplacer $x,y,z$ paramétrés dans l'équation du plan pour trouver $t$, puis $H$.
\item Calculer $MH$, ou appliquer directement la formule :
\[ d(M,\mathcal{P})=\dfrac{|ax_M+by_M+cz_M+d|}{\sqrt{a^2+b^2+c^2}}. \]
\end{enumerate}
\textbf{Distance d'un point $M$ à une droite $d$} passant par $A$, de vecteur directeur $\vec{u}$ :
\begin{enumerate}
\item Chercher $H\in d$ (paramètre $t$) tel que $\overrightarrow{MH}\cdot\vec{u}=0$.
\item Résoudre l'équation en $t$, en déduire $H$, puis calculer $d(M,d)=MH$.
\end{enumerate}
\end{methode}

\begin{popfigure}
\begin{tikzpicture}[scale=0.9, >=Stealth]
% Point to Plane
\begin{scope}
\coordinate (A) at (0,0);
\coordinate (B) at (3,0.3);
\coordinate (D) at (0.5,2.5);
\coordinate (C) at ($(B)+(D)-(A)$);
\draw[popblue, thick, fill=popblue!10] (A) -- (B) -- (C) -- (D) -- cycle;
\coordinate (M) at (1.5,3.5);
\coordinate (H) at (1.5,1.2);
\node[circle, fill, inner sep=1.5pt, label=above:$M$] at (M) {};
\node[circle, fill, inner sep=1.5pt, label=below:$H$] at (H) {};
\draw[->, thick, poporange] (M) -- (H) node[midway, right] {$\overrightarrow{MH}\parallel\vec{n}$};
\draw (1.3,1.4) -- (1.4,1.3) -- (1.5,1.4);
\node at (1.5, -0.5) {Point -- Plan};
\end{scope}
% Point to Line
\begin{scope}[xshift=6cm]
\draw[popgreen, thick] (0,0) -- (3,1);
\coordinate (M) at (1.5,2.5);
\coordinate (K) at (1.2,1.4);
\node[circle, fill, inner sep=1.5pt, label=above:$M$] at (M) {};
\node[circle, fill, inner sep=1.5pt, label=below:$K$] at (K) {};
\draw[->, thick, poporange] (M) -- (K) node[midway, right] {$\overrightarrow{MK}\perp\vec{u}$};
\draw (1.0,1.5) -- (1.1,1.3) -- (1.3,1.4);
\node at (1.5, -0.5) {Point -- Droite};
\end{scope}
\end{tikzpicture}
\legende{Projeté orthogonal : $H$ sur le plan (droite $\parallel\vec{n}$), $K$ sur la droite ($\overrightarrow{MK}\perp\vec{u}$).}
\end{popfigure}

\begin{exoresolu}[Distance d'un point à un plan et à une droite]
Dans un repère orthonormé, on donne le plan $\mathcal{P}:2x-y+2z-4=0$, le point $M(3;2;1)$ et la droite $d:\begin{cases} x=1+t \\ y=2t \\ z=-1+t \end{cases}\ (t\in\mathbb{R})$.
\begin{enumerate}
\item Déterminer le projeté orthogonal $H$ de $M$ sur $\mathcal{P}$ et la distance $d(M,\mathcal{P})$.
\item Déterminer le projeté orthogonal $K$ de $M$ sur $d$ et la distance $d(M,d)$.
\end{enumerate}
\tcblower
\textbf{1.} Un vecteur normal à $\mathcal{P}$ est $\vec{n}\begin{pmatrix} 2 \\ -1 \\ 2 \end{pmatrix}$. La perpendiculaire à $\mathcal{P}$ passant par $M$ a pour représentation :
\[ \begin{cases} x = 3+2t \\ y = 2-t \\ z = 1+2t \end{cases} \quad (t\in\mathbb{R}). \]
$H$ est l'intersection avec $\mathcal{P}$. On remplace dans l'équation du plan :
\[ 2(3+2t)-(2-t)+2(1+2t)-4=0 \Leftrightarrow 6+4t-2+t+2+4t-4=0 \Leftrightarrow 9t+2=0 \Leftrightarrow t=-\frac{2}{9}. \]
D'où $H\left(3-\frac{4}{9}\,;\,2+\frac{2}{9}\,;\,1-\frac{4}{9}\right)$, c'est-à-dire $H\left(\frac{23}{9};\frac{20}{9};\frac{5}{9}\right)$.

Distance : $d(M,\mathcal{P})=\dfrac{|2(3)-2+2(1)-4|}{\sqrt{4+1+4}}=\dfrac{|6-2+2-4|}{3}=\dfrac{2}{3}$.

\emph{Vérification :} $\overrightarrow{MH}\begin{pmatrix} -\frac{4}{9} \\ \frac{2}{9} \\ -\frac{4}{9} \end{pmatrix}$, $MH=\sqrt{\frac{16+4+16}{81}}=\sqrt{\frac{36}{81}}=\frac{6}{9}=\frac{2}{3}$ \checkmark.

\textbf{2.} Un point de $d$ s'écrit $K(1+t\,;\,2t\,;\,-1+t)$, et $\vec{u}\begin{pmatrix} 1 \\ 2 \\ 1 \end{pmatrix}$ dirige $d$. On a :
\[ \overrightarrow{MK}\begin{pmatrix} 1+t-3 \\ 2t-2 \\ -1+t-1 \end{pmatrix}=\begin{pmatrix} t-2 \\ 2t-2 \\ t-2 \end{pmatrix}. \]
$K$ est le projeté orthogonal de $M$ sur $d$ si et seulement si $\overrightarrow{MK}\cdot\vec{u}=0$ :
\[ (t-2)\times 1+(2t-2)\times 2+(t-2)\times 1=0 \Leftrightarrow t-2+4t-4+t-2=0 \Leftrightarrow 6t-8=0 \Leftrightarrow t=\frac{4}{3}. \]
Donc $K\left(\dfrac{7}{3};\dfrac{8}{3};\dfrac{1}{3}\right)$. Alors :
\[ \overrightarrow{MK}\begin{pmatrix} \frac{4}{3}-2 \\[2pt] \frac{8}{3}-2 \\[2pt] \frac{4}{3}-2 \end{pmatrix}=\begin{pmatrix} -\frac{2}{3} \\[2pt] \frac{2}{3} \\[2pt] -\frac{2}{3} \end{pmatrix}, \qquad d(M,d)=MK=\sqrt{\frac{4+4+4}{9}}=\frac{2\sqrt{3}}{3}. \]
\textbf{Conclusion :} $d(M,\mathcal{P})=\dfrac{2}{3}$ et $d(M,d)=\dfrac{2\sqrt{3}}{3}$.
\end{exoresolu}

\begin{aretenir}[Formules de distances (repère orthonormé)]
\begin{itemize}
\item \textbf{Point $M(x_M;y_M;z_M)$ au plan $\mathcal{P}:ax+by+cz+d=0$ :}
\[ d(M,\mathcal{P})=\frac{|ax_M+by_M+cz_M+d|}{\sqrt{a^2+b^2+c^2}}. \]
\item \textbf{Point $M$ à la droite $d$ (point $A$, directeur $\vec{u}$) :}
$K\in d$ tel que $\overrightarrow{MK}\cdot\vec{u}=0$, puis $d(M,d)=MK$.
\item \textbf{Droite $d$ (directeur $\vec{u}$) au plan $\mathcal{P}$ (normal $\vec{n}$) :} si $d\parallel\mathcal{P}$ ($\vec{u}\cdot\vec{n}=0$), distance = distance d'un point de $d$ à $\mathcal{P}$.
\end{itemize}
\end{aretenir}

\begin{attention}[Les 5 erreurs qui coûtent des points au Probatoire]
\begin{enumerate}
\item \textbf{Oublier de vérifier la non-colinéarité.} Avant d'affirmer que trois points définissent un plan, ou que deux vecteurs dirigent un plan, il faut \emph{démontrer} que les vecteurs ne sont pas colinéaires. Sans cette justification, la réponse est incomplète.
\item \textbf{Utiliser le même paramètre pour deux droites.} Pour étudier l'intersection de $d$ (paramètre $t$) et $d'$ (paramètre $t'$), il faut deux lettres différentes. Écrire $t=t'$ d'office revient à supposer le résultat.
\item \textbf{Conclure « non coplanaires » trop vite.} Deux droites sont non coplanaires seulement si elles ne sont \emph{ni parallèles ni sécantes} : il faut avoir testé la colinéarité des vecteurs directeurs \emph{et} résolu (ou montré l'incompatibilité du) système d'intersection.
\item \textbf{Se tromper de vecteur dans la formule de distance.} Dans $d(M,\mathcal{P})=\dfrac{|ax_M+by_M+cz_M+d|}{\sqrt{a^2+b^2+c^2}}$, le dénominateur utilise les coefficients $(a,b,c)$ du plan, pas les coordonnées de $M$. Ne pas oublier la valeur absolue au numérateur ni le $d$ dans l'équation.
\item \textbf{Confondre vecteur directeur et vecteur normal.} Pour une droite, on lit un \emph{vecteur directeur} ; pour un plan $ax+by+cz+d=0$, $(a,b,c)$ est un vecteur \emph{normal} (perpendiculaire), pas directeur. Pour le projeté sur une droite, la condition est $\overrightarrow{MK}\cdot\vec{u}=0$ avec $\vec{u}$ directeur de la droite ; pour le projeté sur un plan, on utilise $\vec{n}$ normal au plan.
\end{enumerate}
\end{attention}

\newpage

\coursec{Exercices}

\courssub{Je vérifie mes connaissances}

\begin{exercice}[QCM : une seule réponse exacte]
L'espace est muni d'un repère $(O\,;\vec{i},\vec{j},\vec{k})$. Pour chaque question, choisir la bonne réponse en justifiant brièvement.
\begin{enumerate}
\item Le point $A$ tel que $\overrightarrow{OA} = 2\vec{i} - 3\vec{j} + \vec{k}$ a pour coordonnées :
\begin{enumerate}
\item $(2\,;-3\,;1)$ \quad \item $(-2\,;3\,;-1)$ \quad \item $(2\,;3\,;1)$ \quad \item $(1\,;-3\,;2)$
\end{enumerate}
\item La représentation $\begin{cases} x = 1 + 2t \\ y = -t \\ z = 3 + t \end{cases}, t \in \mathbb{R}$, est celle d'une droite passant par $A(1\,;0\,;3)$ et dirigée par :
\begin{enumerate}
\item $\vec{u}(2\,;-1\,;1)$ \quad \item $\vec{u}(1\,;0\,;3)$ \quad \item $\vec{u}(2\,;1\,;1)$ \quad \item $\vec{u}(1\,;-1\,;3)$
\end{enumerate}
\item Un vecteur normal au plan $P : 3x - 2y + z - 5 = 0$ est :
\begin{enumerate}
\item $\vec{n}(3\,;-2\,;1)$ \quad \item $\vec{n}(3\,;2\,;1)$ \quad \item $\vec{n}(3\,;-2\,;-5)$ \quad \item $\vec{n}(1\,;-2\,;3)$
\end{enumerate}
\item La distance du point $O(0\,;0\,;0)$ au plan $P : x + 2y + 2z - 6 = 0$ vaut :
\begin{enumerate}
\item $6$ \quad \item $3$ \quad \item $2$ \quad \item $\dfrac{6}{5}$
\end{enumerate}
\end{enumerate}
\end{exercice}

\begin{exercice}[Vrai ou faux, en justifiant]
L'espace est muni d'un repère orthonormal $(O\,;\vec{i},\vec{j},\vec{k})$. Dire si chacune des affirmations suivantes est vraie ou fausse, en justifiant la réponse.
\begin{enumerate}
\item Trois points distincts de l'espace définissent toujours un unique plan.
\item Deux droites de l'espace qui n'ont aucun point commun sont parallèles.
\item Si deux plans ont des vecteurs normaux colinéaires, alors ces plans sont parallèles.
\item Les droites $(d) : \begin{cases} x = 2t \\ y = 1 - t \\ z = 3 \end{cases}$ et $(d') : \begin{cases} x = -4s \\ y = 3 + 2s \\ z = -1 \end{cases}$ sont orthogonales.
\end{enumerate}
\end{exercice}

\begin{exercice}[Lecture dans un repère]
$ABCDEFGH$ est un cube de côté $1$. On munit l'espace du repère $(A\,;\overrightarrow{AB},\overrightarrow{AD},\overrightarrow{AE})$.
\begin{enumerate}
\item Justifier que $(A\,;\overrightarrow{AB},\overrightarrow{AD},\overrightarrow{AE})$ est bien un repère de l'espace.
\item Donner les coordonnées des huit sommets du cube.
\item Donner les coordonnées des vecteurs $\overrightarrow{AG}$, $\overrightarrow{BH}$ et $\overrightarrow{EC}$.
\end{enumerate}
\end{exercice}

\begin{exercice}[Calculs directs]
L'espace est muni d'un repère $(O\,;\vec{i},\vec{j},\vec{k})$. On donne $A(2\,;-1\,;3)$, $B(0\,;4\,;1)$ et $\vec{u}(-1\,;2\,;5)$.
\begin{enumerate}
\item Calculer les coordonnées de $\overrightarrow{AB}$ et de $2\overrightarrow{AB} - 3\vec{u}$.
\item Les vecteurs $\overrightarrow{AB}$ et $\vec{u}$ sont-ils colinéaires ?
\item Déterminer une représentation paramétrique de la droite $(AB)$.
\item Donner une équation cartésienne du plan passant par $A$ et de vecteur normal $\vec{n}(1\,;-2\,;2)$.
\end{enumerate}
\end{exercice}

\courssub{Je m'entraîne}

\begin{exercice}[Positions relatives de deux droites]
L'espace est muni d'un repère $(O\,;\vec{i},\vec{j},\vec{k})$. Dans chaque cas, déterminer la position relative des droites $(d)$ et $(d')$ (sécantes, strictement parallèles, confondues ou non coplanaires) et, lorsqu'elles sont sécantes, les coordonnées de leur point d'intersection.
\begin{enumerate}
\item $(d) : \begin{cases} x = 1 + t \\ y = 2 - t \\ z = 3 + 2t \end{cases}, t \in \mathbb{R}$ \quad et \quad $(d') : \begin{cases} x = 4 - 2s \\ y = 1 + s \\ z = 7 - 4s \end{cases}, s \in \mathbb{R}$.
\item $(d) : \begin{cases} x = 2 + t \\ y = -1 + t \\ z = 1 - t \end{cases}, t \in \mathbb{R}$ \quad et \quad $(d') : \begin{cases} x = 5 - s \\ y = 2 + s \\ z = -2 + s \end{cases}, s \in \mathbb{R}$.
\item $(d) : \begin{cases} x = t \\ y = 1 + 2t \\ z = -t \end{cases}, t \in \mathbb{R}$ \quad et \quad $(d') : \begin{cases} x = 2 + s \\ y = s \\ z = 1 + s \end{cases}, s \in \mathbb{R}$.
\end{enumerate}
\end{exercice}

\begin{exercice}[Intersection de deux plans]
L'espace est muni d'un repère orthonormal $(O\,;\vec{i},\vec{j},\vec{k})$. On considère les plans
\[ P : 2x - y + z - 3 = 0 \qquad \text{et} \qquad P' : x + y - 2z + 1 = 0. \]
\begin{enumerate}
\item Donner un vecteur normal à chacun des plans $P$ et $P'$.
\item Montrer que $P$ et $P'$ sont sécants.
\item Déterminer une représentation paramétrique de leur droite d'intersection $(\Delta)$. (On pourra poser $z = t$ et exprimer $x$ et $y$ en fonction de $t$.)
\item Vérifier que le point $I(1\,;0\,;1)$ appartient à $(\Delta)$.
\end{enumerate}
\end{exercice}

\begin{exercice}[Distance d'un point à un plan et volume]
L'espace est muni d'un repère orthonormal $(O\,;\vec{i},\vec{j},\vec{k})$ où l'unité est le centimètre. On donne le point $M(2\,;-1\,;3)$ et le plan $P : x + 2y - 2z + 5 = 0$.
\begin{enumerate}
\item Calculer la distance du point $M$ au plan $P$.
\item Déterminer les coordonnées du projeté orthogonal $H$ de $M$ sur $P$.
\item Soit $T$ un triangle contenu dans $P$, d'aire \SI{6}{cm^2}. En déduire le volume du tétraèdre de sommet $M$ et de base $T$.
\end{enumerate}
\end{exercice}

\begin{exercice}[Projeté orthogonal sur une droite]
L'espace est muni d'un repère orthonormal $(O\,;\vec{i},\vec{j},\vec{k})$. On considère la droite $(d)$ passant par $A(1\,;1\,;0)$ et dirigée par $\vec{u}(2\,;-1\,;1)$, et le point $M(3\,;0\,;2)$.
\begin{enumerate}
\item Donner une représentation paramétrique de $(d)$.
\item Soit $H$ le point de $(d)$ de paramètre $t$. Exprimer les coordonnées de $\overrightarrow{MH}$ en fonction de $t$.
\item Déterminer $t$ pour que $\overrightarrow{MH} \cdot \vec{u} = 0$, puis en déduire les coordonnées du projeté orthogonal de $M$ sur $(d)$.
\item Calculer la distance de $M$ à la droite $(d)$.
\end{enumerate}
\end{exercice}

\courssub{Je me prépare au Probatoire}

\begin{exercice}[Étude d'un plan et d'une droite — type Probatoire]
L'espace est muni d'un repère orthonormal $(O\,;\vec{i},\vec{j},\vec{k})$. On considère les points
\[ A(1\,;2\,;-1), \qquad B(3\,;0\,;1), \qquad C(0\,;1\,;2), \qquad D(2\,;1\,;0) \]
et le vecteur $\vec{u}(1\,;1\,;1)$.
\begin{enumerate}
\item \begin{enumerate}
\item Calculer les coordonnées des vecteurs $\overrightarrow{AB}$ et $\overrightarrow{AC}$.
\item Montrer que les points $A$, $B$ et $C$ définissent un plan, noté $(ABC)$.
\end{enumerate}
\item \begin{enumerate}
\item Déterminer les coordonnées du produit vectoriel $\overrightarrow{AB} \wedge \overrightarrow{AC}$ (ou, à défaut, chercher un vecteur $\vec{n}(a\,;b\,;c)$ orthogonal à la fois à $\overrightarrow{AB}$ et à $\overrightarrow{AC}$).
\item En déduire une équation cartésienne du plan $(ABC)$.
\end{enumerate}
\item On note $(\Delta)$ la droite passant par $D$ et dirigée par $\vec{u}$.
\begin{enumerate}
\item Donner une représentation paramétrique de $(\Delta)$.
\item Montrer que $(\Delta)$ n'est pas parallèle au plan $(ABC)$.
\item Déterminer les coordonnées du point d'intersection $K$ de $(\Delta)$ et du plan $(ABC)$.
\end{enumerate}
\item Calculer la distance du point $D$ au plan $(ABC)$.
\end{enumerate}
\end{exercice}

\begin{exercice}[Le cube du menuisier — problème de synthèse]
Un menuisier de Bafoussam fabrique un coffre cubique $ABCDEFGH$ de côté $1$ mètre. Pour étudier la pose d'une tige de renfort, il munit l'espace du repère orthonormal $(A\,;\overrightarrow{AB},\overrightarrow{AD},\overrightarrow{AE})$.
\begin{enumerate}
\item Donner les coordonnées des sommets $B$, $D$, $E$ et $G$ dans ce repère.
\item \begin{enumerate}
\item Montrer que les vecteurs $\overrightarrow{BE}$ et $\overrightarrow{BG}$ ne sont pas colinéaires.
\item Déterminer une équation cartésienne du plan $(BEG)$.
\item Vérifier que le point $A$ n'appartient pas au plan $(BEG)$ et calculer la distance de $A$ à ce plan.
\end{enumerate}
\item La tige de renfort est modélisée par la droite $(AG)$.
\begin{enumerate}
\item Donner une représentation paramétrique de la droite $(AG)$.
\item Montrer que la droite $(AG)$ est perpendiculaire au plan $(BEG)$.
\item Déterminer les coordonnées du point d'intersection $\Omega$ de $(AG)$ et du plan $(BEG)$. Que représente $\Omega$ pour le triangle $BEG$ ?
\end{enumerate}
\item Les droites $(AG)$ et $(BE)$ sont-elles sécantes ? Justifier soigneusement.
\end{enumerate}
\end{exercice}

\begin{exercice}[Antenne relais et zone de couverture — type Probatoire]
Une entreprise de télécommunications installe une antenne relais sur une colline. Dans un repère orthonormal $(O\,;\vec{i},\vec{j},\vec{k})$ de l'espace (unité : \SI{100}{m}), le sol est modélisé par le plan $P : 2x + y - 2z + 4 = 0$. Le pied du pylône est le point $F(1\,;2\,;3)$ et le sommet du pylône est le point $S(3\,;3\,;0)$. Un technicien se trouve au point $T(5\,;-1\,;2)$.
\begin{enumerate}
\item \begin{enumerate}
\item Déterminer une représentation paramétrique de la droite $(FS)$ modélisant le pylône.
\item Le pylône est-il perpendiculaire au sol ? Justifier.
\end{enumerate}
\item \begin{enumerate}
\item Calculer la distance du technicien $T$ au plan $P$. Exprimer le résultat en mètres.
\item Déterminer les coordonnées du projeté orthogonal $H$ de $T$ sur le plan $P$.
\end{enumerate}
\item Un câble relie le technicien $T$ au sommet $S$ du pylône.
\begin{enumerate}
\item Déterminer une représentation paramétrique de la droite $(TS)$.
\item Montrer que la droite $(TS)$ est sécante au plan $P$ et déterminer les coordonnées du point d'intersection $I$.
\end{enumerate}
\item On considère le plan $Q : x - y + z - 1 = 0$.
\begin{enumerate}
\item Montrer que les plans $P$ et $Q$ sont sécants.
\item Déterminer une représentation paramétrique de leur droite d'intersection.
\end{enumerate}
\end{enumerate}
\end{exercice}

\newpage

\coursec{Corrigés des exercices}

\begin{corrige}[Exercice 1 — QCM : une seule réponse exacte]
\begin{enumerate}
\item \textbf{Réponse 1 : $(2\,;-3\,;1)$}.\\
Dans le repère $(O\,;\vec{i},\vec{j},\vec{k})$, les coordonnées d'un point $A$ sont les coefficients de $\overrightarrow{OA}$ sur la base $(\vec{i},\vec{j},\vec{k})$. Ici $\overrightarrow{OA} = 2\vec{i} - 3\vec{j} + \vec{k}$, donc $A(2\,;-3\,;1)$.

\item \textbf{Réponse 1 : $\vec{u}(2\,;-1\,;1)$}.\\
Une représentation paramétrique $\begin{cases} x = x_A + t\,u_x \\ y = y_A + t\,u_y \\ z = z_A + t\,u_z \end{cases}$ correspond à une droite passant par $A(x_A,y_A,z_A)$ et de vecteur directeur $\vec{u}(u_x,u_y,u_z)$. Ici $A(1\,;0\,;3)$ (pour $t=0$) et $\vec{u}(2\,;-1\,;1)$.

\item \textbf{Réponse 1 : $\vec{n}(3\,;-2\,;1)$}.\\
Une équation cartésienne de plan $ax+by+cz+d=0$ admet pour vecteur normal $\vec{n}(a\,;b\,;c)$. Pour $P : 3x - 2y + z - 5 = 0$, $\vec{n}(3\,;-2\,;1)$.

\item \textbf{Réponse 3 : $2$}.\\
Distance du point $O(0,0,0)$ au plan $P : x+2y+2z-6=0$ :
\[
d(O,P) = \frac{|1\cdot0 + 2\cdot0 + 2\cdot0 - 6|}{\sqrt{1^2+2^2+2^2}} = \frac{6}{\sqrt{9}} = \frac{6}{3} = 2.
\]
\end{enumerate}
\end{corrige}

\begin{corrige}[Exercice 2 — Vrai ou faux, en justifiant]
\begin{enumerate}
\item \textbf{Faux}. Trois points distincts définissent un unique plan \textbf{si et seulement s'ils ne sont pas alignés}. S'ils sont alignés, une infinité de plans passent par ces trois points.

\item \textbf{Faux}. Deux droites de l'espace sans point commun sont soit \textbf{parallèles} (coplanaires), soit \textbf{non coplanaires} (droites « gauches »).

\item \textbf{Vrai}. Deux plans ont des vecteurs normaux colinéaires $\iff$ leurs vecteurs normaux sont proportionnels $\iff$ les plans sont parallèles (au sens large : parallèles stricts ou confondus).

\item \textbf{Faux}.\\
$(d) : \vec{u} = \begin{pmatrix}2 \\ -1 \\ 0\end{pmatrix}$ \quad $(d') : \vec{v} = \begin{pmatrix}-4 \\ 2 \\ 0\end{pmatrix}$.\\
On a $\vec{v} = -2\vec{u}$, donc les vecteurs directeurs sont colinéaires : les droites sont \textbf{parallèles} (et distinctes car $z=3$ pour $(d)$ et $z=-1$ pour $(d')$). Elles ne sont pas orthogonales.
\end{enumerate}
\end{corrige}

\begin{corrige}[Exercice 3 — Lecture dans un repère]
\begin{enumerate}
\item Le repère $(A\,;\overrightarrow{AB},\overrightarrow{AD},\overrightarrow{AE})$ est un repère de l'espace car les trois vecteurs $\overrightarrow{AB}$, $\overrightarrow{AD}$, $\overrightarrow{AE}$ sont \textbf{non coplanaires} (ils représentent trois arêtes concourantes d'un cube, deux à deux orthogonaux). De plus, comme le cube a pour côté $1$, ces vecteurs sont de norme $1$ et orthogonaux deux à deux : c'est un \textbf{repère orthonormal}.

\item Coordonnées des sommets (unité : côté du cube = 1) :
\[
\begin{array}{lll}
A(0\,;0\,;0), & B(1\,;0\,;0), & C(1\,;1\,;0), \\
D(0\,;1\,;0), & E(0\,;0\,;1), & F(1\,;0\,;1), \\
G(1\,;1\,;1), & H(0\,;1\,;1).
\end{array}
\]

\item Vecteurs :
\[
\overrightarrow{AG} = \begin{pmatrix}1 \\ 1 \\ 1\end{pmatrix},\quad
\overrightarrow{BH} = \begin{pmatrix}-1 \\ 1 \\ 1\end{pmatrix},\quad
\overrightarrow{EC} = \begin{pmatrix}1 \\ 1 \\ -1\end{pmatrix}.
\]
\end{enumerate}
\end{corrige}

\begin{corrige}[Exercice 4 — Calculs directs]
\begin{enumerate}
\item $\overrightarrow{AB} = \begin{pmatrix}0-2 \\ 4-(-1) \\ 1-3\end{pmatrix} = \begin{pmatrix}-2 \\ 5 \\ -2\end{pmatrix}$.\\
$2\overrightarrow{AB} - 3\vec{u} = 2\begin{pmatrix}-2 \\ 5 \\ -2\end{pmatrix} - 3\begin{pmatrix}-1 \\ 2 \\ 5\end{pmatrix} = \begin{pmatrix}-4 \\ 10 \\ -4\end{pmatrix} - \begin{pmatrix}-3 \\ 6 \\ 15\end{pmatrix} = \begin{pmatrix}-1 \\ 4 \\ -19\end{pmatrix}$.

\item $\overrightarrow{AB}(-2\,;5\,;-2)$ et $\vec{u}(-1\,;2\,;5)$ ne sont \textbf{pas colinéaires} car les rapports des coordonnées ne sont pas égaux : $\frac{-2}{-1}=2 \neq \frac{5}{2}=2,5$.

\item Droite $(AB)$ : passe par $A(2\,;-1\,;3)$, dirigée par $\overrightarrow{AB}(-2\,;5\,;-2)$.
\[
(AB) : \begin{cases} x = 2 - 2t \\ y = -1 + 5t \\ z = 3 - 2t \end{cases},\quad t\in\mathbb{R}.
\]

\item Plan passant par $A(2\,;-1\,;3)$, normal $\vec{n}(1\,;-2\,;2)$.
Équation : $1(x-2) - 2(y+1) + 2(z-3) = 0 \iff x - 2 - 2y - 2 + 2z - 6 = 0$.
\[
\boxed{x - 2y + 2z - 10 = 0}
\]
\end{enumerate}
\end{corrige}

\begin{corrige}[Exercice 5 — Positions relatives de deux droites]
\begin{enumerate}
\item $(d) : \vec{u}=\begin{pmatrix}1 \\ -1 \\ 2\end{pmatrix}$ \quad $(d') : \vec{v}=\begin{pmatrix}-2 \\ 1 \\ -4\end{pmatrix}$.
$\vec{u}$ et $\vec{v}$ ne sont pas colinéaires ($\frac{1}{-2} \neq \frac{-1}{1}$). Droites non parallèles.
Recherche intersection :
\[
\begin{cases}
1+t = 4-2s \\
2-t = 1+s \\
3+2t = 7-4s
\end{cases}
\iff
\begin{cases}
t+2s = 3 \quad (1)\\
-t-s = -1 \quad (2)\\
2t+4s = 4 \quad (3)
\end{cases}
\]
$(1)+(2) \Rightarrow s = 2 \Rightarrow t = -1$.
Vérification (3) : $2(-1)+4(2)=6 \neq 4$. \textbf{Système incompatible}.
\textbf{Conclusion :} Les droites sont \textbf{non coplanaires}.

\item $(d) : \vec{u}=\begin{pmatrix}1 \\ 1 \\ -1\end{pmatrix}$ \quad $(d') : \vec{v}=\begin{pmatrix}-1 \\ 1 \\ 1\end{pmatrix}$.
Non colinéaires.
Système :
\[
\begin{cases}
2+t = 5-s \\
-1+t = 2+s \\
1-t = -2+s
\end{cases}
\iff
\begin{cases}
t+s = 3 \quad (1)\\
t-s = 3 \quad (2)\\
t+s = 3 \quad (3)
\end{cases}
\]
L'équation (3) est identique à (1). $(1)+(2) \Rightarrow 2t=6 \Rightarrow t=3 \Rightarrow s=0$, et (3) est vérifiée.
Le système admet la solution unique $t=3, s=0$.
Point d'intersection : $t=3 \Rightarrow M(5\,;2\,;-2)$.
\textbf{Conclusion :} Les droites sont \textbf{sécantes} en $M(5\,;2\,;-2)$.

\item $(d) : \vec{u}=\begin{pmatrix}1 \\ 2 \\ -1\end{pmatrix}$ \quad $(d') : \vec{v}=\begin{pmatrix}1 \\ 1 \\ 1\end{pmatrix}$.
Non colinéaires.
Système :
\[
\begin{cases}
t = 2+s \\
1+2t = s \\
-t = 1+s
\end{cases}
\]
D'après (1) et (3) : $t = 2+s$ et $t = -1-s \Rightarrow 2+s = -1-s \Rightarrow 2s = -3 \Rightarrow s = -1,5 \Rightarrow t = 0,5$.
Vérification (2) : $1+2(0,5) = 2 \neq -1,5$. \textbf{Incompatible}.
\textbf{Conclusion :} Les droites sont \textbf{non coplanaires}.
\end{enumerate}
\end{corrige}

\begin{corrige}[Exercice 6 — Intersection de deux plans]
\begin{enumerate}
\item $\vec{n}_P = \begin{pmatrix}2 \\ -1 \\ 1\end{pmatrix}$ \quad $\vec{n}_{P'} = \begin{pmatrix}1 \\ 1 \\ -2\end{pmatrix}$.

\item $\vec{n}_P$ et $\vec{n}_{P'}$ ne sont pas colinéaires (les rapports $2/1$, $-1/1$, $1/-2$ ne sont pas égaux). Donc les plans $P$ et $P'$ sont \textbf{sécants}.

\item Système :
\[
\begin{cases}
2x - y + z = 3 \quad (1)\\
x + y - 2z = -1 \quad (2)
\end{cases}
\]
Posons $z = t$ (paramètre).
$(1)+(2) \Rightarrow 3x - t = 2 \Rightarrow x = \frac{t+2}{3}$.
D'après (2) : $y = -1 - x + 2t = -1 - \frac{t+2}{3} + 2t = \frac{-3 - t - 2 + 6t}{3} = \frac{5t - 5}{3}$.
Représentation paramétrique de $(\Delta) = P \cap P'$ :
\[
\boxed{(\Delta) : \begin{cases} x = \frac{2}{3} + \frac{1}{3}t \\ y = -\frac{5}{3} + \frac{5}{3}t \\ z = t \end{cases},\quad t\in\mathbb{R}}
\]
(Vecteur directeur $\vec{u}(1\,;5\,;3)$, point $\bigl(\frac{2}{3}\,;-\frac{5}{3}\,;0\bigr)$).

\item Vérification pour $I(1\,;0\,;1)$ :
$P : 2(1) - 0 + 1 - 3 = 0$ \quad $\checkmark$
$P' : 1 + 0 - 2(1) + 1 = 0$ \quad $\checkmark$
$I$ appartient bien aux deux plans, donc à $(\Delta)$.
(On retrouve $I$ pour $t=1$ dans la paramétrisation ci-dessus).
\end{enumerate}
\end{corrige}

\begin{corrige}[Exercice 7 — Distance d'un point à un plan et volume]
Données : $M(2\,;-1\,;3)$, $P : x + 2y - 2z + 5 = 0$, repère orthonormal (unité : cm).
\begin{enumerate}
\item Distance $d(M,P) = \frac{|1\cdot2 + 2(-1) - 2\cdot3 + 5|}{\sqrt{1^2+2^2+(-2)^2}} = \frac{|2 - 2 - 6 + 5|}{3} = \frac{1}{3}$.
\[
\boxed{d(M,P) = \frac{1}{3}\text{ cm}}
\]

\item Droite $(MH)$ perpendiculaire à $P$ : dirigée par $\vec{n}(1\,;2\,;-2)$, passant par $M$.
\[
(MH) : \begin{cases} x = 2 + t \\ y = -1 + 2t \\ z = 3 - 2t \end{cases}
\]
Intersection avec $P$ :
$(2+t) + 2(-1+2t) - 2(3-2t) + 5 = 0 \iff 2+t -2+4t -6+4t +5 = 0 \iff 9t -1 = 0 \iff t = \frac{1}{9}$.
Coordonnées de $H$ :
\[
H\left(2+\frac{1}{9}\,;\,-1+\frac{2}{9}\,;\,3-\frac{2}{9}\right) = \boxed{H\left(\frac{19}{9}\,;\,-\frac{7}{9}\,;\,\frac{25}{9}\right)}
\]

\item Volume du tétraèdre $M$-$T$ (base $T$ dans $P$, sommet $M$) :
$V = \frac{1}{3} \times \text{Aire}(T) \times d(M,P) = \frac{1}{3} \times 6 \times \frac{1}{3} = \frac{2}{3}$.
\[
\boxed{V = \frac{2}{3}\text{ cm}^3}
\]
\end{enumerate}
\end{corrige}

\begin{corrige}[Exercice 8 — Projeté orthogonal sur une droite]
Données : $A(1\,;1\,;0)$, $\vec{u}(2\,;-1\,;1)$, $M(3\,;0\,;2)$.
\begin{enumerate}
\item $(d) : \begin{cases} x = 1 + 2t \\ y = 1 - t \\ z = t \end{cases},\quad t\in\mathbb{R}$.

\item $H \in (d) \Rightarrow H(1+2t\,;\,1-t\,;\,t)$.
$\overrightarrow{MH} = \begin{pmatrix}1+2t-3 \\ 1-t-0 \\ t-2\end{pmatrix} = \begin{pmatrix}2t-2 \\ 1-t \\ t-2\end{pmatrix}$.

\item Condition d'orthogonalité : $\overrightarrow{MH} \cdot \vec{u} = 0$.
$(2t-2)(2) + (1-t)(-1) + (t-2)(1) = 0$
$4t - 4 -1 + t + t - 2 = 0 \iff 6t - 7 = 0 \iff t = \frac{7}{6}$.
Coordonnées de $H$ (projeté orthogonal) :
\[
H\left(1+2\cdot\frac{7}{6}\,;\,1-\frac{7}{6}\,;\,\frac{7}{6}\right) = \boxed{H\left(\frac{10}{3}\,;\,-\frac{1}{6}\,;\,\frac{7}{6}\right)}
\]

\item Distance $d(M,(d)) = MH = \|\overrightarrow{MH}\|_{t=7/6}$.
$\overrightarrow{MH} = \begin{pmatrix}2\cdot\frac{7}{6}-2 \\ 1-\frac{7}{6} \\ \frac{7}{6}-2\end{pmatrix} = \begin{pmatrix}\frac{1}{3} \\ -\frac{1}{6} \\ -\frac{5}{6}\end{pmatrix}$.
$MH = \sqrt{\left(\frac{1}{3}\right)^2 + \left(-\frac{1}{6}\right)^2 + \left(-\frac{5}{6}\right)^2} = \sqrt{\frac{4}{36} + \frac{1}{36} + \frac{25}{36}} = \sqrt{\frac{30}{36}} = \frac{\sqrt{30}}{6}$.
\[
\boxed{d(M,(d)) = \frac{\sqrt{30}}{6}}
\]
\end{enumerate}
\end{corrige}

\begin{corrige}[Exercice 9 — Étude d'un plan et d'une droite — type Probatoire]
\textbf{Barème indicatif (sur 20) :} 1.a (2 pts), 1.b (2 pts), 2.a (4 pts), 2.b (2 pts), 3.a (1 pt), 3.b (2 pts), 3.c (3 pts), 4 (4 pts).\\
\textbf{Points de vigilance :} Justifier la non-colinéarité pour définir un plan (1.b) ; calculer correctement le produit vectoriel ou résoudre le système pour $\vec{n}$ (2.a) ; vérifier $\vec{u}\cdot\vec{n}\neq 0$ pour non-parallélisme (3.b) ; distance point/plan = 0 car $D\in(ABC)$ (4).

\begin{enumerate}
\item \begin{enumerate}
\item $\overrightarrow{AB} = \begin{pmatrix}2 \\ -2 \\ 2\end{pmatrix},\quad \overrightarrow{AC} = \begin{pmatrix}-1 \\ -1 \\ 3\end{pmatrix}$.
\item $\overrightarrow{AB}$ et $\overrightarrow{AC}$ ne sont pas colinéaires ($\frac{2}{-1}\neq\frac{-2}{-1}$). Les points $A,B,C$ ne sont pas alignés, ils définissent un unique plan $(ABC)$.
\end{enumerate}
\item \begin{enumerate}
\item Recherche $\vec{n}(a,b,c)$ orthogonal à $\overrightarrow{AB}$ et $\overrightarrow{AC}$ :
\[
\begin{cases}
2a - 2b + 2c = 0 \\
-a - b + 3c = 0
\end{cases}
\iff
\begin{cases}
a - b + c = 0 \\
a + b - 3c = 0
\end{cases}
\]
En ajoutant : $2a - 2c = 0 \Rightarrow a = c$. Alors $c - b + c = 0 \Rightarrow b = 2c$.
On choisit $c=1 \Rightarrow \vec{n}(1\,;2\,;1)$.
(Produit vectoriel : $\overrightarrow{AB}\wedge\overrightarrow{AC} = \begin{vmatrix}\vec{i}&\vec{j}&\vec{k}\\2&-2&2\\-1&-1&3\end{vmatrix} = \vec{i}(-6+2) - \vec{j}(6+2) + \vec{k}(-2-2) = (-4,-8,-4) \sim (1,2,1)$).
\item Équation cartésienne de $(ABC)$ : $1(x-1) + 2(y-2) + 1(z+1) = 0 \iff x + 2y + z - 4 = 0$.
\[
\boxed{(ABC) : x + 2y + z - 4 = 0}
\]
\end{enumerate}
\item \begin{enumerate}
\item $(\Delta) : D(2\,;1\,;0)$, $\vec{u}(1\,;1\,;1)$.
\[
(\Delta) : \begin{cases} x = 2 + t \\ y = 1 + t \\ z = t \end{cases},\quad t\in\mathbb{R}.
\]
\item $\vec{u}\cdot\vec{n} = 1 + 2 + 1 = 4 \neq 0 \Rightarrow (\Delta)$ n'est pas parallèle à $(ABC)$.
\item Intersection $K$ : substitution dans $x+2y+z-4=0$ :
$(2+t) + 2(1+t) + t - 4 = 0 \iff 2+t+2+2t+t-4=0 \iff 4t=0 \iff t=0$.
$K = (2\,;1\,;0) = D$.
\textbf{Conclusion :} Le point $D$ \textbf{appartient} au plan $(ABC)$. L'intersection est le point $D$ lui-même.
\end{enumerate}
\item Distance de $D$ au plan $(ABC)$ : comme $D\in(ABC)$, la distance est \textbf{nulle}.
\[
\boxed{d(D,(ABC)) = 0}
\]
\end{enumerate}
\end{corrige}

\begin{corrige}[Exercice 10 — Le cube du menuisier — problème de synthèse]
Repère orthonormal $(A\,;\overrightarrow{AB},\overrightarrow{AD},\overrightarrow{AE})$, unité : mètre.
$A(0,0,0)$, $B(1,0,0)$, $D(0,1,0)$, $E(0,0,1)$, $G(1,1,1)$.
\begin{enumerate}
\item $B(1\,;0\,;0)$, $D(0\,;1\,;0)$, $E(0\,;0\,;1)$, $G(1\,;1\,;1)$.

\item \begin{enumerate}
\item $\overrightarrow{BE} = \begin{pmatrix}-1 \\ 0 \\ 1\end{pmatrix}$, $\overrightarrow{BG} = \begin{pmatrix}0 \\ 1 \\ 1\end{pmatrix}$.
Non colinéaires (premières coordonnées $-1$ et $0$).
\item Plan $(BEG)$ : passe par $B(1,0,0)$, vecteurs directeurs $\overrightarrow{BE},\overrightarrow{BG}$.
Vecteur normal $\vec{n} = \overrightarrow{BE}\wedge\overrightarrow{BG} = \begin{vmatrix}\vec{i}&\vec{j}&\vec{k}\\-1&0&1\\0&1&1\end{vmatrix} = \vec{i}(0-1) - \vec{j}(-1-0) + \vec{k}(-1-0) = \begin{pmatrix}-1 \\ 1 \\ -1\end{pmatrix} \sim \begin{pmatrix}1 \\ -1 \\ 1\end{pmatrix}$.
Équation : $1(x-1) -1(y-0) + 1(z-0) = 0 \iff x - y + z - 1 = 0$.
\[
\boxed{(BEG) : x - y + z - 1 = 0}
\]
\item $A(0,0,0)$ : $0 - 0 + 0 - 1 = -1 \neq 0 \Rightarrow A\notin (BEG)$.
Distance $d(A,(BEG)) = \frac{|0-0+0-1|}{\sqrt{1^2+(-1)^2+1^2}} = \frac{1}{\sqrt{3}} = \frac{\sqrt{3}}{3}$.
\[
\boxed{d(A,(BEG)) = \frac{\sqrt{3}}{3}\text{ m}}
\]
\end{enumerate}
\item \begin{enumerate}
\item $(AG) : A(0,0,0)$, $\overrightarrow{AG}(1,1,1)$.
\[
(AG) : \begin{cases} x = t \\ y = t \\ z = t \end{cases},\quad t\in\mathbb{R}.
\]
\item \textbf{Attention :} Vecteur directeur de $(AG)$ : $\vec{u}=(1,1,1)$. Vecteur normal de $(BEG)$ : $\vec{n}=(1,-1,1)$.
$\vec{u}\cdot\vec{n} = 1 - 1 + 1 = 1 \neq 0$.
\textbf{Conclusion :} La droite $(AG)$ n'est \textbf{PAS perpendiculaire} au plan $(BEG)$. (C'est la droite $(AG)$ qui est perpendiculaire au plan $(BDE)$).
\item Intersection $\Omega = (AG) \cap (BEG)$ :
Substitution $x=y=z=t$ dans $x-y+z-1=0 \Rightarrow t - t + t - 1 = 0 \Rightarrow t=1$.
$\Omega(1\,;1\,;1) = G$.
\textbf{Conclusion :} $\Omega$ est le sommet $G$ du triangle $BEG$ (donc un \textbf{sommet} du triangle).
\end{enumerate}
\item Droites $(AG)$ et $(BE)$ :
$(AG) : (t,t,t)$ \quad $(BE) : B + s\overrightarrow{BE} = (1-s,\,0,\,s)$.
Système : $t = 1-s$, $t = 0$, $t = s$.
$t=0 \Rightarrow s=0 \Rightarrow 1-s=1\neq 0$. \textbf{Pas d'intersection}.
Vecteurs directeurs $(1,1,1)$ et $(-1,0,1)$ non colinéaires.
\textbf{Conclusion :} Les droites $(AG)$ et $(BE)$ sont \textbf{non coplanaires}.
\end{enumerate}
\end{corrige}

\begin{corrige}[Exercice 11 — Antenne relais et zone de couverture — type Probatoire]
Repère orthonormal $(O\,;\vec{i},\vec{j},\vec{k})$, unité : $100$ m.
$P : 2x+y-2z+4=0$, $\vec{n}_P(2,1,-2)$.
$F(1,2,3)$, $S(3,3,0)$, $T(5,-1,2)$.
\begin{enumerate}
\item \begin{enumerate}
\item $\overrightarrow{FS} = \begin{pmatrix}2 \\ 1 \\ -3\end{pmatrix}$.
$(FS) : \begin{cases} x = 1 + 2t \\ y = 2 + t \\ z = 3 - 3t \end{cases},\quad t\in\mathbb{R}$.
\item Pylône $\perp$ sol $\iff \overrightarrow{FS} \parallel \vec{n}_P$.
$\overrightarrow{FS}(2,1,-3)$ et $\vec{n}_P(2,1,-2)$ ne sont pas colinéaires (3\textsuperscript{e} coordonnée).
\textbf{Conclusion :} Le pylône n'est \textbf{pas perpendiculaire} au sol.
\end{enumerate}
\item \begin{enumerate}
\item $d(T,P) = \frac{|2\cdot5 + 1(-1) -2\cdot2 + 4|}{\sqrt{4+1+4}} = \frac{|10-1-4+4|}{3} = \frac{9}{3} = 3$ (unités).
Distance réelle : $3 \times 100 = \boxed{300\text{ m}}$.
\item Projeté $H$ de $T$ sur $P$ : droite $(TH)$ dirigée par $\vec{n}_P(2,1,-2)$.
$(TH) : \begin{cases} x = 5 + 2u \\ y = -1 + u \\ z = 2 - 2u \end{cases}$.
Intersection avec $P$ : $2(5+2u) + (-1+u) -2(2-2u) + 4 = 0$
$\iff 10+4u -1+u -4+4u +4 = 0 \iff 9u + 9 = 0 \iff u = -1$.
$H(5-2\,;\,-1-1\,;\,2+2) = \boxed{H(3\,;-2\,;4)}$.
\end{enumerate}
\item \begin{enumerate}
\item $\overrightarrow{TS} = \begin{pmatrix}-2 \\ 4 \\ -2\end{pmatrix} \sim \begin{pmatrix}-1 \\ 2 \\ -1\end{pmatrix}$.
$(TS) : \begin{cases} x = 5 - t \\ y = -1 + 2t \\ z = 2 - t \end{cases},\quad t\in\mathbb{R}$.
\item $\overrightarrow{TS}\cdot\vec{n}_P = (-2)(2) + 4(1) + (-2)(-2) = -4+4+4 = 4 \neq 0 \Rightarrow (TS)$ n'est pas parallèle à $P$, donc \textbf{sécante}.
Intersection $I$ : substitution dans $P$ :
$2(5-t) + (-1+2t) -2(2-t) + 4 = 0$
$\iff 10-2t -1+2t -4+2t +4 = 0 \iff 9 + 2t = 0 \iff t = -\frac{9}{2}$.
$I\left(5+\frac{9}{2}\,;\,-1-9\,;\,2+\frac{9}{2}\right) = \boxed{I\left(\frac{19}{2}\,;-10\,;\frac{13}{2}\right)}$.
\end{enumerate}
\item \begin{enumerate}
\item $Q : x - y + z - 1 = 0$, $\vec{n}_Q(1,-1,1)$.
$\vec{n}_P(2,1,-2)$ et $\vec{n}_Q(1,-1,1)$ non colinéaires $\Rightarrow P$ et $Q$ \textbf{sécants}.
\item Intersection $P \cap Q$ :
\[
\begin{cases}
2x + y - 2z = -4 \quad (1)\\
x - y + z = 1 \quad (2)
\end{cases}
\]
Posons $z = t$.
$(1)+(2) \Rightarrow 3x - z = -3 \Rightarrow 3x = t - 3 \Rightarrow x = \frac{t-3}{3}$.
D'après (2) : $y = x + z - 1 = \frac{t-3}{3} + t - 1 = \frac{t-3+3t-3}{3} = \frac{4t-6}{3}$.
Représentation paramétrique :
\[
\boxed{P\cap Q : \begin{cases} x = -1 + \frac{1}{3}t \\ y = -2 + \frac{4}{3}t \\ z = t \end{cases},\quad t\in\mathbb{R}}
\]
(Vecteur directeur $(1,4,3)$, point $(-1,-2,0)$).
\end{enumerate}
\end{enumerate}
\end{corrige}

\newpage

\coursec{Fiche bilan}

\begin{popfigure}
\begin{tikzpicture}[mindmap, grow cyclic, every node/.style={concept, font=\small},
  root concept/.append style={concept color=popblue, font=\bfseries},
  level 1/.append style={level distance=3.6cm, sibling angle=60},
  level 2/.append style={level distance=2.4cm, sibling angle=45, font=\scriptsize}]
\node[root concept]{Géométrie analytique de l'espace}
  child[concept color=poporange]{ node {Repères et coordonnées}
    child { node {Base $(\vec{\imath},\vec{\jmath},\vec{k})$} }
    child { node {$\overrightarrow{AB}\,(x_B-x_A\,;\,\dots)$} }
  }
  child[concept color=popgreen]{ node {Droites de l'espace}
    child { node {Représentation paramétrique} }
    child { node {Vecteur directeur $\vec{u}$} }
  }
  child[concept color=poppurple]{ node {Positions de deux droites}
    child { node {Parallèles ou sécantes} }
    child { node {Non coplanaires} }
  }
  child[concept color=poppink]{ node {Plans}
    child { node {Représentation paramétrique} }
    child { node {$ax+by+cz+d=0$} }
    child { node {Vecteur normal $\vec{n}$} }
  }
  child[concept color=popgold]{ node {Positions de deux plans}
    child { node {Parallèles} }
    child { node {Sécants} }
  }
  child[concept color=popblue!70]{ node {Projetés et distances}
    child { node {$d(M,\mathcal{P})=\dfrac{|ax_M+by_M+cz_M+d|}{\sqrt{a^2+b^2+c^2}}$} }
    child { node {$d(M,\mathcal{D})$} }
  };
\end{tikzpicture}
\legende{Carte mentale de la leçon : les six grandes idées de la géométrie analytique dans l'espace.}
\end{popfigure}

\begin{bilan}
\textbf{1. Repères et coordonnées}
\begin{itemize}
  \item Un \cle{repère} de l'espace est un quadruplet $(O\,;\vec{\imath},\vec{\jmath},\vec{k})$ où les trois vecteurs sont \cle{non coplanaires}. Il est \cle{orthonormé} si les vecteurs sont unitaires et orthogonaux deux à deux.
  \item $\overrightarrow{AB}\begin{pmatrix} x_B-x_A \\ y_B-y_A \\ z_B-z_A \end{pmatrix}$ ; milieu de $[AB]$ : $I\left(\dfrac{x_A+x_B}{2}\,;\,\dfrac{y_A+y_B}{2}\,;\,\dfrac{z_A+z_B}{2}\right)$.
  \item Distance : $AB=\sqrt{(x_B-x_A)^2+(y_B-y_A)^2+(z_B-z_A)^2}$ (repère orthonormé).
  \item Colinéarité : $\vec{u}$ et $\vec{v}$ colinéaires $\iff$ il existe $k\in\mathbb{R}$ tel que $\vec{v}=k\vec{u}$ $\iff$ leurs coordonnées sont proportionnelles.
\end{itemize}

\textbf{2. Droite de l'espace}

Droite passant par $A(x_A;y_A;z_A)$, de vecteur directeur $\vec{u}\begin{pmatrix} \alpha \\ \beta \\ \gamma \end{pmatrix}$ ($\vec{u}\neq\vec{0}$) :
\[ M(x;y;z)\in(AB) \iff \overrightarrow{AM}=t\,\vec{u} \iff \begin{cases} x=x_A+\alpha t \\ y=y_A+\beta t \\ z=z_A+\gamma t \end{cases} \quad (t\in\mathbb{R}) \]

\textbf{3. Positions relatives de deux droites}
\begin{itemize}
  \item Vecteurs directeurs \textbf{colinéaires} $\Rightarrow$ droites \textbf{parallèles} (confondues si elles ont un point commun, strictement parallèles sinon).
  \item Vecteurs directeurs \textbf{non colinéaires} : résoudre le système des représentations paramétriques (deux paramètres distincts $t$ et $t'$).
  \begin{itemize}
    \item Solution unique $\Rightarrow$ droites \textbf{sécantes} (coplanaires).
    \item Pas de solution $\Rightarrow$ droites \textbf{non coplanaires}.
  \end{itemize}
\end{itemize}

\textbf{4. Plan : représentations}

Plan passant par $A$, dirigé par $\vec{u}$ et $\vec{v}$ \textbf{non colinéaires} :
\[ M\in\mathcal{P} \iff \overrightarrow{AM}=t\vec{u}+t'\vec{v} \iff \begin{cases} x=x_A+\alpha t+\alpha' t' \\ y=y_A+\beta t+\beta' t' \\ z=z_A+\gamma t+\gamma' t' \end{cases} \quad (t,t'\in\mathbb{R}) \]
Un plan est aussi défini par \textbf{trois points non alignés}, ou par \textbf{un point et un vecteur normal}.

\textbf{5. Équation cartésienne d'un plan, vecteur normal}
\begin{center}
\begin{tabularx}{\linewidth}{|l|X|}
\hline
\rowcolor{popblueL} \textbf{Objet} & \textbf{Résultat à connaître par c\oe ur} \\
\hline
Vecteur normal & $\vec{n}\begin{pmatrix} a \\ b \\ c \end{pmatrix}\neq\vec{0}$ est normal à $\mathcal{P}$ $\iff$ $\mathcal{P}$ : $ax+by+cz+d=0$ \\
\hline
Trouver $d$ & Remplacer $(x;y;z)$ par les coordonnées d'un point $A\in\mathcal{P}$ \\
\hline
Appartenance & $M\in\mathcal{P} \iff ax_M+by_M+cz_M+d=0$ \\
\hline
Plans parallèles & Vecteurs normaux colinéaires (confondus ou strictement parallèles) \\
\hline
Plans sécants & Vecteurs normaux non colinéaires ; l'intersection est une \textbf{droite} \\
\hline
Plans perpendiculaires & $\vec{n}\cdot\vec{n}'=0$ (repère orthonormé) \\
\hline
\end{tabularx}
\end{center}

\textbf{6. Projetés orthogonaux et distances} (repère orthonormé)
\begin{itemize}
  \item \cle{Projeté orthogonal} de $M$ sur $\mathcal{P}$ : point $H$ de $\mathcal{P}$ tel que $\overrightarrow{MH}$ soit colinéaire à $\vec{n}$. Méthode : paramétrer la droite $(M,\vec{n})$, puis injecter dans l'équation de $\mathcal{P}$.
  \item Distance d'un point à un plan :
  \[ d(M,\mathcal{P})=MH=\dfrac{|ax_M+by_M+cz_M+d|}{\sqrt{a^2+b^2+c^2}} \]
  \item Projeté orthogonal de $M$ sur une droite $\mathcal{D}(A,\vec{u})$ : point $H\in\mathcal{D}$ tel que $\overrightarrow{MH}\cdot\vec{u}=0$ ; alors $d(M,\mathcal{D})=MH$.
\end{itemize}

\textbf{Méthodes express}
\begin{enumerate}
  \item \emph{Montrer que deux droites sont sécantes} : vecteurs directeurs non colinéaires, puis système à solution unique ; donner le point d'intersection.
  \item \emph{Équation d'un plan par trois points} : vérifier que les points sont non alignés, poser $ax+by+cz+d=0$ et résoudre le système obtenu en remplaçant par les trois points.
  \item \emph{Vecteur normal à partir de deux vecteurs directeurs} : chercher $\vec{n}(a;b;c)$ avec $\vec{n}\cdot\vec{u}=0$ et $\vec{n}\cdot\vec{v}=0$.
\end{enumerate}
\end{bilan}

\begin{aretenir}[Checklist avant le Probatoire]
\begin{itemize}
  \item[$\square$] Je sais vérifier qu'un triplet de vecteurs forme une base de l'espace (non coplanarité).
  \item[$\square$] Je sais calculer les coordonnées de $\overrightarrow{AB}$, d'un milieu et une distance $AB$.
  \item[$\square$] Je sais écrire la représentation paramétrique d'une droite à partir d'un point et d'un vecteur directeur.
  \item[$\square$] Je sais tester la colinéarité de deux vecteurs (coordonnées proportionnelles).
  \item[$\square$] Je sais discuter la position relative de deux droites : parallèles, sécantes ou non coplanaires.
  \item[$\square$] Je sais écrire la représentation paramétrique d'un plan (un point, deux vecteurs non colinéaires).
  \item[$\square$] Je sais déterminer une équation cartésienne $ax+by+cz+d=0$ d'un plan passant par trois points.
  \item[$\square$] Je sais lire un vecteur normal sur l'équation d'un plan et l'utiliser pour étudier le parallélisme de deux plans.
  \item[$\square$] Je sais déterminer l'intersection de deux plans sécants (c'est une droite).
  \item[$\square$] Je sais calculer le projeté orthogonal d'un point sur un plan et la distance $d(M,\mathcal{P})=\dfrac{|ax_M+by_M+cz_M+d|}{\sqrt{a^2+b^2+c^2}}$.
  \item[$\square$] Je sais calculer le projeté orthogonal d'un point sur une droite grâce à $\overrightarrow{MH}\cdot\vec{u}=0$.
  \item[$\square$] Je pense à vérifier les hypothèses : vecteur directeur non nul, vecteurs non colinéaires, points non alignés.
\end{itemize}
\end{aretenir}

\findecours

\end{document}
